% 本文件由 LaTeX 排版 Agent 根据有效 translation.json 自动生成。
% author=Miscellaneous; work_id=0715; title=宗徒宪章

\chapter{宗徒宪章}

% structure: p0001 source=h1 target=omitted reason=page-title-already-rendered-by-parent

\Emph{卷一}\newline{}\Emph{卷二}\newline{}\Emph{卷三}\newline{}\Emph{卷四}\newline{}\Emph{卷五}\newline{}\Emph{卷六}\newline{}\Emph{卷七}\newline{}\Emph{卷八}\par

\SourceRecord{Translated by James Donaldson.；Ante-Nicene Fathers；Vol. 7.；Edited by Alexander Roberts, James Donaldson, and A. Cleveland Coxe.；Buffalo, NY: Christian Literature Publishing Co.,；1886.；<http://www.newadvent.org/fathers/0715.htm>.}

% structure: p0001 source=h1 target=section reason=page-title

\section{宗徒宪章（卷一）}

\Emph{论平信徒。}\par

% structure: p0003 source=h2 target=subsection reason=relative-heading-rank; base=h2

\subsection{第一节 总则}

宗徒与长老，致所有从外邦中相信了主耶稣基督的人：愿恩宠与平安，由全能的天主，透过我们的主耶稣基督，在认识祂的事上，多多地加给你们。\par

天主教会是天主的栽种和祂心爱的葡萄园\BibleRef{依 5:7}，其中包含那些相信了祂无谬神圣宗教的人，他们因信德成为祂永恒国度的嗣子，分享祂的神性影响和圣神的通传，藉耶稣受装备并将祂的敬畏接受到心中，享受宝贵的、无罪的基督之血洒布的益处，并得以自由地称全能的天主为父，同作祂爱子的共同嗣子和分享者。你们这些享有祂许诺的人，请聆听这神圣的教导，因为它是按你们救主的命令传授的，且符合祂荣耀的话语。天主的子女们，你们要留心，在一切事上服从天主，在一切事上中悦我们的主基督。因为若有人追随不义，行违背天主旨意的事，这样的人必被天主视为不服从的外邦人。\par

% structure: p0006 source=h3 target=subsubsection reason=relative-heading-rank; base=h2

\subsubsection{论贪欲。}

一、所以，要戒绝一切非法的欲望和不义。因为法律上记载着：\BibleQuote{不可贪恋你近人的妻子、田地、仆婢、牛驴，并你近人所有的一切。}\BibleRef{出 20:17} 因为对这一切的贪恋都出自那恶者。凡贪恋近人的妻子或仆婢的，心里已犯了奸淫和偷窃；若不悔改，将被我们的主耶稣基督定罪：愿荣耀藉着祂归于天主，直到永远。阿们。因为祂在福音中这样说，总结、确认并成全了法律的十诫：\BibleQuote{法律上记载：不可奸淫。但我告诉你们，即我在法律中藉梅瑟所说的，但现在我亲自对你们说：凡注视近人的妻子而贪恋她的，心里已与她犯了奸淫。} 凡心里贪恋近人妻子的，就被定为犯奸淫。但贪恋牛或驴的人，岂不图谋偷窃它们，据为己用，牵走它们吗？再者，贪恋田地并持守此心态的人，岂不邪恶地设法挪移地界，强迫占有者无偿地让出些什么？正如先知某处所说：\BibleQuote{祸哉，那些以房接房，以田连田，以致剥夺近人所有之物的人。} 因此祂说：\BibleQuote{难道你一人要住在地上吗？这些事已传人万军之主的耳中。} 别处又说：\BibleQuote{凡挪移邻舍地界的，必受诅咒；众百姓要说：阿们。}\EditorialNote{申命纪 27:17} 所以梅瑟说：\BibleQuote{不可挪移你邻舍的地界，即你祖先所立的。}\BibleRef{申 19:14} 因此，对这样的人，必有从天主而来的恐怖、死亡、审判和定罪。至于那些服从天主的人，有一条天主的法律，单纯、真实、活泼，即：\BibleQuote{你不愿别人对你做的，你也不要做给别人。}\BibleRef{多 4:16} 你不愿有人以邪念注视你的妻子以败坏她，所以也不要以恶意注视近人的妻子。你不愿你的衣服被人拿走，所以也不要拿走别人的。你不愿被人打、辱骂、侮辱，所以也不要那样对待别人。\par

% structure: p0008 source=h3 target=subsubsection reason=relative-heading-rank; base=h2

\subsubsection{我们不应报复伤害，也不应为自己申冤。}

二、若有人咒骂你，你要祝福他。因为户籍纪上记载：\BibleQuote{祝福你的，必蒙祝福；咒骂你的，必受咒骂。}\BibleRef{户 24:9} 同样，福音上记载：\BibleQuote{要祝福那咒骂你们的人。}\BibleRef{路 6:28} 受伤害时，不要报复，要耐心忍受；因为圣经这样说：\BibleQuote{不要说：我要报复敌人对我的伤害；而要忍受，让主为你伸冤，使那伤害你的人遭报。}\BibleRef{箴 20:22} 因为祂又在福音中说：\BibleQuote{要爱你们的仇敌，善待那恨你们的人，为那凌辱你们、迫害你们的祈祷；这样，你们才可以成为你们天父的子女，因为祂使太阳照耀恶人也照耀善人，降雨给义人也给不义的人。}\BibleRef{玛 5:44} 所以，亲爱的，我们要遵守这些诫命，好让我们因遵守它们而成为光明之子。所以，你们天主的仆人和子女，要彼此担待。\par

% structure: p0010 source=h2 target=subsection reason=relative-heading-rank; base=h2

\subsection{第二节 给男性的诫命}

% structure: p0011 source=h3 target=subsubsection reason=relative-heading-rank; base=h2

\subsubsection{论自身的装饰及由此产生的罪。}

要让丈夫对自己的妻子不傲慢、不狂妄，反要同情、慷慨，乐意唯独取悦自己的妻子，尊敬她、殷勤待她，努力使她喜悦；（三、）不可如此装饰自己，以致引诱别的女子到你这里来。因为若你被她胜过，与她犯了罪，天主的永死将降到你身上；你必将受可感知而痛苦的折磨。或者你没有做出这样的恶行，而是摆脱她、拒绝她，在此情况下你并非完全无罪，即使你没有犯那罪行本身，而只在于因你的装扮引诱了那女人贪恋你。因为是你导致那女子如此受影响，她因贪恋你而犯了与你通奸的罪；然而你并非同等有罪，因为你没有派人去找那个被你迷惑的女子，你也没有贪恋她。因此，既然你没有将自己交付给她，你将在主你的天主那里蒙受怜悯，祂曾说：\Quote{你不可奸淫}，和\Quote{你不可贪恋。}因为若有这样的女子，因看见你或不合时宜地遇见你，心中受打击，派人到你那里，但你作为一个虔诚的人拒绝了她，她若因你的美貌、年轻和装扮而心中受伤并爱上了你，你将因她的过犯而被定罪，因为你成了她绊倒的缘由\BibleRef{玛 18:7}，并要遭受祸害。因此，你要向主天主祈祷，免得因此灾祸临到你；因为你不可讨人喜欢而犯罪，而要讨天主喜欢，以获得圣洁的生活，并分享永远的安息。天主和本性赐予你的美貌，不要再加修饰；反而在人前谦逊地减少它。因此，不可让头上的头发留得太长，而要剪短；免得因精心梳理头发、留长并抹油，招致那些被迷惑或迷惑人的女子。你也不可穿过于精美的衣服去诱惑人；也不可怀着邪恶的狡诈，讲究过于精美的袜子和鞋子，只当穿着合乎庄重和实用的尺度。你也不可在手指上戴金戒指；因为所有这些装饰都是淫荡的记号，你若以不当的方式追求这些，就不像一个好人应有的行为：因为你不——一个信徒和属天主的人——不可让头上的头发留长并梳拢在一起，也不可让它们披散开来，也不可使之蓬松，也不可通过精心梳理和编织使之卷曲发亮；因为这与法律相悖，法律在其附加的诫命中这样说：\Quote{你不可为自己做卷发和圆形剃发。}男人也不可毁掉胡须的毛发，反自然地改变男人的形象。因为法律规定：\Quote{你不可损坏你的胡须。}因为造物主天主使这为女人合宜，但判定它不适用于男人。但你若做这些事去讨人喜欢，违背法律，你在按自己肖像造你的天主面前将是可憎的。因此，你若愿意蒙天主悦纳，就要戒绝所有祂憎恶的事，不做任何不悦乐祂的事。\par

% structure: p0013 source=h3 target=subsubsection reason=relative-heading-rank; base=h2

\subsubsection{我们不应过分好奇那些生活邪恶的人，而应专注于自己的本分工作。}

四、你不可如同流浪者和漂泊者，无正当理由在街上闲逛，窥探那些行为邪恶的人。而要专心于自己的职业和工作，努力做天主所喜悦的事。并且牢记基督的圣言，不断默想它们。因为圣经这样对你说：\Quote{你应日夜默想祂的法律；当你行在田野时，坐在家中时，躺下时，起来时，都要这样，好使你在一切事上有智慧。}\BibleRef{苏 1:8}不但如此，即使你富有，不需要靠职业维持生计，也不可成为一个游荡者，随意外出；而要去到一些信徒和同信仰者那里，与他们谈论天主的活泼圣言。——\par

% structure: p0015 source=h3 target=subsubsection reason=relative-heading-rank; base=h2

\subsubsection{我们应当阅读哪些圣经。}

五、或者你若在家，阅读法律书、列王纪和先知书；咏唱达味的诗篇；并勤读福音，它是其他经书的圆满。\par

% structure: p0017 source=h3 target=subsubsection reason=relative-heading-rank; base=h2

\subsubsection{我们应当远离所有教会之外之人的著作。}

六、要远离一切外邦人的书籍。因为这些外邦人的言论、法律或假先知与你何干？它们会动摇不坚定者的信仰。难道你从天主的法律中发现了什么缺陷，以至于要去求助于那些外邦人的寓言？如果你喜欢读历史，有\BibleBook{列王}{纪}；如果喜欢智慧或诗歌，有\BibleBook{先知}{书}、\BibleBook{约伯}{传}和\BibleBook{箴言}{篇}，你会在其中发现比所有外邦诗人和诡辩家更深邃的智慧，因为这是主、唯一智慧的天主的话。如果你想要歌唱，有\BibleBook{圣咏}{集}；如果你想要万物的起源，有\BibleBook{创世}{纪}；如果你想要法律和规条，有主天主荣耀的法律。因此，要彻底远离所有奇怪和魔鬼的书籍。不仅如此，当你读法律时，不要认为自己必须遵守附加的诫命；虽然不是全部，但有一些。你阅读这些只是为了历史知识，为了了解它们，并光荣天主使你从如此众多和巨大的束缚中解脱出来。你要自己分辨哪些规条来自自然律，哪些是后来添加的，或者哪些是金牛犊事件后在旷野中向以色列人引入和赐下的附加规条；因为法律中包含主天主在百姓陷入偶像崇拜并制造像埃及亚匹斯那样的牛犊之前所说的话——那就是十诫。至于他们犯罪后进一步加在他们身上的束缚，不要引到自己身上；因为我们的救主降临的唯一目的就是要使那些应受这些束缚的人从为他们保留的愤怒中解脱出来，祂要成全法律和先知，并且要废除或改变那些附加在法律之上的次要束缚。因此，祂向我们呼喊说：\BibleQuote{凡劳苦和负重担的，你们都到\Emph{我跟前}来，我要使你们得安息。}\BibleRef{玛 11:28}因此，当你读了与福音和先知相符的法律后，也要读\BibleBook{列王}{纪}，以便从中知道哪些君王是正义的，他们如何蒙天主祝福，永生的应许如何从祂那里与他们同在；而那些背离天主的君王，不久便因他们的背道而受到天主正义的审判而灭亡，失去了祂的生命，不是得到安息，而是继承永罚。因此，阅读这些书籍会使你在信德上大大坚强，并在基督内受到造就，你是祂的身体和肢体。此外，当你在公共场合行走并想要沐浴时，要使用专为男性准备的浴室，以免你以不雅的方式向妇女暴露身体，或看到不适合男性的景象，从而要么你自己被诱惑，要么你诱惑并引那些容易屈服于这种诱惑的妇女。因此，要小心避免这些事，免得你给自己的灵魂带来陷阱。\par

% structure: p0019 source=h3 target=subsubsection reason=relative-heading-rank; base=h2

\subsubsection{论恶妇。}

七、让我们学习圣言在\BibleBook{智慧}{书}中所说的话：\BibleQuote{我儿，你要遵守我的话，将我的诫命珍藏在你心中。对智慧说：你是我的姊妹，并要使明智成为你的亲人：好能保护你远离陌生邪恶的女人，如果她用甜言蜜语与你交谈。因为她从她家的窗口向街上观望，要在愚昧之子中，看到一个无知的少年，在市场上行走，在她家附近的街口，在黄昏的昏暗或夜晚的寂静黑暗中说话。一个貌似淫妇的女人遇见他，她偷走了年轻人的心。她游荡散漫，她的脚不停在家中；她有时在外面，有时在街市上，在每个角落埋伏。于是她抓住他，亲吻他，厚着脸皮对他说：我这里有平安祭，今日我还了我的愿；因此我出来迎接你，切切寻求你的面貌，竟找到了你。我已用床单铺好了我的床，用埃及的绣花巾装饰了它。我以番红花薰了我的床，以肉桂薰了我的房屋。来吧，让我们饱享爱情直到早晨；来吧，让我们彼此亲爱欢娱。}等等。接着他又说：\BibleQuote{她用许多花言巧语诱惑他，用她嘴唇的圈套强迫他。他像一只愚笨的鸟跟随她。}又说：\BibleQuote{不要听从邪恶的女人；因为淫妇的嘴唇好像蜂房滴蜜，在喉咙里暂时甘甜，但以后你必发现她比苦胆更苦，比双刃剑更锋利。}又说：\BibleQuote{你要赶快离开，不要停留；不要眼目注视她；因为她使许多人受伤倒下，被她杀害的无数众多。}\BibleQuote{不然，}他说，\BibleQuote{你终必在最后懊悔，当你的皮肉和身体被消耗时，你必说：我何等恨恶训诲，我的心藐视义人的责备！我没有听从我师傅的声音，也没有侧耳听我教师的训诲。我几乎陷于一切邪恶。}我们不再多引；即使我们有所遗漏，你也应明智地从圣经中拣选最宝贵的内容，以此坚固自己，弃绝一切恶事，好使你在永远的生命中，在天主面前成为圣洁。\par

% structure: p0021 source=h2 target=subsection reason=relative-heading-rank; base=h2

\subsection{第三节 对妇女的诫命}

% structure: p0022 source=h3 target=subsubsection reason=relative-heading-rank; base=h2

\subsubsection{论妻子对丈夫的服从，以及她必须温柔和端庄。}

VIII. 妻子当顺服自己的丈夫，因为\BibleQuote{丈夫是妻子的头}\BibleRef{格前 11:3}。但那行义之人的头是基督，而\Quote{基督的头是天主}，即祂的父。因此，妻子啊，仅次于全能者、我们的天主和父、今世及来世的主、一切气息与权能的创造者，以及祂的爱子、我们的主耶稣基督（愿荣耀藉着祂归于天主），你当敬畏你的丈夫，尊重他，唯独取悦他，在生活的各项事务中使自己蒙他悦纳，好使你的丈夫因你而被称为有福，正如撒罗满的智慧所言：\Quote{谁能找到贤德的妇人？她的价值远胜珍珠。她丈夫的心信赖她，决不缺少利益；她一生使丈夫有益无损。她寻找羊毛和细麻，以双手劳作。她好像商船，从远方运来粮食。天未亮她就起来，把食物分给家中的人，将当作的工分派给使女。她想得田地，就买来；用手所得之利栽种葡萄园。她以能力束腰，使臂膀有力。她觉得经营有利；她的灯终夜不灭。她伸手帮助困苦人，周济穷人。她不因下雪为家人担心，因为全家都穿着朱红衣服。她为自己制作绣花毯子；她的衣服是细麻和紫色布做的。她丈夫在城门口与本地的长老同坐，为众人所认识。她做细麻布衣裳出卖，又将腰带卖给商人。她以能力和威仪为衣，想到日后的景况就喜乐。她开口就发智慧，仁慈的法律在她舌上。她观察家务，并不吃闲饭。她的儿女起来称她有福，她的丈夫也称赞她说：才德的女子很多，惟独你超过一切。艳丽是虚假的，美容是虚浮的；惟敬畏天主的妇人必得称赞。愿她享受操作所得的；愿她的工作在城门口使她受称赞。}\Quote{贤德的妻子是丈夫的冠冕。}又\Quote{许多女子建立了家室。}你已学到，一位明智而可爱的妻子从天主那里得到何等的赞美。妻子啊，你若愿成为信者之一，并悦乐主，就不要为取悦其他男子而装饰你的美貌；也不要追求精细的刺绣、衣服或鞋子，以引诱那些被此类事物所迷惑的人。因为即使你行这些恶事并非有意犯罪，而只是为装饰和美丽，你仍不能逃避将来的惩罚，因为你迫使别人如此注视你以致贪恋你，并且你没有谨慎避免自己犯罪及给人绊跌的机会。但若你屈服并犯下罪行，那么你既对自己的罪有份，也成了毁灭他人灵魂的原因。此外，当你与一人行淫后，开始绝望，便会再次偏离本分，跟随他人，变得麻木不仁；正如圣言所说：\BibleQuote{恶人陷入深渊，便藐视一切，随后羞辱和责备临到他。}\BibleRef{箴 18:3} 因为这样的妇人后来受伤，便毫无节制地陷害愚人的灵魂。因此，让我们学习圣言如何战胜这样的妇人，它说：\BibleQuote{我痛恨那比死亡更甚的、诱惑人心的妇人是网罗，她的手是锁链。}\BibleRef{训 7:26} 又在另一处说：\BibleQuote{妇人美貌而无见识，如同金环戴在猪鼻上。}\BibleRef{箴 11:22} 又说：\Quote{如蛀虫在木材中，恶妇毁坏自己的丈夫。}又说：\Quote{宁可住在房顶的角上，也不与争吵的妇人同住。}因此，身为基督徒的妇女，你们不要效法这类人。但你们这些愿忠于自己丈夫的人，务要小心只取悦他一人。你们在街上时，当蒙头；因为这样的遮盖能使你们避免被闲人注目。不可涂抹你们的面容，因为那是天主的化工；你们身上没有一处需要装饰，因为天主所造的一切都甚好。但轻浮地装饰已经好的事物，是对造物主恩赐的侮辱。你们外出时当眼目向下，蒙面而行，正符合妇女的体统。\par

% structure: p0024 source=h3 target=subsubsection reason=relative-heading-rank; base=h2

\subsubsection{妇女不可与男子共浴。}

IX. 也要避免与男子同浴的失序行为，因为恶者的网罗甚多。基督徒妇女不可与双性人同浴；因为若她必须蒙面，以端庄躲避陌生男子，她又怎能赤身与男子一同进入浴室呢？但若浴室专为妇女所用，则当有秩序地、端庄地、有节制地沐浴。然而不可无故沐浴，不可过多、过频，也不可在正午，如有可能，不可每日沐浴；第十时辰当为合宜沐浴的确定时间。因为身为基督徒妇女，你应当恒常避免那多眼的好奇心。\par

% structure: p0026 source=h3 target=subsubsection reason=relative-heading-rank; base=h2

\subsubsection{论好争吵及喧嚷的妇女。}

十、但至于争辩之心，务必对所有的人都抑制它，尤其是对你们的丈夫；免得他若是不信者或外教人，有了恶表的机会或亵渎天主，而你们也分受从天主来的祸。因为祂说：\Quote{祸哉，那因他而使我的名在外邦人中受亵渎的人！}另外，免得你们的丈夫若是基督徒，因着他所知道的圣经，被迫说出那写在\WorkTitle{智慧篇}中的话：\BibleQuote{住在旷野，比与一个好争论和易怒的妇人同住更好。}所以你们作妻子的，要借着你们的谦逊和温良，对教会外所有的人——无论男女——显示你们的虔敬，为使他们悔改并在信德上进步。我们既然已警告了你们，并简短地教导了你们，我们视你们为姐妹、女儿和肢体；你们自己既有智慧，终其一生坚持无可指责的生活。追求认识那使你们得以进入我们主的国、中悦祂的知识，并如此安息直到永远。阿们。\par

\SourceRecord{Translated by James Donaldson.；Ante-Nicene Fathers；Vol. 7.；Edited by Alexander Roberts, James Donaldson, and A. Cleveland Coxe.；Buffalo, NY: Christian Literature Publishing Co.,；1886.；<http://www.newadvent.org/fathers/07151.htm>.}

% structure: p0001 source=h1 target=section reason=page-title

\section{宗徒宪章（卷二）}

\Emph{论主教、司铎与执事。}\par

% structure: p0003 source=h2 target=subsection reason=relative-heading-rank; base=h2

\subsection{第一段　论候选主教之审核}

% structure: p0004 source=h3 target=subsubsection reason=relative-heading-rank; base=h2

\subsubsection{主教必须精于圣言并有经验。}

一、关于主教，我们从主那里获悉：凡在各堂区将要被祝圣为主教的牧者，必须无可指摘、无可责备，远离人间常见的各种邪恶，年岁不低于五十；因为这样的人大体上已摆脱了青年时期的放荡，以及外邦人的诽谤，有时某些假弟兄也会将责难加于许多人，他们不顾福音中天主的圣言：\BibleQuote{凡说废话的人，在审判之日，必向主交账。}\BibleRef{玛 12:36} 又说：\BibleQuote{凭你的话，要定你为义；凭你的话，也要定你有罪。}\BibleRef{玛 12:37} 因此，倘若可能，他应受过良好教育；若未受过教育，至少应娴熟于圣言，且年岁相当。但若在一个小堂区找不到年长者，则可选一位在邻里间有好名声、被他们视为堪当主教职务的年轻人——他自青年时代便以温良和规矩行事，如同一位年长者——经过考核和普遍好评后，可在平安中祝圣。因为撒罗满十二岁就作了以色列王，约史雅八岁就正义地执政\BibleRef{列下 22:1}，同样，约阿士七岁就治理百姓\BibleRef{编下 24:1}。因此，就算那人年轻，也当温良、和善、安静。因为主天主藉依撒意亚说：\BibleQuote{但我看顾谁呢？必看顾那谦卑和安静，且时时因我的话战栗的人。}\BibleRef{依 66:2} 同样在福音中也说：\BibleQuote{温良的人是有福的，因为他们要承受土地。}\BibleRef{玛 5:5} 他也当仁慈；因为又说：\BibleQuote{怜悯人的人是有福的，因为他们要受怜悯。}\BibleRef{玛 5:7} 他也当缔造和平；因为又说：\Quote{缔造和平的人是有福的，因为他们要称为天主的子女。} 他也当良心清白，远离一切邪恶、恶毒和不义；因为又说：\BibleQuote{心里洁净的人是有福的，因为他们要看见天主。}\BibleRef{玛 5:8}\par

% structure: p0006 source=h3 target=subsubsection reason=relative-heading-rank; base=h2

\subsubsection{主教及其余圣职人员应具备何种品格。}

二、因此，他应节制、明智、端庄、坚定、稳定，不嗜酒；不暴戾，却温和；不斗殴，不贪财；\BibleQuote{不可做新归化的人，免得他被骄傲冲昏，而落在魔鬼所受的判决和罗网里；因为凡自高自大的，必被贬抑。}\BibleRef{弟前 3:6} 主教应当是这样的人，他曾是\BibleQuote{只有一个妻子的丈夫}\BibleRef{弟前 3:2}，且那妻子也没有别的丈夫，\BibleQuote{善于管理自己的家}\BibleRef{弟前 3:4}。在他接受祝圣、就任主教职位之前，应如此考查：他是否庄重、忠信、端庄？他是否有一位庄重而忠信的妻子，或从前有过这样的一位？他是否虔敬地教育了子女，并\BibleQuote{以主的训诫和劝勉来教养他们}\BibleRef{弗 6:4}？他的家仆是否敬畏他、尊重他，并都服从他？因为若那些为世俗之事直接接近他的人尚且好斗不顺服，那么其他不属他家的人，当在他管理之下时，又怎能顺服他呢？\par

% structure: p0008 source=h3 target=subsubsection reason=relative-heading-rank; base=h2

\subsubsection{祝圣前应在哪些事上考查主教。}

三、还应考查他是否在今世的事务上无可指摘；因为经上记载：\Quote{要仔细查究那将要被祝圣为司祭之人的一切过失。}\par

% structure: p0010 source=h2 target=subsection reason=relative-heading-rank; base=h2

\subsection{第二段　论主教的品格与教导}

因此，他也应不轻易发怒；因为智慧说：\Quote{愤怒能毁灭智慧人。}他还当仁慈，宽厚、有爱心；因为我们的主说：\BibleQuote{如果你们彼此相爱，众人因此就可认出你们是我的门徒。}\BibleRef{若 13:35} 他也当乐于施舍，爱护寡妇和旅客；乐于服侍、服务和照顾；忠于职守；并知道谁最值得帮助。\par

% structure: p0012 source=h3 target=subsubsection reason=relative-heading-rank; base=h2

\subsubsection{慈善分配不应给予每位寡妇，有时有丈夫的妇女更应优先；任何贪饕、酗酒、懒惰者都不应得到分配。}

四、因为若有寡妇能自给，而另有一位妇女虽非寡妇，却因生病、抚育众多子女或双手无力而贫困，则他应以慈善之手援助后者。但若有人因贪饕、酗酒或懒惰而贫困，他就不配得任何援助，也不配被视为天主的教会的一员。因为圣经论及此类人说：\BibleQuote{懒惰人把手藏在怀里，甚至不能送回到口边。}\BibleRef{箴 19:24} 又说：\BibleQuote{懒惰人抱着双手，自食其肉。}\BibleRef{训 4:5} \Quote{因为凡酗酒的和纵欲的，必致贫穷；凡贪睡的人，必衣衫褴褛。}在另一处：\Quote{如果你把你的眼睛放在饮酒和杯盏上，你以后行走必比杵更赤裸。}确实，懒惰是饥荒之母。\par

% structure: p0014 source=h3 target=subsubsection reason=relative-heading-rank; base=h2

\subsubsection{主教在审判中不可徇情面；应有温和性情，生活节制。}

五、主教不可偏袒人；不可因敬畏或谄媚富人而违背正义，也不可忽视或欺压穷人。因为天主对\Person{梅瑟}说：\Quote{不可偏袒富人，也不可在穷人的案件中偏护穷人：因为审判是天主的。}\BibleRef{肋 19:15}又说：\Quote{你要完全按着公义追求公义。}主教应节俭，饮食知足，以便常保持清醒，乐于教导和劝诫无知者；不可饮食奢侈，纵情享乐，喜好美味。他应耐心温和地劝诫，自己深通经书，默思并勤习主的书卷，常常诵读，以便能正确解释圣经，按先知和法律诠释福音，且使法律和先知的诠释与福音一致。因为主\Person{耶稣}说：\Quote{你们查考经典，因为你们以为其中有永生，正是这些为我作证。}\BibleRef{若 5:39}又说：\Quote{\Person{梅瑟}曾经论及我而写作。}\BibleRef{若 5:46}但最重要的是，他应仔细区分原初的法律和附加的规条，指出哪些是信徒的法律，哪些是不信者的束缚，以免任何人陷入那些束缚。因此，主教啊，你要谨慎研读圣言，以便能精确解释一切，并饱以丰富的道理养育你的子民，以法律的光明照亮他们；因为天主说：\Quote{趁着还有时间，你们要以知识的光明照亮自己。}\BibleRef{欧 10:12}\par

% structure: p0016 source=h3 target=subsubsection reason=relative-heading-rank; base=h2

\subsubsection{主教不可贪图不义之财，亦不可作保人或律师。}

VI. 主教不可贪图不义之财，尤其在外邦人面前，宁可受辱也不可加害于人；不可贪婪，不可强夺；不可偷窃；不可奉承富人，不可憎恨穷人；不可毁谤，不可作假见证；不可轻易发怒；不可争吵；不可被俗世事务缠身；不可为任何人作保，也不可在钱财诉讼中作原告；不可虚荣；不可心持两端，也不可一口两舌；不可轻易听信谗言或毁谤；不可虚伪；不可参与异教节庆；不可追求虚妄的欺骗；不可贪恋世俗之事，也不可贪爱钱财。因为这一切都与天主相悖，令魔鬼喜悦。主教也当恳切地将这一切诫命交付给平信徒，劝勉他们效法自己的行为。因为主说：\BibleQuote{你们要使以色列子民成为圣洁。}\BibleRef{肋 15:31} 他应当明智、谦卑、善于以主的训诲劝戒人，心地善良，弃绝一切世俗的邪恶计谋和异教的贪欲；应当有秩序、敏锐地观察并防备恶人，但仍作众人的朋友；正义、明辨；凡在人间可称赞的品德，主教都当自身具备。因为如果牧者自身在一切恶事上无可指责，他必能迫使自己的门徒，并藉自己的生活方式促使他们成为自己行为的可贵效法者。正如先知某处所说：\BibleQuote{司铎如何，百姓也如何。}\BibleRef{欧 4:9} 因为我们的主、导师、天主子耶稣基督，祂先实行，然后教导，正如路加某处所说：\BibleRef{宗 1:1}\BibleQuote{耶稣开始所行和所教导的。}因此祂说：\BibleQuote{谁若实行并教导，他在天国里将称为大的。}\BibleRef{玛 5:19} 因为你们主教是人民的向导和守望者，正如你们自己有基督为你们的向导和守望者。所以你们要成为天主子民的好向导和守望者。因为主藉厄则克耳对你们每个人说：\Quote{人子，我立你为以色列家族的守望者；你要听我口中的话，并要留意，从我这里宣告。当我向恶人说：你必死；你若不警告他离开恶道，那恶人必因自己的罪孽而死，我却要向你追讨他的血债。你若警告恶人离开他的道，使他转离，他却不转离，他必因自己的罪孽而死，你却救了你的灵魂。}\newline{}\Quote{同样，若有战争的刀剑临近，人民设立一个守望者，他看见刀剑临近，却不警告他们，刀剑临到，夺去他们中的一人，那人虽因自己的罪孽被夺去，但我要向守望者追讨他的血债，因为他没有吹号角。但若他吹了号角，听见的人不受警告，刀剑临到夺去他，他的血就归在他自己头上，因为他听见号角却不受警告。但受警告的人救了灵魂；守望者因发出了警告，必定存活。}这里的刀剑是审判，号角是神圣的福音，守望者是主教，他被设立在教会中，有责任藉宣讲作证并严厉警告那审判。若你不向人民宣告并作证，那些不知之人的罪就要归于你。因此你们要勇敢地警告和责备未受教导的人，教导无知的人，坚固有理解的人，领回迷途的人。弟兄们，若我们在同一场合重复同样的事，并不算错。因为屡次听闻，有些人可望感到羞愧，至少行些善事，避免恶行。因为天主藉先知说：\Quote{向她们作证这些事，也许她们会听你的声音。}又说：\Quote{也许他们会听，也许他们会顺从。}梅瑟也对人民说：\BibleQuote{如果你们听从上主你们的天主，行祂眼中看为善为正的事。}\BibleRef{出 15:26} 又说：\BibleQuote{以色列，请听：上主我们的天主是唯一的上主。}\BibleRef{申 6:4} 我们的主在福音中也屡次记载说：\Quote{有耳可听的，听吧！}智慧的撒罗满说：\BibleQuote{我儿，要听你父亲的训诲，不可离弃你母亲的法则。}\BibleRef{箴 1:8} 然而直到今日，人仍未听从；因为他们虽似听了，却没有听对，正如他们离弃了独一的真天主，被引入毁灭性的危险异端中，对此我们日后还要再谈论。\par

% structure: p0018 source=h2 target=subsection reason=relative-heading-rank; base=h2

\subsection{第三部分  主教应如何对待无辜者、有罪者及悔罪者}

% structure: p0019 source=h3 target=subsubsection reason=relative-heading-rank; base=h2

\subsubsection{受洗者应有的品格}

VII. 亲爱的，你们要知道，那些受圣洗归于我们主耶稣之死的人，有义务不再继续犯罪；因为死人不能再行恶，同样，与基督同死的人也不能作恶。因此，弟兄们，我们不相信任何领受了生命之洗的人仍继续犯放荡的罪行。凡在圣洗之后犯罪者，除非悔改并离弃其罪，否则将被判定下地狱之火。\par

% structure: p0021 source=h3 target=subsubsection reason=relative-heading-rank; base=h2

\subsubsection{关于被诬告者或已被定罪者}

VIII. 但若有人因外教人的恶意迫害，因为他不再与他们同流合污，放纵无度，就当知道，这样的人是天主所祝福的，正如我主在福音中所说：\Quote{几时人为了我的缘故辱骂你们，迫害你们，捏造一切坏话毁谤你们，你们是有福的。你们欢喜踊跃罢，因为你们在天上的赏报是丰厚的。}\BibleRef{玛 5:11} 因此，若有人被毁谤、被诬告，这样的人是有福的；因为经上说：\Quote{被弃绝的人，天主不试炼他。} 但若有人被证实行了恶事，这样的人不仅伤害自己，还使整个教会身体及其道理受到亵渎，好像我们基督徒不实行自己所宣称的美好正当事，我们自己也要被主责备：\Quote{因为他们说而不做。}\BibleRef{玛 23:3} 因此，主教在充分定罪后，必须勇敢地拒绝这样的人，除非他们改变生活。\par

% structure: p0023 source=h3 target=subsubsection reason=relative-heading-rank; base=h2

\subsubsection{主教不应收受贿赂。}

IX. 因为主教不仅自己不可得罪人，也不可徇私；要以温柔教导那些犯罪的人。但如果他自己良心不正，且因贪图不义之财和收受贿赂而徇私，宽容公开的罪人，容他留在教会中，他便无视天主和我们主的声音，祂说：\Quote{你们要施行准确的审判。}\Quote{审判时不可徇私：不可使不敬虔的人成义。}\Quote{不可因任何人接受礼物；因为礼物蒙蔽智慧人的眼目，歪曲义人的言语。}\BibleRef{出 23:8} 祂在别处又说：\Quote{从你们中间除掉那恶人。} 撒罗满也在《箴言》中说：\Quote{从会众中赶出那捣乱的人，争吵就与他一同离开。}\BibleRef{箴 22:10}\par

% structure: p0025 source=h3 target=subsubsection reason=relative-heading-rank; base=h2

\subsubsection{主教若因错误的判决而宽恕罪人，自己也有罪。}

X. 但那些不考虑这些事的人，会违背正义，宽恕该受惩罚的人，就如撒乌耳宽恕了阿加格，厄里宽恕了他的儿子们，他们\Quote{不认识上主}。这样的人亵渎了自己的尊严，也亵渎了在他辖区的天主教会。这样的人在天主和圣洁的人面前被视为不义，因为他给许多新受洗的人和慕道者，以及男女青年提供了绊脚石，关于这些人有祸了，且\Quote{把一块磨石拴在他的颈上，}\BibleRef{玛 18:6} 使他沉入海底，因他的罪过。因为看到他们的治理者是怎样的，因他的邪恶和对正义的疏忽，他们会变得怀疑，并沉溺于同样的毛病，被迫与他一同灭亡；如同与雅洛贝罕联合的百姓，以及与科辣黑同谋的人。但如果罪人看到主教和执事是清白无辜的，羊群是纯洁的，他要么不敢藐视他们的权威，而根本不敢进入天主的教会，如同被自己的良心打击；要么他若无所顾忌，胆敢进入，就会立刻被定罪，如同乌匝在约柜旁伸手扶它，以及阿干偷窃了当灭之物，又如革哈齐贪图纳阿曼的银钱，立刻受罚；要么他会受到牧者的劝诫，被引向悔改。因为当他逐个审视整个教会，在主教和他所照管的百姓中看不到任何瑕疵，他便会困惑，心中刺痛，然后谦和地带着羞愧和许多眼泪离开，羊群得以保持纯洁。他将流着泪转向天主，悔改自己的罪，并怀有希望。是的，整个羊群看到他的眼泪，就会受到教训，因为罪人借着悔改避免了灭亡。\par

% structure: p0027 source=h3 target=subsubsection reason=relative-heading-rank; base=h2

\subsubsection{主教应如何审判罪人。}

XI. 因此，主教啊，要努力在你的行为上纯洁，并装饰你的职位和尊严，这是代天主在人间的角色，因为你被安置在众人之上，包括司铎、君王、统治者、父亲、儿女、教师，以及一般所有服从你的人：因此当你说话时，要坐在教会中，如同有权审判罪人。因为对你们主教们说：\Quote{凡你们在地上所捆绑的，在天上也要被捆绑；凡你们在地上所释放的，在天上也要被释放。}\BibleRef{玛 18:18}\par

% structure: p0029 source=h3 target=subsubsection reason=relative-heading-rank; base=h2

\subsubsection{关于主教应如何对待忏悔者的训示。}

XII. 因此，主教啊，要像天主一样以权威审判，但要接纳忏悔者；因为天主是慈悲的天主。责斥犯罪的人，劝诫不悔改的人，鼓励坚持善行的人，接纳忏悔者；因为上主天主曾起誓应许赐予忏悔者赦免他们所行的过犯。因为祂藉厄则克耳说：\Quote{你要对他们说：我指着我的永生说，吾主上主决不喜欢恶人死亡，而喜欢恶人离开他的恶道，得以生存。以色列家啊，你们转离你们的恶道罢！何必死亡呢？}\BibleRef{则 33:11} 这里的话给罪人带来希望，如果他们悔改，就有救恩的希望，免得他们因绝望而屈服于过犯；而是因着救恩的希望，他们可以悔改，流着泪向天主祈求自己的罪，从心里悔改，平息祂对他们的愤怒；这样他们就会从祂那里得到宽恕，如同从慈悲的父亲那里一样。\par

% structure: p0031 source=h3 target=subsubsection reason=relative-heading-rank; base=h2

\subsubsection{我们应谨慎，不可试探任何罪恶的道路。}

XIII. 然而，那些尚属无辜者非常需要保持这种状态，不要尝试何为罪，以免他们需要经历烦恼、忧伤以及为了获得宽恕而有的哀哭。因为，人啊，当你犯罪时，你如何知道你在今生还能存活多少日子，使你有时机悔改？因为你离开此世的时刻是不确定的；如果你死在罪中，那么为你就不再留有悔改的机会；正如天主藉\Person{达味}所说：\Quote{在坟墓里，有谁承认祢？}因此，我们理应在尽本分上准备好，这样我们才能无忧地等待进入另一个世界。所以，圣言也藉着明智的\Person{所罗门}劝诫你说：\BibleQuote{在你的出口之先预备你的工作，在田间提前备好一切，}\BibleRef{箴 24:27}免得你旅途所需有所欠缺；正如福音中所提到的五个糊涂童女，她们因着熄灭了自己神圣知识的灯，而被关在新郎的婚筵之外。所以，那珍视自己灵魂安全的人，将小心地脱离危险，保持无罪的自由，这样他就能为自己保存先前善功的益处。因此，你要如同执行天主的审判一样来审判。因为经上说：\BibleQuote{审判是属于主的。}\BibleRef{申 1:17}所以，首先，要权柄地定那有罪之人的罪；然后，要努力以仁慈、怜悯和接纳他的意愿，将他带回家，如果他要改变自己的生活方式并成为悔改者，就应许他救恩；当他悔改并接受责罚后，就要接纳他：要记念我们的主曾说：\BibleQuote{对一个悔改的罪人，天堂也要欢乐。}\BibleRef{路 15:7}\par

% structure: p0033 source=h3 target=subsubsection reason=relative-heading-rank; base=h2

\subsubsection{论那些声称不可接纳悔改者进入教会的人。论一个义人，即使他与罪人来往，也不会与他一同丧亡。论没有人会因他人而受罚，但每个人必须为自己的行为交账。论我们必须扶助信德软弱的人；并论主教不可受在俗人中任何滋事之人的支配。}

XIV. 倘若你拒绝接受那悔改的人，你便将他暴露给那些埋伏毁灭的人，忘记了\Person{达味}所说的：\BibleQuote{不要将向祢认罪的我灵魂交付给毁灭的野兽。}因此，\Person{耶肋米亚}在劝人悔改时说：\BibleQuote{那跌倒的，岂不能起来？那转身的，岂不能回来？我的百姓为何以无耻的退转而退后？他们又硬着心。\BibleRef{耶 8:4} 转回来吧，悖逆的儿女，我要医治你们的悖逆。\BibleRef{耶 3:22}}因此，毫无疑虑地接受那悔改的人。不要被那些不仁慈的人阻止，他们说我们不可被这样的人玷污，甚至不可与他们说话：因为这样的建议是来自那些不认识天主及其眷顾的人，是不讲理的审判者，是不仁慈的畜生。这些人不知道，我们应当避免与犯罪者交往，不是言语上，而是行动上：因为\BibleQuote{义人的义归于自己，恶人的恶归于自己。\BibleRef{则 18:20}}又说：\BibleQuote{若一国犯罪得罪我，严重背叛，我伸手攻击它，折断它的粮杖，使饥荒临到它，灭绝其中的人和牲畜：即使有\Person{诺厄}、\Person{约伯}、\Person{达尼尔}这三人在其中，他们只能因自己的义救自己的灵魂——这是主天主说的。\BibleRef{则 14:13}}\WorkTitle{圣经}最清楚地表明，一个义人与恶人交往，并不与他一同灭亡。因为在现世，义人和恶人在日常生活的交往中混杂在一起，但不是在神圣的共融中：在这方面，天主的朋友和宠儿并不犯罪。因为他们只是效法\BibleQuote{他们在天的父，祂使太阳上升，照义人也照不义的人，降雨给义人也给不义的人；\BibleRef{玛 5:45}}义人并不因此遭遇危险。因为得胜的和失败的都在同一赛场上，但只有勇敢跑完赛程的人才能得到花冠；并且\BibleQuote{若不按规矩竞技，就不能得到花冠。\BibleRef{弟后 2:5}}因为各人都要为自己交账，天主不会将义人与恶人一同毁灭；因为祂有一个不变的法则，即无辜者永不受罚。因为祂既没有使\Person{诺厄}与世人一同淹死，也没有烧死\Person{罗特}与索多玛人，也没有毁灭\Person{辣哈布}与耶里哥居民同尽。你若想知道这事在我们中间如何，\Person{犹达斯}曾是我们中的一员，与我们同享职分；\Person{西满}术士也接受了主的印记。但两人都显为邪恶，前者上吊自尽，后者违反自然地在空中飞行时被摔在地上。此外，\Person{诺厄}和他的儿子们都在方舟里；但只有\Person{含}被发现有罪，在他的儿子身上受了惩罚。但如果父亲不因子女受罚，子女也不因父亲受罚，那么显然妻子不因丈夫受罚，仆人不因主人受罚，一个亲属不因另一个亲属受罚，一个朋友不因另一个朋友受罚，义人不因恶人受罚。但各人都要为自己的行为交账。因为\Person{诺厄}没有为世界受罚，\Person{罗特}没有为索多玛人遭火烧，\Person{辣哈布}没有为耶里哥居民被杀，以色列也没有为埃及人受罚。因为不是同住，而是思想上的认同，才会使义人与恶人一同被定罪。因此，我们不可听从那些呼求死亡、憎恨人类、喜爱指控、以漂亮借口置人于死地的人。因为一人不会为另一人而死，但\BibleQuote{各人被自己的罪恶绳索缠住。\BibleRef{箴 5:22}}还有：\BibleQuote{看哪，那人和他的作为都在他面前。\BibleRef{依 62:11}}如今我们应当帮助那些与我们同在、处于危险中、跌倒的人，并尽我们所能，以我们的劝诫使他们醒悟，从而救他们脱离死亡。因为\BibleQuote{健康的人不需要医生，而是有病的人；\BibleRef{玛 9:12}}因为\BibleQuote{你们在天之父不愿意这些小子中的一个丧亡。\BibleRef{玛 18:14}}因为我们不应确立心硬之人的意愿，而应确立宇宙的天主和圣父的意愿，这意愿是由我们的主耶稣基督启示给你们的，愿光荣归于祂，直到永远。阿们。\par

因为这不公平，你身为主教，是头，竟向尾屈从，即向在平信徒中某个煽动叛乱的人屈从，以致毁灭他人；你只应服从天主。因为治理下属是你的特权，而非受他们治理。因为儿子因生育过程服从父亲，却不管辖父亲；奴隶因法律服从主人，却不管辖主人；学生不管辖老师，士兵不管辖国王，任何平信徒不管辖他的主教。因为若有人与恶人交谈，为以圣言教导他们，便以为会被玷污或共担他们的罪，这是毫无理由的；\OriginalTerm{厄则克耳}{Ezekiel}仿佛有意阻止心术不正者的猜疑，这样说：\Quote{你们在以色列地方怎么说这句俗语：父亲吃了酸葡萄，儿子的牙酸倒呢？主天主说：我指着我的永生起誓，你们在以色列中，必不再有这俗语的用途。因为所有的灵魂都属于我；父亲的灵魂怎样，儿子的灵魂也怎样，都一样；犯罪的灵魂，它必死亡。但人若是义人，行公道正义}（先知接着列举其余德行，然后总结说：\Quote{这样的人是义人}），\Quote{他必因自己的义德而生活，这是主天主说的。他若生一个儿子，是强盗，是流人血的，不行他义父的道}（先知接着说了下面的话，又总结说）：\Quote{他必不能存活；他行了这一切恶事，他必要死亡；他的血要归在他身上。他们还要问你说：为什么？儿子不承受父亲的罪孽吗？或是他自己的义德，因他自己行了正义和仁慈？你对他们说：犯罪的灵魂，它必死亡。父亲不承受儿子的罪孽，儿子也不承受父亲的罪孽。义人的义德归在自己身上，恶人的邪恶也归在自己身上。}稍后他说：\Quote{义人若转离他的义德，而行恶事，他所行的一切义德，因他所行的一切恶事，都不再被记念；他因所行的恶事，和他所犯的罪，必要死亡。}稍后又加说：\Quote{恶人若转离他所行的恶事，行公道正义，他保存了自己的灵魂，他转离了他所行的一切不敬；他必存活，不至死亡。}后来又说：\Quote{主天主说：以色列家啊，我必按你们各人的行为审判你们。}\par

% structure: p0036 source=h3 target=subsubsection reason=relative-heading-rank; base=h2

\subsubsection{司铎不可忽视过犯，也不可急于惩罚}

十五．你们这些我们亲爱的儿子，请看我们的主天主是何等仁慈而又公义；何等恩慈良善，却必然\Quote{不豁免有罪的人}：\BibleRef{鸿 1:3}尽管他欢迎悔改的罪人，使他重生，不留下任何怀疑的余地给那些愿意严厉审判、全然拒绝罪人、不肯赐予他们能引向悔改的劝勉的人。针对这样的人，天主借\OriginalTerm{依撒意亚}{Isaiah}对主教们说：\Quote{安慰，安慰我的百姓，你们司祭；向耶路撒冷说宽容的话。}因此，你们听了他的话，就该鼓励那些犯罪的人，引导他们悔改，赐予他们希望，不要徒然以为因着你们对他们的爱，你们就会分担他们的过犯。要欣然接纳悔改者，为他们喜乐，以慈悲和怜悯的心肠审判罪人。因为若有人走在河边，快要绊倒，你非但不出手相助，反而推他跌入河中，你就成了杀害你弟兄的凶手；你本该在他将要跌倒时伸出援手，免得他无可挽救地灭亡，这样既让百姓得警戒，又使犯罪者不至完全丧亡。主教，你的责任既不可忽视百姓的罪，也不可拒绝悔改者，免得你不谙熟地摧毁主的羊群，或玷污他赐予他百姓的新名，而你自己受到斥责，如同那些古时的牧人，天主曾借\OriginalTerm{耶肋米亚}{Jeremiah}如此论及他们：\Quote{许多牧人毁坏了我的葡萄园，践踏了我的产业。}\BibleRef{耶 12:10}又在另一处说：\Quote{我的愤怒向牧人燃起，对羊群我必要惩治。}\BibleRef{匝 10:3}又在别处说：\Quote{你们是亵渎我名的司祭。}\BibleRef{拉 1:6}\par

% structure: p0038 source=h3 target=subsubsection reason=relative-heading-rank; base=h2

\subsubsection{论悔改，其方式及规则}

第十六、当你看见犯罪者时，要严厉命令将他驱逐出去；当他出去时，也要让执事严厉对待他，然后他们出去寻找他，把他拘留在教会之外；他们进来时，让他们为你求情。因为我们的救主自己也曾为犯罪者向祂的圣父求情，正如福音所记载：\BibleQuote{父啊，宽赦他们，因为他们不知道他们所做的是什么。}\BibleRef{路 23:34} 然后命令犯罪者进来；如果你经查问发现他已痛悔，并适合被接回教会，就按他罪过的轻重——如两、三、五或七周——使他受苦守斋后，释放他，并对他说些适合责斥、教训和劝勉罪人改过的话，好使他私下继续谦卑，祈求天主怜悯他，说：\Quote{主，你若细察罪孽，谁能站立得住？因为你有宽宥。} 在 \WorkTitle{创世纪} 中对该隐所说的话也属于这类宣告：\Quote{你犯了罪；安静吧；}即不要继续犯罪。因为罪人应当为自己的罪羞愧，天主给梅瑟关于米黎盎的神谕足以证明这一点，当时梅瑟祈求她得到宽赦。因为天主对他说：\Quote{如果她的父亲吐唾沫在她脸上，她岂不应当羞愧？让她在营外关锁七天，然后才让她进来。} 因此，当犯罪者表示悔改时，我们也应当这样做——即按他们罪过的比例，将他们隔离一段确定的时间，然后像父亲对待儿女一样，在他们的悔改后重新接纳他们。\par

% structure: p0040 source=h3 target=subsubsection reason=relative-heading-rank; base=h2

\subsubsection{主教必须无可指责，并作他属下的榜样}

第十七、但如果主教本身是犯罪者，他如何还能追究别人的罪过？或者他如何能责备别人，无论是他本人或是他的执事，若他们因徇私或受贿，而良心都不清白？因为当官长求问，审判官受贿时，审判不能达到完美；但若两者都是\BibleQuote{贼的伙伴，不顾给寡妇伸冤}\BibleRef{依 1:23}，那些在主教训导之下的人便不能支持或维护他：因为他们会对他说福音所记载的：\BibleQuote{为什么你看见你弟兄眼中的木屑，却不理会自己眼中的大梁？}\BibleRef{路 6:41} 因此，主教与其执事应当惧怕任何这样的事；即不要给它们机会。因为犯罪者若看见别人也做同样坏的事，就会被鼓励去做同样的事；于是恶者从单一事例找到机会，在别人身上运作——愿天主禁止——这样羊群将被毁灭。因为犯罪者越多，他们所作的祸害也越大：因为未受纠正的罪愈演愈烈，且蔓延到别人；因为\BibleQuote{一点面酵能使整个面团发酵}\BibleRef{迦 5:9}，一个小偷能使整个民族染上可憎的事，\BibleQuote{死苍蝇使香膏变臭}\BibleRef{训 10:1}，\BibleQuote{君王若听从谎言，他的一切臣仆都是邪恶的}\BibleRef{箴 29:12}。同样，一只疥癣的羊，若不从健康的羊中分离，会将同样的疾病传染给其余；一个染上瘟疫的人，人人都应躲避；一条疯狗，对它所触碰的每个人都是危险的。因此，如果我们忽视将违法者从天主的教会中分离，我们将使\BibleQuote{主的殿成为贼窝}\BibleRef{玛 21:13}。因为主教的责任不是在犯罪者面前沉默，而是责备他们、劝诫他们、降伏他们、以守斋磨难他们，好使其余人产生虔诚的敬畏：因为正如祂所说：\BibleQuote{使以色列子民成为虔诚的}\BibleRef{肋 15:31}。因为主教必须是一个以劝诫阻止罪恶的人，树立义德的榜样，宣讲天主所预备的美善，并宣告在审判之日将要降临的忿怒，免得他轻视并忽略天主的种植；并且由于他的疏忽，听到 \WorkTitle{欧瑟亚书} 中所说的话：\Quote{你为何对不义默不作声，并收取其果实？}\par

% structure: p0042 source=h3 target=subsubsection reason=relative-heading-rank; base=h2

\subsubsection{主教必须注意使其人民不犯罪，因他被立为他们中的守望者}

XVIII. 因此，主教应将他的关怀扩展到各种人：对于那些没有犯罪的人，使他们继续保持无罪；对于那些犯罪的人，使他们悔改。因为主对你说：\BibleQuote{你们小心，不要轻视这些小子中的一个。}\BibleRef{玛 18:10} 你也有责任向忏悔者赐予赦免。因为一旦犯罪者真诚地发自灵魂说：\Quote{我得罪了主，}圣神就回答说：\Quote{主也赦免了你的罪；放心，你不会死。}因此，主教啊，请意识到你职位的尊荣，既然你领受了束缚的权力，也领受了释放的权力。既然有释放的权力，要认识你自己，并在今世举止合宜，知道你要交出一份巨大的账目。因为，如圣经所说：\BibleQuote{托付谁多，就向谁索取更多。}\BibleRef{路 12:48} 因为没有人是无罪的，除了那为我们而成为人的那一位；经上记载：\Quote{无人不受污秽；即使只有一天。}为此，古代圣人和族长的生活和行为被描述出来；不是要我们因阅读而指责他们，而是要我们自己悔改，并希望我们也能获得赦免。因为他们的瑕疵对我们既是保障也是劝诫，因为我们由此学习，当我们犯罪时，如果悔改，就必得赦免。经上记载：\BibleQuote{谁能夸口说自己有洁净的心？谁敢断言自己无罪？}\BibleRef{箴 20:9} 因此，无人无罪。所以，尽你最大的努力成为无可指摘的人；并关心你羊群中的每一部分，免得有人因你而跌倒，因而丧亡。因为平信徒只关心自己，而你却关心众人，因为你肩负更重的负担，背负更沉的担子。经上记载：\BibleQuote{主对梅瑟说：你和亚郎要担当司祭职的罪孽。}\BibleRef{户 18:1} 既然你要为众人交账，就关心众人吧。保护那些健全的，劝诫那些犯罪的；当你用禁食惩罚他们后，用赦免使他们舒缓；当犯罪者流泪请求重新接纳时，接纳他，让整个教会为他祈祷；当你通过覆手接纳他后，允许他此后留在羊群中。但对于那些昏睡和疏忽的人，要努力劝化、坚固、警告和医治他们，要意识到这样做会获得多么大的赏报，而若疏忽则会招致多大的危险。因为厄则克耳对那些不关心人民的监督者这样说：\Quote{祸哉以色列的牧人！他们只牧养自己；牧人不牧养羊，只牧养自己。你们喝奶，穿羊毛，宰杀肥壮的，却不牧养羊。瘦弱的你们没有养壮，有病的你们没有医治，受伤的你们没有包扎，被赶散的你们没有领回，迷失的你们没有寻找；却用暴力严厉惩罚他们；羊因没有牧人而四散，成了田野一切野兽的食物。}他又说：\Quote{牧人没有寻找我的羊；牧人牧养自己，却不牧养我的羊。}稍后又说：\Quote{看，我要攻击那些牧人，从他们手中追讨我的羊，使他们不再牧养我的羊，牧人也不再牧养自己；我要从我手中救出我的羊，不再作他们的食物。}他还对人民说：看，我必在羊与羊之间，在公绵羊与公山羊之间施行审判。你们吃了肥美的草场，却用脚践踏剩下的草场，羊就吃了你们脚践踏的，这在你们眼中是小事吗？稍后他又说：\Quote{你们要知道我是主，你们是我草场的羊；你们是我的人，我是你们的天主——上主天主说。}\par

% structure: p0044 source=h3 target=subsubsection reason=relative-heading-rank; base=h2

\subsubsection{一个疏忽自己羊群的牧人将被定罪，而不愿被牧人引领的羊应受惩罚。}

XIX. 主教们，请听；平信徒们，请听，天主是如何说的：\Quote{我必在公羊与公羊之间，在羊与羊之间施行审判。}他对牧人说：\Quote{你们必因你们的无能，并因毁灭羊群而受审判。}也就是说，我必在一位主教与另一位主教之间，在一位平信徒与另一位平信徒之间，在一位统治者与另一位统治者之间施行审判（因为这些羊和这些公羊不是无理性的，而是有理性的受造物）：免得有朝一日，一个平信徒说：我是羊，不是牧人，我不管自己；让牧人管吧，因为他将独自为我交账。因为正如那只不跟随好牧人的羊，暴露在狼群中，走向毁灭；同样，那只跟随坏牧人的羊，也面临无法避免的死亡，因为它的牧人会吞噬它。因此，必须留心避免毁灭性的牧人。\par

% structure: p0046 source=h3 target=subsubsection reason=relative-heading-rank; base=h2

\subsubsection{被治理者应如何服从派定管理他们的主教。}

二十．至於善牧，平信徒應尊敬他、愛戴他、敬畏他，如同自己的主、自己的師傅、天主的大司祭、虔敬的導師。因為聽他的，就是聽基督；拒絕他的，就是拒絕基督；不接受基督的，就不接受祂的天主和父。因為祂說：\BibleQuote{聽你們的，就是聽我；拒絕你們的，就是拒絕我；拒絕我的，就是拒絕那派遣我來的。}\BibleRef{路 10:16}同樣，主教應愛平信徒如同自己的兒女，以溫情的熱誠培育他們、珍愛他們；如同雞蛋，為了孵化雛鳥；或如雛鳥，抱在懷中，養育成鳥。要勸誡眾人；責備一切需要責備的人；責備，而非打擊；使他們羞愧而折服，卻不推翻他們；警告他們，使之悔改；叱責他們，使之改善和走上更好的生活道路。要照顧強壯的，即堅固那已在信德上堅強的人；和平地牧養百姓；堅固軟弱的，即用勸勉鼓勵那受試探的；醫治有病的，即用教導醫治那因疑慮而在信德上軟弱的人；包紮破碎的，即用安慰的勸誡包紮那因罪而迷失、受傷、壓傷、破碎、偏離正路的人；解除他們的過犯，給予希望；藉此使他們恢復力量，歸回教會，領回羊群。領回那被驅逐的，即不可允許那因罪惡而被懲罰逐出的人長久被隔離；而要接納他，領回他，使他歸回羊群，即歸回無瑕教會的子民。尋找那迷失的，即不可讓那因罪過眾多而對救恩絕望的人完全喪亡。尋找那沉睡、昏沉、懶惰的，以及那因沉睡過深而忘卻自己生命、離自己羊群甚遠、有落入狼群被吞噬之險的人。用勸誡帶回他，激勵他警醒；灌輸希望，不許他說某些人所說的：\Quote{我們的罪惡壓在我們身上，我們在它們中消瘦，我們怎能存活？}因此，主教要盡可能視罪過為己有，對罪人說：你只管回頭，我願為你承受死亡，正如我們的主為我和所有人承受死亡一樣。因為\BibleQuote{善牧為羊捨命；但僱工，不是牧人，羊不是他自己的，看見狼來了——就是魔鬼——就撇下羊逃跑，狼就抓住羊。}\BibleRef{若 10:11}因此，我們必須知道，天主對犯罪的人極其慈悲，且發誓應許悔改。但犯罪的人若不認識天主關於悔改的應許，不明白祂的寬容和忍耐，更因從未從你學習而不曉得宣講悔改的聖經，就因自己的愚昧而喪亡。但——如同慈悲的牧人、殷勤餵養羊群的人——你要尋找並數點你的羊群。尋找那迷失的，\BibleRef{瑪 18:12}如同主天主我們仁慈的父，派遣了自己的兒子，善牧和救主，我們的主耶穌，命令祂\Quote{離開那九十九隻在山上，去尋找那迷失的，找到後，扛在肩上，帶回羊群，因找到迷失的而歡喜。}同樣，主教啊，你要順從，尋找那迷失的，引導那偏離正路的，領回那走入歧途的：因為你有權柄領回他們，並藉赦免釋放那些心碎的人。因為藉著你，我們的救主對那因罪惡感到沮喪的人說：\Quote{你的罪赦了；你的信德救了你；平安去吧。}而這個平安和安穩的港灣就是基督的教會；你一旦釋放了他們的罪，就將他們領回教會，使他們如今健康無可指摘，懷有希望，勤奮，熱心行善。如同一位熟練而仁慈的醫生，醫治一切在罪惡之路上迷失的人；因為\BibleQuote{健康的人不需要醫生，而是有病的人。因為人子來，是為拯救和尋找那迷失的。}\BibleRef{瑪 9:12}既然你是主的教會的醫生，就要為每個病人的情況提供合適的療法。醫治他們，用一切可能的方法治好他們，使他們健康地歸回教會。牧養羊群，\BibleQuote{不要傲慢輕蔑，如同作威作福}，\BibleRef{則 34:4}而要如同溫良的牧人，\BibleQuote{將羔羊聚集在懷中，溫柔地引導那懷孕的。}\BibleRef{瑪 20:25}\par

% structure: p0048 source=h3 target=subsubsection reason=relative-heading-rank; base=h2

\subsubsection{不聽取雙方陳述就判斷，或在一個人被定罪之前就決定懲罰，是危險的事。}

XXI. 要温柔、和蔼、温和、不诡诈、不虚伪；不固执、不无礼、不严厉、不傲慢、不无情、不自大、不谄媚、不怯懦、不三心二意、不凌辱属下的百姓、不隐瞒神圣的法律和悔改的应许、不急于驱逐和开除，而要稳健，不喜欢严酷，不鲁莽。不要接受少于三个证人的证据来判定任何人，而且这些证人必须是名声可靠的人；要查问他们是否出于恶意或嫉妒而指控；因为有许多人喜欢作恶、言语轻率、诽谤、憎恨弟兄，以驱散\Person{基督}的羊群为己任；你若不经仔细审查就接受他们的证词，便会分散你的羊群，任凭它们被狼吞噬，即被魔鬼和恶人吞噬，甚至不是人，而是披着人形的野兽——被外邦人、犹太人和无神论的异端者吞噬。因为那些毁灭性的狼很快会向任何被赶出教会的人下手，视他为交付给它们吞噬的羔羊，把他的毁灭当作自己的获益。因为那\BibleQuote{他们的父亲魔鬼是杀人者}。\BibleRef{若 8:44}凡因你判断疏忽而被不公地分离的，将充满忧伤，沮丧，从而要么流落到外邦人那里，要么陷入异端，最终完全与教会和天主的希望隔绝，陷于不敬之中，你便成了他丧亡的罪魁祸首；因为不应对犯罪者急于赶出，而在他回头时缓慢接纳；反倒急于切断，在他哀伤并应被医治时却无情。因为神圣的圣经论及这样的人说：\BibleQuote{他们的脚奔跑作恶；他们急于流人血。毁灭与悲惨在其路上，和平之道他们不曾知晓。天主的敬畏不在他们眼前。}然而和平之道乃是我们的救主\Person{耶稣基督}，祂教导我们说：\BibleQuote{你们宽恕，也必蒙宽恕；你们给予，也必给予你们；}\BibleRef{路 6:37}也就是说，赐予罪的赦免，你们的过犯也必得赦免。同样，祂也以祈祷教导我们向天主说：\BibleQuote{宽免我们的罪债，如同我们也宽免得罪我们的人。}\BibleRef{玛 6:12}因此，若你不宽恕得罪你的人，怎能期望自己的罪得赦免？岂不因你在祈祷中假装宽恕而实际不宽恕，反使自己更受束缚？当你口说宽恕却实际不宽恕时，岂不会被自己的话质问？因为要知道，那赶出行径不恶的人、或不接纳回头的人，便是杀害弟兄的凶手，流他的血，如同\Person{加音}流了他兄弟\Person{亚伯尔}的血一样，那血\BibleQuote{从地上向天主喊冤，}\BibleRef{创 4:10}且将受追究。因为义人若被任何人无辜杀害，便永远与天主同享安息。同样，那无缘无故被主教分离的人也如此。那将无辜之人当作害群之马驱逐的，比杀人者更凶残。这样的人不体察天主的仁慈，也不记念祂对悔改者的恩慈，也不注目那些曾是大罪人、因悔改而获赦免的榜样。因此，那抛弃无辜者的人比杀害肉身者更残忍。同样，那不接纳悔改者的人，就是驱散\Person{基督}的羊群，实在是反对祂。因为天主审判罪人固然公义，但接纳回头者时也同样仁慈。因为\Person{达味}，那合天主心意的，在他的诗歌中将仁慈和审判都归于天主。\par

% structure: p0050 source=h3 target=subsubsection reason=relative-heading-rank; base=h2

\subsubsection{\Person{达味}、\Person{尼尼微人}、\Person{希则克雅}和其子\Person{默纳舍}是悔改的杰出榜样：\WorkTitle{犹大王默纳舍的祈祷}}

XXII. 主教啊，你也应当将前人的榜样放在眼前，并巧妙地将它们应用于那些需要严厉或安慰话语的人身上。此外，在你的审判中，遵循天主的旨意是合理的；正如天主对待罪人和那些回头的人一样，你在审判中也应如此行事。天主不是曾通过\Person{纳堂}责备\Person{达味}的过犯吗？然而，当他一说悔改，天主就救他脱离死亡，说：\Quote{放心吧，你必不死。}\BibleRef{撒下 12:13} 同样，当天主使\Person{约纳}被海和鲸鱼吞下，因为他拒绝向\Place{尼尼微}人宣讲，但当他从鲸鱼腹中向天主祈祷时，天主就救他的性命免于败坏。当\Person{希则克雅}一时骄傲，但他一以哀恸祈祷，天主就赦免了他的过犯。但是，主教们啊，请听一个在此场合有益的实例。因为\WorkTitle{列王纪下}和\WorkTitle{编年纪下}这样记载：\Quote{\Person{希则克雅}死后，他的儿子\Person{默纳舍}继位。他登基时十二岁，在\Place{耶路撒冷}作王五十五年；他母亲名叫\Person{赫斐漆巴}。他在主眼前行恶，没有远离主在\Place{以色列}子民面前所消灭的外邦人的可憎之事。\Person{默纳舍}回来建造了他父亲\Person{希则克雅}所拆毁的高丘；他为\Person{巴耳}立了柱像，为\Person{巴耳}筑了祭台，又造了木偶，如同\Place{以色列}王\Person{阿哈布}所行的。他在主的殿中立了祭台，主曾对\Person{达味}和他的儿子\Person{撒罗满}说：我要将我的名放在那里。\Person{默纳舍}设立了祭台，并在那里事奉\Person{巴耳}，说：我的名要永远存留。他在主殿宇的两个院子内为天上的万象筑了祭台；他使自己的儿女在名为\Place{本希农谷}的地方经火；他求问邪术，与交鬼的和行巫术的来往，与行邪术的、观兆的和忒辣芬打交道。他在主眼中大大行恶，惹祂发怒。他铸造了一个偶像，就是他所作的木偶的像，立在主的殿中，主曾在\Place{耶路撒冷}圣城中选择立祂的名永远在那里，并说：我必不再将我的脚从\Place{以色列}地挪开，这地是我赐给他们祖宗的；只要他们谨慎遵行我一切所吩咐的，并我仆人\Person{梅瑟}所吩咐他们的一切诫命。但他们没有听从。\Person{默纳舍}引诱他们行恶，比主从\Place{以色列}子民面前所消灭的列国更甚。主藉祂仆人众先知论到\Person{默纳舍}和祂的百姓说：因为\Place{犹大}王\Person{默纳舍}行了这一切可憎的恶事，比在他以前的\Place{阿摩黎人}更甚，又使\Place{犹大}因他的偶像犯罪，所以主\Place{以色列}的天主如此说：看哪，我必使灾祸临到\Place{耶路撒冷}和\Place{犹大}，凡听见的人，必耳鸣。我必用\Place{撒玛黎雅}的准绳和\Person{阿哈布}家的线铊，拉在\Place{耶路撒冷}上；我必擦净\Place{耶路撒冷}，如同人擦净书板，转过来擦净。我必倾覆它，将我的余民交给他们的仇敌，使他们成为所有仇敌的掠物和战利品，因为他们在我眼中行了这一切恶事，从我将他们祖先领出\Place{埃及}地直到今日，他们惹我发怒。此外，\Person{默纳舍}流了许多无辜的血，直到\Place{耶路撒冷}满溢，从这边到那边，除了他的罪之外，他还使\Place{犹大}犯罪，在主眼前行恶。主使\Place{亚述}王的将帅来攻击他，他们用链子锁住\Person{默纳舍}，用铜镣铐住他，带他到\Place{巴比伦}。他在监狱中被铁链捆绑，用铁链锁住全身。给他吃的是粗糠做的饼，量给的是极少的，喝的是醋调的水，也是量给的，仅够维持生命。他处在痛苦和极大的患难中。当他被强力折磨时，他求告主他的天主的面，在他祖先的天主面前极其谦卑。他向主祈祷说：主，全能的天主，我们祖先\Person{亚巴郎}、\Person{依撒格}、\Person{雅各伯}和他们正义后裔的天主，祢创造了天地并其中的一切装饰，祢以祢命令的话语束缚了大海，祢封闭了深渊，并用祢可畏而荣耀的名封印了它，所有人都畏惧祢，在祢的能力面前战栗。因为祢荣耀的威严不可承受，祢向罪人的愤怒威胁无法忍受。但祢的仁慈应许无法估量、不可测度。因为祢是至高的主，大有怜悯，恒久忍耐，极仁慈，且愿意赦免人恶。主啊，照着祢的厚恩，祢应许了悔改和赦免给那些得罪祢的人，并且因祢无限的仁慈为罪人设立了悔改，使他们可以得救。因此，主啊，祢是义人的天主，并没有为\Person{亚巴郎}、\Person{依撒格}、\Person{雅各伯}这些没有得罪祢的义人设立悔改；但祢为我这个罪人设立了悔改。因为我犯罪比海沙更多。主啊，我的过犯增多了，我的过犯增多了，我不配抬头望天，因为我的不义众多。我被许多铁链压伤，因为我惹动了祢的怒气，在祢面前行了恶，设立了可憎之物，增加了罪行。因此，我屈膝以内心向祢恳求恩宠。主啊，我犯了罪，我犯了罪，我承认我的不义，所以我谦卑恳求祢，饶恕我，主啊，饶恕我，不要因我的不义毁灭我。不要永远向我发怒，为我保留灾祸；也不要定我下到地的深层。因为祢是悔改之人的天主，甚至天主，祢将在我身上显明祢的良善；因为祢要照着祢的大怜悯拯救我这不配的人。因此，我要一生一世永远赞美祢；因为天上万军都赞美祢，荣耀归于祢，直到永远。阿们。主垂听了他的声音，怜悯了他。有一团火焰在他周围显现，所有围绕他的铁锁链都脱落了；主医治了\Person{默纳舍}的患难，带他回到\Place{耶路撒冷}他的国度。\Person{默纳舍}就知道唯有主是天主。他一生尽心尽力单独敬拜主天主；他被称为义人。他除掉主殿中的外邦神祇和偶像，以及他在主殿中建造的所有祭台，和在\Place{耶路撒冷}的所有祭台，把它们都丢出城外。他重修了主的祭台，在其上献了平安祭和感谢祭。\Person{默纳舍}劝勉\Place{犹大}事奉主\Place{以色列}的天主。他与祖先同眠而死，平安无事；他的儿子\Person{阿孟}继位。\Person{阿孟}在主眼中行恶，完全像他父亲\Person{默纳舍}执政初期所行的一切。他惹动了主他的天主发怒。}\par

我们亲爱的孩子们，你们已经听过主天主如何暂时惩罚了那迷恋偶像并杀害许多无辜者的人，然而当他悔改时，主接纳了他，赦免了他的过犯，并将王位归还给他。因为祂不仅宽恕悔改者，还使他们恢复原有的尊荣。\par

% structure: p0053 source=h3 target=subsubsection reason=relative-heading-rank; base=h2

\subsubsection{阿孟可作为故意犯罪者的鉴戒}

二十三、没有比崇拜偶像更严重的罪，因为这是对天主的亵渎。然而，即使是这样的罪，在真诚悔改后也得到了宽恕。但若有人故意犯罪，存心试探天主是否会惩罚恶人，这样的人将得不到赦免，即使他对自己说：\Quote{一切都好，我要随从我邪恶的意念行事。} 这样的人就是阿孟，默纳舍的儿子。因为圣经说：\BibleQuote{阿孟起了邪恶的叛逆念头，说道：我的父亲从小就犯了极大的叛逆罪，但在老年时悔改了；现在我要随从我的灵魂所欲，之后再归向天主。他在天主眼中行了恶，超过他以前所有的人。天主迅速将他从美好的土地上彻底除灭。他的臣仆合谋，在他家中杀了他，他只统治了两年。}\par

% structure: p0055 source=h3 target=subsubsection reason=relative-heading-rank; base=h2

\subsubsection{基督耶稣我们的主来是为了拯救悔改的罪人}

二十四、因此，你们这些平信徒，要小心，免得你们中有人心中存有阿孟的意念，而突然被剪除而灭亡。同样，主教当尽一切努力，使那些尚未犯罪的人不致陷入罪恶，并医治和接纳那些从罪恶中回转的人。但如果他冷酷无情，不肯接纳悔改的罪人，他就是得罪了主他的天主，自以为比天主的公义更公义，不肯接纳天主因基督的缘故所接纳的人。为此，天主差遣祂的儿子降世为人，为了人类；为此，天主喜悦那造男造女的主，由女人所生；为此，天主没有怜惜祂，没有免于十字架、死亡和埋葬，而是容许那位本性不能受难的、祂的爱子、天主圣言、祂伟大议会的天使死去，为要拯救那些当死的罪人。那些不接纳悔改者的人，正是激怒了祂。因为祂不以我——玛窦——为耻，我从前是税吏；祂接纳了伯多禄，当他因害怕三次否认主，但以悔改和痛哭使主息怒后，主立他为祂羊群的牧者。此外，祂立了我们的同宗徒保禄，使他从迫害者变为宗徒，并宣称他是蒙拣选的器皿，尽管他先前曾加给我们许多祸害，并亵渎了祂的圣名。祂也对另一名犯罪的妇女说：\BibleQuote{你的罪，虽多，已得赦免，因为你爱得多。} \BibleRef{路 7:47} 当长老们将另一名犯罪的妇女带到祂面前，将判决留给祂并离开后，我们的主，那洞察人心的主，问她是否长老们定了她的罪，当她回答没有时，祂对她说：\BibleQuote{你去吧，我也不定你的罪。} \BibleRef{若 8:11} 啊，主教们，这位耶稣——我们的救主、我们的君王、我们的天主——应当作为你们的榜样摆在面前；你们应当效法祂，温和、安静、慈悲、怜悯、和平、无欲、善于教导、勤于劝化、乐于接纳和安慰；不打人、不轻易发怒、不伤害人、不傲慢、不骄傲、不嗜酒、不醉酒、不虚华浪费、不贪美味、不挥霍，使用天主的恩赐不当作别人的，而当作自己的，如同被委派的忠心管家，将来要向天主交账。\par

% structure: p0057 source=h2 target=subsection reason=relative-heading-rank; base=h2

\subsection{第四节 关于所收集资源的处理：用以供养圣职人员并救济穷人}

主教应认为那些符合需要和体面的食物和衣服已经足够。他不可将主的东西当作别人的而滥用，当有节制，\BibleQuote{因为工人配得他的报酬} \BibleRef{路 10:7}。他不可在饮食上奢侈，或喜爱无用的家具，只当满足于维持生计所需。\par

% structure: p0059 source=h3 target=subsubsection reason=relative-heading-rank; base=h2

\subsubsection{关于初熟之物和什一之物，以及主教应如何享用它们或分施给别人}

XXV. 他应按照天主的命令，使用所赐的什一与初果，作为天主之人；同样，他应公正地分配那些因穷人、孤儿、寡妇、困苦者和患难中的外乡人而带来的自愿奉献，因为他将那位把管理权交托给他的天主视为他账目的审查者。要公义地分给所有匮乏者，你们自己可以使用属于主的事物，但不可滥用；可以吃用，但不可独自吃尽：要与匮乏者分享，从而在天主面前显为无可指摘。因为如果你们独自享用，你们必受天主的斥责，祂对这些贪得无厌、独吞一切的人说：\BibleQuote{你们吃奶，穿羊毛}；\BibleRef{则 34:3} 又在另一处说：\BibleQuote{难道你们要独自住在地上吗？} \BibleRef{依 5:8} 因此，律法中命令你们说：\BibleQuote{你要爱你的邻人如你自己。} \BibleRef{肋 19:18} 我们现在说这些，并非认为你们不可享用自己劳苦的果实；因为经上记着：\BibleQuote{牛在场上踹谷时，不可笼住它的嘴。} \BibleRef{申 25:4} 而是说你们应当以节制和公义行事。所以，正如牛在打谷场上劳作而不被笼嘴，确实吃，但并非吃尽；同样，你们在打谷场上劳作者，即在天主的教会中劳作者，也当从教会得食。这也是肋未人的情形，他们曾在会幕（证言）中服务，而会幕在各方面都是教会的预像。不仅如此，它的名字本身就暗示那会幕预先被指定作为教会的证言。因此，侍奉会幕的肋未人也得以分享所有人民献给天主的物品——即礼物、供物、初果、什一、祭品和献仪，不受干扰，他们连同妻子、儿女。因为他们的职分是服侍会幕，所以在以色列子民中在土地上没有分得产业，因为人民的献仪就是肋未的份，是他们支派的产业。因此，诸位主教，你们对于你们的人民而言，就是司铎和肋未人，服侍圣所，即圣而公天主教会；你们站在上主你们的天主的祭台前，并借大司祭\Person{耶稣基督}向祂献上合理无血的祭献。你们对于平信徒而言，是先知、首领、治理者、君王；是天主与其忠信子民之间的中保，领受并宣告祂的圣言，熟谙圣经。你们是天主的声音，祂旨意的见证人，担负众人的罪，并为众人代祷；你们已经听到，天主的话严厉警告你们，若你们向人隐藏知识的钥匙，若不向你们辖下的子民宣告祂的旨意，他们就要丧亡；若你们妥善服侍圣所，你们将从天主得着确定的赏报，以及不可言喻的尊荣与荣耀。因为正如你们担负重担，你们也当得着食物和其他必需品的供应作为果实。因为你们效法主基督；正如祂在受难时，\Quote{在木头上亲自担负了我们众人的罪}，无辜者为应受惩罚者，同样，你们也应当将人民的罪当作自己的罪。因为关于我们的救主，\Person{依撒意亚}说：\BibleQuote{祂承担了我们的罪，为我们受苦。} \BibleRef{依 53:4} 又说：\BibleQuote{祂承担了许多人的罪，并为我们的过犯被交付。} \BibleRef{依 53:12} 因此，正如你们是别人的榜样，你们也有基督作你们的榜样。因此，正如祂关怀所有人，你们也当关怀你们辖下的平信徒。难道你们以为主教的职务是容易或轻省的重担吗？因此，正如你们承担重量，你们也有权利在众人之前分享果实，并分给那些匮乏的人，因为你们要向那毫不偏袒地审查你们账目的祂交账。因为侍奉教会的人应当由教会供给，作为司铎、肋未人、首领和天主的仆人；正如\BibleBook{户籍}{纪}中论司铎所记：\BibleQuote{上主对\Person{亚郎}说：你和你儿子，以及你家族的人，应担当司铎圣物方面的罪过。} \BibleRef{户 18:1} \BibleQuote{看，我把以色列子民奉献给我的初果交给你管理；我把这些赐给你和你的子孙，作为永久的酬劳。这些要归你们，从至圣之物中，从祭献中和礼物中，从一切祭品和赎过祭、赎罪祭中；凡他们因一切圣物所献给我的，都归你和你的儿子。你们要在圣所中吃这些。} 稍后又说：\BibleQuote{凡最好的油、酒、麦子，他们献给上主的初熟之物，我都赐给你；凡初熟之物，我都赐给你，凡献与我的圣物，也都赐给你。凡在以色列人中所献的人身和牲畜的首胎，洁净和不洁净的，以及祭品中的胸和右腿，都归司铎和他们的家属，以及肋未人。}\par

你们这些平信徒，天主的特选教会，请听。因为从前百姓被称为\Quote{天主的子民} \BibleRef{出 19:5}和\Quote{圣洁的国民} \BibleRef{希 12:23}。因此，你们就是圣洁而神圣的\Quote{天主的教会}，被登记在天上，\Quote{君尊的司祭，圣洁的国民，特选的民族}，为天主主装饰的新娘，伟大的教会，忠信的教会。现在请留心听从前所说的：祭品和什一之物属于基督——我们的大司祭，以及那些为祂服务的人。救恩的什一就是耶稣名字的第一个字母。请听，圣而公教会，你逃脱了十灾，领受了十诫，学习了法律，持守了信德，相信了耶稣，认识了十的数字，相信了耶稣名字的第一个字母约塔，并以祂的名字命名，被建立，并在祂光荣的圆满中发光。从前是祭献的，如今是祈祷、转求和感恩。从前是初果、什一、供物和礼物，如今是祭品，由圣主教们通过耶稣基督——为他们而死的那位——献给天主。因为他们是你们的大司祭，如同司铎是你们的司铎，你们现在的执事代替你们的肋未人；同样，你们的读经员、歌咏员、守门者、女执事、寡妇、贞女和孤儿也是如此；但超越这一切的那位是大司祭。\par

% structure: p0062 source=h3 target=subsubsection reason=relative-heading-rank; base=h2

\subsubsection{论各级圣职人员按何种模式和尊荣由天主委任。}

主教是圣言的施行者，知识的保管者，在你们神圣敬拜的各个部分中，你们与天主之间的中保。他是虔敬的导师；并且在天主之后，他是你们的父亲，曾藉水和圣神重生你们，使你们成为嗣子。他是你们的统治者和治理者；他是你们的君王和掌权者；他在天主之后，是你们地上的神，理应受你们的尊敬。因为关于他和像他这样的人，天主宣告说：\BibleQuote{我曾说：你们是神，都是至高者的子女。}并且\BibleQuote{不可毁谤神。}\BibleRef{出 22:28} 因为让主教管理你们，他受有天主的权柄，这权柄要执行于圣职人员之上，并治理全体百姓。让执事服侍他，如同基督服侍祂的父；让他凡事无可指责地服侍他，如同基督不自作主张，只做那讨祂父喜悦的事。也让女执事受你们尊敬，如同圣神之位，她不可在没有执事的情况下行事或说话；同样，护慰者也凭自己不说什么或做什么，而是等候祂的旨意，将光荣归于基督。正如我们若没有圣神的教导就不能相信基督，同样，任何女子没有女执事也不可接近执事或主教。让司铎被你们视为代表我们宗徒，让他们做神圣知识的导师；因为我们的主派遣我们时曾说：\BibleQuote{你们要去使万民成为门徒，因父、及子、及圣神之名给他们授洗，教训他们遵守我所吩咐你们的一切。}\BibleRef{玛 28:19} 让寡妇和孤儿被视为代表全燔祭台；让贞女受尊敬，代表香祭台和香本身。\par

% structure: p0064 source=h3 target=subsubsection reason=relative-heading-rank; base=h2

\subsubsection{论人擅自闯入任何司铎职务是可怕的事，如科辣黑及其同党、撒乌耳和乌齐雅所作。}

因此，正如从前不属于肋未支派的外族人，不可奉献任何东西，也不可在没有司铎的情况下接近祭台，同样，你们也不可没有主教做任何事；因为若有人没有主教而行某事，他所行的徒然无效。因为撒乌耳在没有撒慕尔的情况下献祭时，曾被告知：\BibleQuote{这为你没有益处。}\BibleRef{撒上 13:13} 同样，任何平信徒若没有司铎而行某事，也是徒劳。又正如乌齐雅王，他不是司铎，却妄行司铎的职务，因此因他的悖逆而患了癞病；同样，任何轻慢天主、狂妄冒犯祂的司铎、不义地夺取那尊荣的平信徒，必不免受罚：他不效法基督，\BibleQuote{基督并非自取光荣成为大司祭，}\BibleRef{希 5:5} 而是等待听见父说：\BibleQuote{主起了誓，绝不反悔：你永为司祭，按照默基瑟德的品位。} 如果基督没有父的认可就不自取光荣，那么人若没有从上级接受这尊荣，却擅自闯入司铎职，去做只有司铎才可做的事，他岂敢呢？跟随科辣黑的人，虽然他们是肋未支派，不也是因为起来反对梅瑟和亚郎，干预不属于他们的事而被火焚烧了吗？达堂和阿彼兰活活下到地狱；那发芽的棍杖制止了群众的轻举妄动，并显示出谁是天主所傅的大司祭。所以，弟兄们，你们应当把你们的祭品和奉献带到主教那里，如同带到你们的大司祭那里，或亲自带去，或由执事代送；不仅带来这些，还要把你们的初果、什一和自愿奉献都带给他。因为他知道谁是处于困苦中的人，他按适宜给予每个人，这样就没有人在同一天或同一周内两次或多次接受施舍，而另一人却一无所有。因为比较合理的做法是供应那些真正处于困境者的需要，而不是那些表面如此的人。\par

% structure: p0066 source=h3 target=subsubsection reason=relative-heading-rank; base=h2

\subsubsection{论宴席，以及邀请圣职人员的各等级者应如何对待他们。}

二十八．如果有人决定邀请年长妇女参加爱宴或筵席——如同我们的救主所称呼的，\BibleRef{路 14:13}——让他们最常派人去让执事知道有困难的那位妇女。但在筵席中，应当为牧者保留他所应得之物，即初果，即使他本人不在席间，因为他是你们的司铎，且为尊敬那位将司铎职托付于他的天主。至于给每位年长妇女的，应给执事双倍，以尊敬基督。也应为长老预留双份，因为他们是不断劳苦于圣言和教训的人，因我们主的宗徒的缘故，他们代表宗徒的地位，作为主教的谋士和教会的冠冕。因为他们就是教会的公议会和元老院。若有读经者在场，他应得一份，以尊敬先知；歌咏者和守门者也应得相同。因此平信徒应以他们的礼物向每个不同等级致以适当的尊敬和最大的敬意。但他们不应事事打扰他们的主教，而应通过那些服侍他的人，即执事，来传达他们的愿望，因为他们可以更自由地与他们交往。因为我们也不能直接向全能的天主说话，而只能通过基督。同样地，让平信徒通过执事向主教表明他们所有的愿望，并相应地按他的指示行事。因为在古时的圣殿中，没有司铎就不能献上或举行任何圣物。先知在某处说：\BibleQuote{司铎的嘴唇应保持知识，人们应从他的口中寻求法律，因为他是全能之主的使者。}\BibleRef{拉 2:7}因为若连崇拜魔鬼的人，在他们可憎、可厌和不洁的仪式中，直到今日仍模仿神圣的规矩（这真是个宽泛的比较，他们的可憎之事与天主的圣敬之间差距甚大），他们在对崇拜的戏仿中，若不通过他们所谓的司铎，就既不献祭也不做任何事，反而视他为他们的石制偶像之口，等待看他将给他们什么命令。凡他命令他们的，他们就做，没有他他们什么都不做；他们尊敬他，他们所谓的司铎，视他的名为可敬，以敬拜无生命的雕像，并为了崇拜邪恶的鬼神。那么，如果这些外邦人，他们荣耀虚妄的假神，并将希望寄托于虚无，尚且努力模仿神圣的规矩，那么你们这些拥有最确定的信德和毫无疑问的望德，并期待荣耀、永恒、永不落空的应许的人，难道不该更合理地尊敬上主天主于那些管理你们的人身上，并视你们的主教为天主之口吗？\par

% structure: p0068 source=h3 target=subsubsection reason=relative-heading-rank; base=h2

\subsubsection{主教与执事之尊荣}

二十九．因为如果\Person{亚郎}因从\Person{梅瑟}向\Person{法郎}宣告天主的话而被称为先知；\Person{梅瑟}自己因同时为君王和大司祭也被称为\Person{法郎}的神，正如天主对他说：\BibleQuote{我使你成为\Person{法郎}的神，你的兄弟\Person{亚郎}要作你的先知。}\BibleRef{出 7:1}为什么你们不也视圣言的中间人为先知，并尊敬他们如同神呢？\par

% structure: p0070 source=h3 target=subsubsection reason=relative-heading-rank; base=h2

\subsubsection{平信徒应以何种方式服从执事}

三十．因为现时执事是你们的\Person{亚郎}，主教是\Person{梅瑟}。因此，如果\Person{梅瑟}被主称为神，那么你们中应尊敬主教如同神，执事如同他的先知。因为\Person{基督}不做什么没有祂的圣父，执事也不做什么没有他的主教；正如圣子没有圣父就一无所是，执事没有主教也一无所是；正如圣子服从圣父，同样每位执事服从他的主教；正如圣子是圣父的使者与先知，同样执事是主教的使者与先知。因此，凡他需与任何人进行的一切事，都应当禀明主教，并由他最后裁定。\par

% structure: p0072 source=h3 target=subsubsection reason=relative-heading-rank; base=h2

\subsubsection{执事不得在没有主教的情况下做任何事}

三十一．他不可在没有他的主教的情况下做任何事，也不可在没有他同意的情况下施予任何东西。因为如果他在主教不知情的情况下给某人东西，如同给一个困难的人，他这样做必然导致对主教的责备，并指责他漠不关心困难的人。但任何用言语或行为责备他主教的人，就是反对天主，不听祂对他们说的话：\BibleQuote{你不要毁谤神。}\BibleRef{出 22:28}因为祂制定那法律，不是针对木石的神祇——那些是可憎的，因为它们被假称为神——而是针对司铎和审判官，祂也曾对他们说：\BibleQuote{你们是神，都是至高者的子女。}\par

% structure: p0074 source=h3 target=subsubsection reason=relative-heading-rank; base=h2

\subsubsection{执事不得在没有主教同意的情况下施予任何东西，因为那将导致对主教的责备}

XXXII. 所以，执事啊，如果你知道有人处于困境，就提醒主教，然后施予他；但不可暗中行事，以免招致对他的指责，免得激起对他的怨言；因为怨言不是针对他，而是针对主天主；执事和其余的人将听到亚郎和米黎盎在反对梅瑟时所听到的话：\Quote{你们怎么敢反对我的仆人梅瑟呢？} \BibleRef{户 12:8} 此外，梅瑟对起来反对他的人说：\Quote{你们的怨言不是针对我们，而是针对主我们的天主。} \BibleRef{出 16:8} 因为如果有人称一个平信徒为\OriginalTerm{拉卡}{Raka}，\BibleRef{玛 5:22} 或愚人，尚且不能无罪，因为这是伤害基督的名号，那么，任何人怎敢反对自己的主教呢？主藉着主教的覆手将圣神赐予你们中间，藉着主教你们学习了神圣的教义，认识了天主，信了基督，藉着主教你们被天主所认识，藉着主教你们被以喜乐之油和明智之膏所傅抹，藉着主教你们被宣称为光明之子，藉着主教主在你们的洗礼中藉主教的覆手作证，并对你们每一个人发出神圣的声音说：\Quote{你是我的儿子，我今天生了你？} 人啊，通过你的主教，天主收养你为祂的孩子。儿子啊，你要承认那只曾是母亲的手。要爱他，他在天主之后成了你的父亲，要尊敬他。\par

% structure: p0076 source=h3 target=subsubsection reason=relative-heading-rank; base=h2

\subsubsection{应当如何尊敬主教，以及视他们为属灵父母而予以敬重。}

XXXIII. 因为如果神的圣言论到我们属肉身的父母说：\Quote{应孝敬你的父亲和你的母亲，好使你获得幸福；} \BibleRef{出 20:12} 又说：\Quote{凡辱骂父亲或母亲的，应处死刑；} \BibleRef{出 21:17} 那么，圣言岂不更应劝勉你们尊敬属灵的父母，爱他们如同你们的恩人和与天主同在的使者？他们藉水重生了你们，以圣神的圆满装备了你们，以言语如同乳汁喂养了你们，以教义滋养了你们，以训诫坚定了你们，将基督的救恩之体和宝血分施给你们，解除了你们的罪，使你们分沾了神圣而纯洁的圣体，接纳你们为天主恩许的分享者和同继承人！要敬畏他们，以各种尊敬尊重他们；因为他们从天主那里获得了生与死的权能：审判罪人，判定他们受永火之死，以及赦免悔改者的罪，恢复他们到新生命。\par

% structure: p0078 source=h3 target=subsubsection reason=relative-heading-rank; base=h2

\subsubsection{司铎应优先于统治者和君王。}

XXXIV. 你们要视这些人配得作为你们的统治者和君王，向他们进贡如同向君王进贡；因为你们应当供养他们和他们的家庭。正如撒慕尔在《列王纪上》为人民制定了关于君王的条例，梅瑟在《肋未纪》中为司祭制定了条例，同样我们也为你们制定了关于主教的条例。因为如果在那里百姓按比例将次要的服务分配给这样一位伟大的君王，那么，主教现在不是更应当从你们那里接受天主为他自己及其所属的其他圣职人员的生活所规定的东西吗？但如果我们还可以进一步说，主教应比古代其他人接受更多：因为旧时的君王只管理军队的事务，受托处理战争与和平以保全人的身体；而这位主教受托在天主面前行使司铎职，以保全身体和灵魂免于危险。因此，灵魂比身体越珍贵，司铎职就越超越王权。因为司铎职捆绑和释放那些应受惩罚或宽恕的人。所以，你们应当爱主教如同父亲，敬畏他如同君王，尊敬他如同主人，将你们的果实和双手的工作带给主教，作为对你们的祝福，将你们的初熟之果、十一之物、奉献和礼物交给他，如同给天主的司铎：你们的小麦、酒、油、秋季果实、羊毛以及主天主赐给你们的一切的初熟之果。你们的奉献将被悦纳，作为上主天主前馨香的祭品；主必降福你们双手的工作，并使土地的美物增多。\Quote{施予者必蒙祝福。} \BibleRef{箴 11:26}\par

% structure: p0080 source=h3 target=subsubsection reason=relative-heading-rank; base=h2

\subsubsection{法律和福音都规定奉献。}

三十五. 现在你们应当知道，虽然主已经将你们从附加的束缚中释放出来，将你们带出它们，使你们获得舒畅，并且祂不许你们将无理性的受造物献为赎罪祭、洁净祭、替罪羊、以及不断的洗涤和洒水，但祂绝没有豁免你们欠给司铎的那些祭献，也没有豁免你们对穷人行善。因为主在福音中对你们说：\BibleQuote{除非你们的正义超过经师和法利塞人的正义，你们决进不了天国。}\BibleRef{玛 5:20} 在这方面，你们的正义将超过他们的，如果你们更关心司铎、孤儿和寡妇；正如经上记载：\Quote{他施舍，他给了穷人；他的正义永远常存。}又说：\BibleQuote{借着正义的行为和信德，罪过得以涤除。}\BibleRef{箴 16:6}又说：\BibleQuote{慷慨的灵魂是有福的。}\BibleRef{箴 11:25}因此，你们要照主所规定的去做，将归于司铎的东西给他，即你们禾场的初果、酒榨的初果、以及赎罪祭，就如同给那在需要涤除和赦免的人与天主之间的中保。因为你们的责任是给予，他的责任是管理，作为教会事务的管理者和处置者。然而，你们不得追究你们的主教，也不得监视他的管理——他如何做、何时、给谁、何处、做得好或坏或一般；因为他有位会向他追究的那一位，即主天主，是祂将这项管理交在他手中，并认为他配得上如此大尊荣的司铎职。\par

% structure: p0082 source=h3 target=subsubsection reason=relative-heading-rank; base=h2

\subsubsection{十诫的叙述，以及它们在此如何规范我们。}

三十六. 你们要把天主的敬畏摆在眼前，常记天主的十诫——全心全意爱独一的主天主；不可留心于偶像或任何其他存在，以为它们是死的神、无理性的存在或魔鬼。要思量天主那繁复的造化，那通过基督而开始的。你们要遵守安息日，因为那位停止创世工作却不停止眷顾工作的那位；这是为默思法律的休息，不是为双手的闲懒。摒弃一切非法的欲望、一切毁灭人的事物、以及一切愤怒。孝敬父母，即为你们生命的源头。爱近人如自己。将生活必需品分给有需要的人。避免发假誓、常发誓以及虚誓；因为你们不会算为无罪。不可空手出现在司铎面前，要不断奉献你们的自愿祭献。此外，不要离开基督的教会；要在一切工作之前于早晨去那里，并在晚上再次聚集，感谢天主保存了你们的生命。在你们的召叫中要勤勉、恒常和劳苦。向主献上你们的自愿祭献；因为祂说：\BibleQuote{你要以诚实的劳动果实尊崇主。}\BibleRef{箴 3:9} 如果你们不能向\OriginalTerm{科爾班}{Corban}投下什么像样的东西，至少施舍给外方人一、二、五个小钱。\BibleQuote{你要为自己积攒天上的宝藏，那里没有虫蛀，也没有贼挖洞偷窃。}\BibleRef{玛 6:20} 在做这事时，不要判断你们的主教或任何平信徒邻人；因为如果你判断你的弟兄，你就成了一个判断者，而没有受任何人的任命，因为判断权只托付给了司铎。因为对他们说：\BibleQuote{你们要施行公义的判断；}\BibleRef{匝 7:9} 又说：\Quote{你们要证显自己是精确的兑换银钱者。} 因为这并未托付给你们；相反，对那些没有长官或职员尊位的人说：\Quote{你们不要判断，你们就不会受判断。}\par

% structure: p0084 source=h2 target=subsection reason=relative-heading-rank; base=h2

\subsection{第五节　论控告，以及对控告者的处理}

% structure: p0085 source=h3 target=subsubsection reason=relative-heading-rank; base=h2

\subsubsection{论控告者与假控告者，以及审判者不宜轻信或轻疑，而应详察。}

三十七. 但主教的职责是正确判断，正如经上记载：\BibleQuote{你们要施行公义的判断；}\BibleRef{若 7:24} 又在别处说：\BibleQuote{你们为什么连自己也不判断什么是正义的呢？}\BibleRef{路 12:57} 因此，你们要像熟练的兑换银钱者：因为正如他们拒绝劣币，收兑良币，同样，主教的职责是保留无可指责的人，而医治或——如果到了不可救药的地步——抛弃那些有可责之处的人，以免急于割除，也不轻信一切控告；因为有时有人出于激情或嫉妒，坚持对弟兄作假控告，如同在巴比伦的\Person{苏撒纳}案中两个长老所做的，以及埃及妇人在\Person{若瑟}案中所做的。因此，作为天主的人，不要轻率接受这样的控告，免得你夺去无辜者的性命，杀害义人；因为那接受这种控告的人是愤怒而非和平的制造者。哪里有愤怒，那里就没有主；因为那愤怒——即撒殚的朋友——借着假弟兄不公义地激动起来的愤怒，绝不容许教会内有同心合意。所以，当你认识这样的人是愚妄的、好争吵的、易怒的、喜欢惹是生非的，不要相信他们；而要观察他们的真实面目，当你从他们口中听到任何反对弟兄的话：因为谋杀在他们眼中不算什么，他们陷害人的方式出乎意料。因此，要仔细考察控告者，智慧地观察他的生活方式，他是何许人、何等样人；如果你发现他是个诚实的人，就按照我们的主的教导\BibleRef{玛 18:15}，带那被控告的人，在没有旁人的时候责备他，使他悔改。但如果他不听劝，你就带一两个人同去，好指正他的过错，用温和与教诲劝诫他；因为\BibleQuote{智慧将安息在善人心中，但在愚人心内却不被认识。}\BibleRef{箴 14:32}\par

% structure: p0087 source=h3 target=subsubsection reason=relative-heading-rank; base=h2

\subsubsection{论罪人应私下受责，而悔改者应被接纳，照我们主的规范。}

三八、因此，若他因你们三人的口而被说服，那便好了。但若有人硬着心，\BibleQuote{你就告诉教会；如果他连教会也不听从，就把他当作外邦人和税吏对待吧；} \BibleRef{玛 18:17} 就不要再接他进入教会如同基督徒一样，而要拒绝他如同外邦人。但若他愿意悔改，就接他回来。因为教会并不接纳外邦人或税吏共融，除非他们对以前的恶行一一悔改；因为我们的主耶稣，天主的基督，已为按悔改接纳人设立了地方。\par

% structure: p0089 source=h3 target=subsubsection reason=relative-heading-rank; base=h2

\subsubsection{悔改的典范}

三九、因为我\Person{玛窦}，就是在这教义中向你们讲话的十二人之一，是位宗徒；我自己从前曾是个税吏，但因着信德而现在蒙受了怜悯，对从前行为悔改了，并获得了宗徒和圣言宣道者的殊荣。\Person{匝凯}，主曾因他的悔改和向他祈祷而接纳了他，他也同样起先是个税吏。此外，那些来听主悔改之道的士兵和税吏群，也听见先知\Person{若翰}在为他们付洗后说了这话：\BibleQuote{除了给你们指定的以外，不要再多做什么。} \BibleRef{路 3:13} 同样，生命也不拒绝外邦人，若他们悔改并抛弃不信。因此，凡被证实犯有任何恶行而不悔改的，就视他如同税吏或外邦人。但若他后来悔改，从错误中转回，那么，正如我们接纳外邦人当他们愿意悔改时，确实接他们进入教会听道，但不让他们领受共融直到他们领受了圣洗的印记并成为完美的基督徒；同样我们也只让这样的人进入听道，直到他们显出悔改的果实，好叫他们因听道不至完全而不可挽回地丧亡。但不要让他们参与祈祷的共融；在诵读法律、先知和福音之后，就让他们离开，好叫他们因这样的离去而改善自己的生活，努力每日参与公共集会并恒常祈祷，以便他们最终也能被接纳，并使看见他们的人受到触动，更加谨慎，免得陷入同样的状况。\par

% structure: p0091 source=h3 target=subsubsection reason=relative-heading-rank; base=h2

\subsubsection{对犯过一次或两次罪的人，我们不应不留情}

四十、但主教，你切莫立刻厌恶那陷入一两次过失的人，也不可将他们逐出主的圣言，也不可拒绝他们与常人交往，因为主也不拒绝和税吏与罪人一同吃饭；当法利塞人为此指责他时，他说道：\BibleQuote{健康的人不需要医生，而是有病的人才需要。} \BibleRef{玛 9:12} 所以，你要与那些因罪而与你们分离的人一同生活和相处；照顾他们，安慰他们，坚固他们，对他们说：\BibleQuote{你们软弱的手，无力的膝，应坚强起来。} \BibleRef{依 35:3} 因为我们应当安慰那些忧伤的人，激励灰心的人，免得因过度的忧伤而陷入混乱，因为\Quote{灰心的人极其慌乱。}\par

% structure: p0093 source=h3 target=subsubsection reason=relative-heading-rank; base=h2

\subsubsection{我们应以何种方式接纳悔改者；应如何对待犯罪者，以及何时应将他们从教会中剪除}

四一。但若有人回头，结出悔改的果实，你就要接纳他祈祷，如同那个失落的小儿子，那个荡子，他与娼妓耗尽了父亲的财产，放猪，想吃豆荚而不能得。这儿子悔改后回到父亲那里，说：\Quote{我得罪了天，也得罪了你，我不配再称为你的儿子；}\BibleRef{路 15:21}父亲满心慈爱，用音乐迎接他，给他穿上旧袍、戴上戒指、穿上鞋，宰了肥牛犊，与朋友们一同欢乐。所以，主教啊，你也该照样做。正如你教导并施洗外邦人后接纳他，同样，你要让众人为这人祈祷，借着覆手将他恢复到羊群中原有的位置，如同被悔改洁净的人；那覆手对他而言代替圣洗：因为我们覆手时，圣神赐给了信徒。若有一位始终屹立不摇的弟兄因你与他和好而责备你，你要对他说：\Quote{你常同我在一起，凡我所有的都是你的。应当欢乐庆祝，因为你这兄弟死而复生，失而复得。} 天主不仅接纳忏悔者，还恢复他们原有的尊严，圣\Person{达味}是充分的见证。他犯了\Person{乌黎雅}的事之后，向天主祈祷说：\Quote{求祢还原我救恩的喜乐，以慷慨的圣神扶持我。} 又说：\Quote{求祢转面不看我的罪过，涂去我的一切过犯。天主，求祢在我内再造一颗纯洁的心，在我五内注入坚强的精神。不要将我从祢面前抛弃，不要从我身上收回祢的圣神。} 因此，要像慈悲的医生，医治所有犯罪的人，使用使人得救的疗法；不仅切割、烧灼或使用腐蚀剂，也要包扎、放入药布、使用温和的愈合药，并洒上安慰的话语。若是空洞的伤口或巨大的裂口，用合适的膏药滋养，使之填满，与其余完好肌肉齐平。若有污秽，用腐蚀性的粉末——即责备的话语——洁净。若有赘肉，用尖锐的膏药——审判的威胁——消除。若蔓延开来，就烧灼，切除腐肉，以禁食使其死灭。但若在你做了这一切之后，看出从头到脚没有热敷、油或绷带的余地，疾病肆虐，妨碍一切治疗，如同坏疽腐蚀整个肢体；那么，经深思熟虑，并与其他高明医生商议，切除腐化的肢体，免得整个教会身体被腐蚀。因此，不要轻易仓促切除，也不要轻易求助于多齿的锯；而要先使用柳叶刀切开伤口，抽出引起疼痛的内在原因，使身体免于疼痛。但若你看见有人无可救药，已麻木不仁，就含悲忍痛地将这不可救药的人从教会中割除。因为：\Quote{从你们中间除去那恶人。}\BibleRef{申 17:7} 又说：\Quote{要叫以色列子民畏惧。}\BibleRef{肋 15:31} 又说：\Quote{审判时，不可看顾富人的情面。}\BibleRef{申 1:17} 又说：\Quote{不可在穷人争讼的事上偏袒穷人，因为审判属于上主。}\BibleRef{出 23:3}\par

% structure: p0095 source=h3 target=subsubsection reason=relative-heading-rank; base=h2

\subsubsection{审判者不可徇私。}

四二。但若诬告是假的，你们这些牧者与执事，由于徇私或受贿，情愿取悦魔鬼，因而将无辜的被告逐出教会，你们将在主的日子交账。因为经上记载：\Quote{不可杀害无辜和义人。}\BibleRef{出 23:7} \Quote{不可拿礼物打灵魂，因为礼物蒙蔽智者的眼目，败坏义人的言语。} 又说：\Quote{他们因接受礼物而释放恶人，夺去义人的义。}\BibleRef{依 5:23} 因此，要小心，不可不公正地定任何人的罪，从而助长恶人。因为\Quote{祸哉，那些称恶为善，称善为恶，以苦为甜，以甜为苦，以光为暗，以暗为光的人。}\BibleRef{依 5:20} 所以，要谨慎，免得你们成了徇私之人，从而落在这主的声音之下。因为若你们不公正地定别人的罪，便是定自己的罪。主说：\Quote{你们用什么判断来判断，也要受什么判断；你们怎样定人的罪，也要怎样被定罪。}\BibleRef{玛 7:2} 因此，若你们不徇私判断，就会揭穿那作假见证陷害邻舍的控告者，证明他是诬告者、恶毒的人、杀人者，他指控那人如同恶人，言语不定，自相矛盾，被自己的话缠住；因为他的嘴唇是他危险的罗网。你一旦证明他说谎，就要严厉审判他，把他交给火剑，照他恶意对待兄弟的做法对待他，因为他尽其所能杀了他的兄弟，抢先占据了法官的耳朵。\BibleRef{申 19:19} 经上又记载：\Quote{凡流人血的，他的血也必被人流。}\BibleRef{创 9:6} 又说：\Quote{你们要从自己中间除去那无辜的血，就是无故流出的血。}\BibleRef{申 19:13}\par

% structure: p0097 source=h3 target=subsubsection reason=relative-heading-rank; base=h2

\subsubsection{假控告者应如何受罚。}

XLIII. 因此，你该把他当作杀害自己弟兄的凶手逐出会众。过些时候，若他说自己悔改了，你当以斋戒克苦他，然后给他覆手，接纳他，但仍要防备他，免得他再次扰乱任何人。但若他再次被接纳后仍然如此滋事，不肯停止与弟兄争吵争斗，出于争辩之心挑剔过失，就把他当作有害的人逐出，免得他摧残天主的教会。因为这样的人是城中的骚乱者；他虽身在教会内，却不成其为教会，只是一个多余无用的肢体，尽其所能给基督的身体抹黑。因为如果有人身上长有赘余的肢体——如多出的手指或肉瘤——他们因着不雅而将其割除，于是丑陋消失，那人在外科医生的帮助下恢复天然的美形；那么，你们这些教会的牧者（因为教会是一个完美的身体，是那些敬畏主、在爱中信从天主的人的健全肢体）岂不更应当如此行，若其中发现一个心怀恶意的多余肢体，使身体其余部分变得不雅，并以叛乱、争斗、诽谤来扰乱它，引起恐惧、骚乱、污点、诽谤、控告、混乱，并做魔鬼类似的作为，如同被魔鬼指派以诽谤、大乱、争斗、分裂来羞辱教会！这样的人，若第二次被逐出教会，就理当完全从主的会众中被剪除。如今主的教会将比先前更美，因为它曾有一个多余而令它不悦的肢体。因此，从此它免受责备和羞辱，成为纯净，没有这些邪恶、诡诈、辱骂、不仁慈、背信的人；没有这些\Quote{憎恨良善者、喜爱宴乐者}\BibleRef{弟后 3:3}，贪图虚荣者，欺骗者，假装有智慧者；没有这些以驱散、甚至彻底离散主的羊群为业的人。\par

% structure: p0099 source=h2 target=subsection reason=relative-heading-rank; base=h2

\subsection{第六节 信友的争议应由主教裁决，并使信友和好}

因此，主教啊，连同你属下的圣职人员，要努力正确地剖分真理之言。因为主说：\Quote{如果你们与我作对，我也要与你们作对。}\BibleRef{肋 26:27} 又在另一处说：\Quote{与圣洁的人，你将成为圣洁；与完全的人，你将成全；与乖戾的人，你将乖戾。}因此，你们要圣洁地行走，好使你们在主面前更显得配受称赞，而非受仇敌的谴责。\par

% structure: p0101 source=h3 target=subsubsection reason=relative-heading-rank; base=h2

\subsubsection{执事要减轻主教的负担，并亲自处理较小的事务。}

XLIV. 你们主教们要同心合意，彼此和睦；彼此同情，友爱弟兄，以关怀牧养子民；同心教导你们属下之人持守同样的情感，在同样的事上抱有同样的意见，\Quote{使你们中间没有分裂，使你们成为一体、一灵，在同样的心意和同样的判断上完全结合}\BibleRef{格前 1:10}，按照主的任命。执事要将一切事呈报主教，如同基督对圣父所行的一样。但他自己有能力处理的事，就让他亲自处理，从主教那里接受权柄，如同主从圣父那里接受了创造和眷顾的权能。至于重大之事，应由主教裁决；但执事应是主教的耳、目、口、心、魂，使主教不致为许多事务分心，只专注于那些更重大的事，正如耶特洛为梅瑟所设定，并他的建议被采纳一样。\par

% structure: p0103 source=h3 target=subsubsection reason=relative-heading-rank; base=h2

\subsubsection{争论和争吵与基督徒不相称。}

XLV. 所以，基督徒没有与人争辩，实为一种高尚的称誉；但若因某种处理或诱惑而与人发生争辩，就当尽力使之平息，即使因此损失一些东西也在所不惜，并且不可将案件提到外邦人的法庭。事实上，你们甚至不可允许今世的统治者对你们的人民作出判决；因为魔鬼藉此对天主的仆人图谋祸害，给我们带来羞辱，好像我们中间没有一个\Quote{有智慧的人能判断弟兄之间的事}，或裁决他们的争议。\par

% structure: p0105 source=h3 target=subsubsection reason=relative-heading-rank; base=h2

\subsubsection{信友不应在外邦人面前诉讼，也不应召外邦人为证人指控信友。}

XLVI. 因此，不要让外邦人知道你们彼此之间的分歧，也不要接受不信者作为反对你们的证人，不要受他们审判，也不要因进贡或畏惧而欠他们什么；而要\BibleQuote{凯撒的归凯撒，天主的归天主}\BibleRef{玛 22:21}，作为贡品、税银或人口税，正如我们的主藉着交付一枚钱币而免于困扰。因此，宁可选择忍受伤害，并努力追求那些促进和平的事，不仅是在弟兄中间，也是在不信者中间。因为在此世的事务中遭受损失，你们必能确保在虔敬之事上不受损失，并要按基督的命令度虔诚的生活。但如果弟兄之间彼此诉讼——愿天主禁止——你们这些作掌权者应当从中明白，这样的人不是在作主内的弟兄，而是作公敌；并且会发现其中一方是温和、柔顺、光明的子女，而另一方却是无怜悯、傲慢且贪婪的。因此，应当责备受谴责的人，使他被隔离，让他因憎恨弟兄而受惩罚。之后，当他悔改时，应接纳他；这样，当他们学会了智慧，就会减轻你们审判的负担。彼此宽恕过犯也是本分——不是审判者的本分，而是争吵者的本分；正如主在我\Person{伯多禄}问祂时决定的：\BibleQuote{我的弟兄得罪我，我该宽恕他多少次？直到七次吗？}祂回答说：\BibleQuote{我不对你说直到七次，而是七十个七次。}\BibleRef{玛 18:21} 因为主愿意我们真正作祂的门徒，对任何人都从不怀恨；例如，无节制的愤怒、无怜悯的激情、不义的贪婪、不和解的仇恨。你们当藉着教导，使恼怒者转为友谊，使不睦者归于和好。因为主说：\BibleQuote{缔造和平的人是有福的，因为他们要称为天主的子女。}\BibleRef{玛 5:9}\par

% structure: p0107 source=h3 target=subsubsection reason=relative-heading-rank; base=h2

\subsubsection{基督徒的审判应在一周的第二天举行。}

XLVII. 你们的审判应在每周的第二天举行，这样，若你们的判决有任何争议，从安息日之前有一段间隔，你们就能纠正争议，并使那些在主日彼此相争的人和好。执事和长老也应出席你们的审判，作为天主的人，以正义审判，不徇情面。因此，当双方都来到时，按法律所说\BibleRef{申 19:17}，有争议的人应分别站在法庭的中间；你们听了他们之后，应圣洁地投票，在主教判决之前努力使双方成为朋友，以免对犯法者的判决传到世上；要知道，在法庭中，天主的基督作为祂判决的见证者和确证者。但如果有人被任何人控告，他们的名声受损，似乎没有在主内正直行事，你们同样应听取双方——控告者和被控告者；但不带偏见，也不仅听一方，而是按正义，如同对永生或永死作出判决。因为天主说：\BibleQuote{你要公正地追求正义。}\BibleRef{申 16:20} 因为被你们公正地惩罚和隔离的人，被排斥于永生和荣耀之外；他在圣洁的人中变得不光彩，并被天主定罪。\par

% structure: p0109 source=h3 target=subsubsection reason=relative-heading-rank; base=h2

\subsubsection{不应为每一项过犯处以相同惩罚，而应针对不同犯者施以不同惩罚。}

XLVIII. 不要对每一样罪判处相同的刑罚，而要适合每一项罪，以极大的智慧区分各种过犯，大的与小的。对待恶行用一种方式，对待恶言用另一种方式；单纯的恶意又不同。对于侮辱或猜疑也是如此。有些人你们仅以威胁就能约束；有些人你们应以罚款济贫来惩罚；有些人你们应以禁食来克制；还有些人你们应按照其罪过的严重程度予以隔离。因为法律并没有对每一样罪处以相同的刑罚，而是对得罪天主、得罪司祭、得罪圣殿或得罪祭献的罪有不同考虑；对得罪君王、统治者、士兵或同民的罪也不同；同样，对得罪仆人、财产或牲畜的罪也有分别。再者，对故意犯罪与无心过失的估价也不同。因此刑罚也不同：或钉十字架而死，或用石头砸死，或罚款，或鞭打，或遭受自己所加于他人的同样祸害。因此，你们也要对不同过犯分配不同刑罚，免得发生不义，激起天主的义怒。因为你们作为工具作出了何等不义的审判，就要从天主那里承受同样的报应。\BibleQuote{因为你们用什么审判来审判，你们也要受什么审判。}\BibleRef{玛 7:2}\par

% structure: p0111 source=h3 target=subsubsection reason=relative-heading-rank; base=h2

\subsubsection{控告者和证人的品性应如何。}

XLIX．因此，当你坐堂审判，当事人双方都在场时（因为在他们彼此和好之前，我们不会称他们为弟兄），要仔细查问来到你面前的人；首先要查问控告者：这是他第一次控告别人，还是他以前也曾控告过别人？他们的这场争讼和控告是否出于某种争吵？控告者的生活如何？即使他良心清白，也不要只凭他一个人的话就相信，因为这违背法律；他要有其他人与他一同作证，且这些人的生活方式相同。正如法律说：\BibleQuote{凭两三个人的口供，一切事都可成立。}\BibleRef{申 19:15}但我们为什么说证人的人品也需要查问？因为常有两人或多人恶意作证，串通说谎，如同在\Place{巴比伦}的两长老对\Person{苏撒纳}，在\Place{撒玛黎雅}的悖逆之子对\Person{纳波特}，在\Place{耶路撒冷}的犹太群众对我们的主和祂的首位殉道者\Person{斯德望}一样。所以，证人应当温和、不轻易发怒、公正、仁慈、谨慎、节制、远离邪恶、忠信、虔诚；因为这样的人的证词因他们的品德而稳固，因他们的生活方式而真实。至于品性不同的人，不要接受他们的证词，即使他们看似一致地指控被告；因为法律规定：\BibleQuote{不可随众行恶；不可随从虚浮的传言；不可随众屈枉正直。}\BibleRef{出 23:2}你也应当特别了解被告：他的为人及其生活方式如何？他的生活是否有好名声？他是否无可指摘？他是否热心于圣洁？他是否爱护寡妇、爱护客旅、爱护穷人、爱护弟兄？他是否贪不义之财？他是否挥霍无度或浪费成性？他是否清醒、不奢侈，或是个酒徒、贪食者？他是否富有同情心和慈善心？\par

% structure: p0113 source=h3 target=subsubsection reason=relative-heading-rank; base=h2

\subsubsection{过去的过错有时会使后来的指控变得可信。}

L．因为如果被告以前常做恶事，那么现在对他的指控在某种程度上就显得真实，除非公义明显地为他辩护。因为也可能他从前虽曾犯错，但对现在所被指控的罪行却是无辜的。因此，你要极其谨慎地处理这些情形，使你对认罪的犯人所宣告的判决稳妥可靠。如果他被逐出后，请求宽恕，俯伏在主教面前，承认自己的过失，你要接纳他。但也不要让诬告者不受惩罚，以免他诽谤其他善人生活，或鼓励别人效法他。同样，也不要让被定罪的人轻易脱身，以免别人陷入同样的罪行。因为作恶的证人必受惩罚，犯罪的人也必受责罚。\par

% structure: p0115 source=h3 target=subsubsection reason=relative-heading-rank; base=h2

\subsubsection{不可只听一面之词就审判。}

LI．我们之前说过，不应只听一方的话就下判决；因为如果你只听了一方而另一方不在场，无法为对他的指控辩护，你就轻率地投票定他的罪，那么在天主公义的审判者面前，你将被视为那人丧亡的根源，并与诬告者同罪。因为\BibleQuote{如抓住狗尾的人，也是如此主持不义审判的。}\BibleRef{箴 26:17}但如果你效法\Place{巴比伦}的长老，他们作证反对\Person{苏撒纳}，不义地判她死刑，你就要承担他们的判决和定罪。因为主藉\Person{达尼尔}从恶人手中拯救了\Person{苏撒纳}，却将那些染指她血的长老投入火中，并藉\Person{达尼尔}责备你们说：\BibleQuote{你们这些以色列子民竟如此愚昧吗？不经审查，不查明真相，就定了以色列女子的罪吗？你们要回到审判所，因为这些人作了假证反对她。}\par

% structure: p0117 source=h3 target=subsubsection reason=relative-heading-rank; base=h2

\subsubsection{异教法庭在定罪之前所持的谨慎给基督徒提供了良好的榜样。}

LII．再看看世上的法庭：我们看到杀人犯、奸淫者、巫士、盗墓者和窃贼被提交审判；主审官收到控告者的指控后，会询问罪犯是否属实。即使罪犯没有否认罪行，他们也并不立即将他处刑；而是经过数日，在全体会议并拉上帷幔的情况下审查他。那将要宣布死刑判决和投票的人，举起双手向太阳，郑重声明自己对这人的死无涉。尽管他们是外邦人，不认识天主，也不知道天主会因那些公正被定罪的人而追讨人的罪，他们尚且避免不义的判决。\par

% structure: p0119 source=h3 target=subsubsection reason=relative-heading-rank; base=h2

\subsubsection{基督徒不应彼此争执。}

LIII. 但你们既知道我们的天主是谁，也知道祂的判断是如何，你们怎能忍受作出不义的审判呢？因为你们的判决立即会被天主知道。如果你们判断正义，你们将配得正义的赏报，无论现在或将来；如果不义，你们也将遭受同样的待遇。因此，弟兄们，我们劝告你们，宁可赢得天主的称赞，而不是责备；因为天主的称赞对人是永生，正如祂的责备是永死。所以你们要做正义的审判者，和平的缔造者，没有怒气。因为\BibleQuote{凡无缘无故向弟兄发怒的，应受审判。} \BibleRef{玛 5:22} 但假如因某人的阴谋你们对某人生气，\BibleQuote{不可让太阳在你们的愤怒中落下；} \BibleRef{弗 4:26} 因为\Person{达味}说：\BibleQuote{你们应发怒，但不可犯罪；} 也就是说，要尽快和解，免得你们的愤怒持续太久，变成根深蒂固的仇恨，从而产生罪恶。\BibleQuote{心怀仇恨之人的灵魂是致死的，} \Person{撒罗满}说。但我们的主和救主\Person{耶稣基督}在福音中说：\BibleQuote{你若在祭台前献你的礼物，在那里想起你的弟兄有什么怨你的事，就把你的礼物留在祭台前，先去与你的弟兄和好，然后再来献你的礼物给天主。} \BibleRef{玛 5:23} 如今给天主的礼物就是每个人的祈祷和感恩。因此，如果你对弟兄有什么不满，或者他对你有什么不满，你的祈祷便不会被垂听，你的感恩也不会被悦纳，因为那隐藏的愤怒。但弟兄们，你们有责任不断祈祷。然而，因为天主不听那些与弟兄不和睦、有不义争吵的人，即使他们一小时祈祷三次，我们也有责任化解所有的敌意和心胸狭隘，以便我们能以纯洁无玷的心祈祷。因为主命令我们甚至要爱我们的仇敌，绝不可恨我们的朋友。立法者说：\BibleQuote{你不可恨任何人；你不可在心里恨你的弟兄。你应确实谴责你的弟兄，不要因他而犯罪。} \BibleRef{肋 19:17} \BibleQuote{你不可恨埃及人，因为你在他的地作过客。你不可恨厄东人，因为他是你的弟兄。} \BibleRef{申 23:7} \Person{达味}说：\BibleQuote{如果我以恶报那以恶待我的人。} 所以，如果你愿意作一个基督徒，就要遵守主的法律：\BibleQuote{解开一切邪恶的绳索；} \BibleRef{依 58:6} 因为主已赐给你权柄，赦免你弟兄得罪你的罪，直到\BibleQuote{七十个七次}，\BibleRef{玛 18:22} 即四百九十次。因此，你曾多少次赦免你的弟兄，如今却不愿意这样做，你也曾听见\Person{耶肋米亚}说：\BibleQuote{你们各人不要心里怀恨邻人的恶？} \BibleRef{匝 8:17} 但你记恨受伤害，保持敌意，进入审判，怀疑祂的愤怒，你的祈祷受阻。不，如果你赦免了你的弟兄四百九十次，还要更多地增加你的温和行为，为你自己的益处而行善。即使他不这样做，你仍要为了天主而尽力宽恕你的弟兄，\BibleQuote{好成为你们在天之父的子女，} \BibleRef{玛 5:45} 并且当你祈祷时，作为天主的朋友得以垂听。\par

% structure: p0121 source=h3 target=subsubsection reason=relative-heading-rank; base=h2

\subsubsection{主教必须藉其执事提醒民众他们负有和平共处的义务。}

LIV. 所以，主教啊，当你读完读经、咏唱圣咏和圣经训导之后要去祈祷时，让执事站在你旁边，大声说：\Quote{愿没有人彼此争吵；愿没有人带着虚伪而来；} 恐怕你们中间有人发现有争议，他们便良心受责，向天主祈祷，与弟兄和好。因为如果我们进入某人的家，要说：\BibleQuote{愿平安临于此家。} \BibleRef{玛 10:12} 如同平安之子将平安赐予那些相称的人，正如经上所记：\BibleQuote{祂来向你们这些近处的人，以及那些远处的人，凡主所认识属于祂的，宣讲了和平。} 那么，那些进入天主的教会的人，更当首先为天主的平安祈祷。但如果他为别人祈求平安，他自己更当在其中，如同光明之子；因为内中没有平安的人不配将平安赐予他人。因此，首要的是我们必须在自己的内心享有平安；因为那在自己心中找不到纷乱的人，不会与别人争吵，而是和平、友善，聚集主的子民，与祂同工，以增加在同心合意中得救的人数。因为那些策划敌意、纷争、争竞和诉讼的人，是邪恶的，与天主疏远。\par

% structure: p0123 source=h3 target=subsubsection reason=relative-heading-rank; base=h2

\subsubsection{列举天主眷顾的若干实例，以及从世界起初的每一个时代，天主如何召叫众人悔改。}

LV. 因为天主从起初就是慈悲的天主，藉义人和先知召叫每一代悔改。祂藉\Person{亚伯尔}、\Person{塞姆}和\Person{舍特}，也藉\Person{厄诺斯}和曾被提升的\Person{哈诺客}，教导洪水以前的人；藉\Person{诺亚}教导洪水时期的人；藉好客的\Person{罗特}教导\Place{索多玛}的居民；藉\Person{默基瑟德}、圣祖们和天主所爱的\Person{约伯}教导洪水以后的人；藉\Person{梅瑟}教导埃及人；藉\Person{梅瑟}、\Person{若苏厄}、\Person{加肋布}、\Person{丕乃哈斯}和其余的人教导以色列人；藉天使和先知教导律法以后的人，又藉祂由\Person{童贞玛利亚}而有的\OriginalTerm{降生成人}{Incarnation}教导同一批人；藉祂的前驱\Person{若翰}教导祂肉身显现前不久的人，又在基督诞生后藉同一人教导他们，说：\BibleQuote{你们悔改吧！因为天国临近了；}\BibleRef{玛 3:2}祂受难以后，藉我们十二宗徒和蒙选的器皿\Person{保禄}教导那些后来的人。因此，我们蒙恩成为祂显现的见证人，同我们的主的兄弟\Person{雅各伯}、其他七十二位门徒以及祂的七位执事一起，亲耳从我们的主耶稣基督口中听到，并凭着确切的知识宣告：\BibleQuote{什么是天主的旨意，就是那善良、可喜悦和成全的旨意}\BibleRef{罗 12:2}这旨意已藉耶稣启示给我们，使无人丧亡，而所有人一心相信祂，一致赞美祂，并因此永远活着。\par

% structure: p0125 source=h3 target=subsubsection reason=relative-heading-rank; base=h2

\subsubsection{论天主的旨意是要人在信仰事宜上同众天使一心一意。}

LVI. 因为这就是我们的主教导我们在祈祷时向祂的父所说的：\BibleQuote{愿祢的旨意承行于地，如于天焉；}\BibleRef{玛 6:10}好使正如无形之掌权者的天上本性都一致光荣天主，同样在地上，所有的人都用同一的口和同一的心意，藉祂的\OriginalTerm{独生子}{only-begotten}基督，光荣那唯一、独一和真实的天主。因此，祂的旨意是要人同心赞美祂，一致朝拜祂。因为在基督内，祂的旨意是要藉祂得救的人众多；但你不要给祂或教会带来任何损失或减少，也不要因你而毁掉哪怕一个灵魂，那本可藉悔改得救；因此，那灵魂不仅因自己的罪而丧亡，也因你的背信而丧亡，你因此应验了经上所记：\BibleQuote{不与我聚集的，就是分散的。}\BibleRef{玛 12:30}这样的人是羊群的驱散者，是仇敌，是天主的敌人，是那些羔羊的毁灭者，他们的牧者是主，而我们是从各民族和各方言中经过许多劳苦、危险、不断的劳作、守夜、禁食、躺卧地上、迫害、鞭打、监禁，才聚集了他们，为承行天主的旨意，并用欢欣和蒙选的人充满宴席厅，使他们坐在祂的桌前，那就是圣而公的教会，向那藉我们召叫他们进入生命的天主歌唱赞美。而你，尽其所能地驱散了他们。你们平信徒也要彼此和睦，像智慧人一样努力增加教会，并挽回、驯服和修复那些看似野性的人。因为按照天主的应许，这是最大的赏报：\BibleQuote{你若将贵重与卑贱的分开，你就可以作我的口。}\BibleRef{耶 15:19}\par

% structure: p0127 source=h2 target=subsection reason=relative-heading-rank; base=h2

\subsection{第七节　论在教会中聚会}

% structure: p0128 source=h3 target=subsubsection reason=relative-heading-rank; base=h2

\subsubsection{论教会与圣职人员的准确描述，以及圣职与平信徒在隆重聚会中为敬拜天主每人所当行之事}

LVII. 但是，主教啊，你当成为圣洁、无可指摘、不斗殴、不轻易发怒、不残忍；而是建设者、劝化者、善于教导、容忍邪恶、温和、良善、忍耐、乐于劝勉、乐于安慰，如同天主的人。\par

当你召集教会聚会时，如同大船的指挥官，要尽可能精心安排聚会，吩咐执事像水手一样，以应有的谨慎和庄重为弟兄们预备座位，如同为乘客预备。首先，建筑物应为长方形，其头部朝东，在东端两侧设有更衣室，这样就像一艘船。在中央安放主教宝座，在他两旁让司铎团坐下；执事穿着紧窄的束腰衣服，站在近旁，因为他们如同船上的水手和舵手；关于平信徒，让他们安静有序地坐在另一边。妇女应自行就座，也要保持静默。在中间，让读经员站在高处：让他诵读\BibleBook{梅瑟}{五书}、农的儿子\BibleBook{若苏厄}{书}、\BibleBook{民长}{纪}、\BibleBook{列王}{纪}、\BibleBook{编年}{纪}，以及被掳归回后所写的书；此外，还有\BibleBook{约伯}{书}、\BibleBook{撒罗满}{书}和十六位先知的书。当两段读经分别诵读之后，让另一人歌唱\Person{达味}的\BibleBook{达味}{圣咏集}，众人在诗节结尾处应和。此后，诵读我们的\BibleBook{宗徒}{大事录}，以及我们的同工\Person{保禄}在圣神引导下写给各教会的\BibleBook{保禄}{书信}；再后，让一位执事或司铎诵读\BibleBook{}{福音}，即我、\Person{玛窦}和\Person{若望}所传给你们的，以及\Person{保禄}的同工\Person{路加}和\Person{玛尔谷}所领受并留给你们的。诵读福音时，所有司铎、执事和全体百姓都当站立，极其静默，因为经上记载：\BibleQuote{以色列，你要静默细听。}\BibleRef{申 27:9}又说：\BibleQuote{你且站在那里，听。}\BibleRef{申 5:31}接着，让司铎们依次（而非全体同时）劝勉百姓，最后是主教，作为指挥官。守门者应站在男子的入口处，观察他们；女执事也站在妇女的入口处，如同船员。因为在会幕和天主的殿中，都有同样的描述和样式。若发现有人坐在不属于他的位置上，执事如同船头管理员应予以斥责，并带他到适当的位置；因为教会不仅像一艘船，也像一个羊栈。牧人将各类牲畜，我指山羊和绵羊，按种类和年龄分别安置，但每一类仍与同类聚集；教会中也当如此。青年应单独就座，若有位置；若没有，便站着。年长者则依次就座。对于站着的儿童，让他们的父母领到身边。较年轻的妇女也当单独就座，若有位置；若没有，便站在妇女后面。已婚并有子女的妇女应单独安排；但童贞女、寡妇和年长妇人应在所有人之前站立或就座；让执事安排座位，使每个进来的人都能到自己的位置，不致坐在入口处。同样，执事监督众人，无人窃语、打盹、嬉笑或点头；因为所有人在圣堂中都应智慧、清醒、警醒地站立，专注于主的言语。此后，等慕道者和补赎者退出后，全体同心合意地站立，面向东方，向天主祈祷，因为祂曾升到至高之天，朝向东方；同时也纪念伊甸园古时位于东方，第一个人因听从蛇的诱惑、违抗天主的命令而从那里被驱逐。至于执事，祈祷结束后，让其中一些服侍圣体圣事的奉献，以敬畏之心服侍主的身体；另一些看管群众，使他们保持安静。但站在大司祭旁边的执事应向众人说：\BibleQuote{你们中任何人不可与另一人有争执；不可有人假冒伪善而来。}然后，男子与男子，妇女与妇女，互行主的亲吻。但不可用诡诈行之，如同\Person{犹达斯}以亲吻出卖主一样。之后，让执事为整个教会、整个世界及其各部分与果实祈祷；为司铎和统治者、为大司祭和君王、并为宇宙的和平祈祷。此后，让大司祭为百姓祈求和平，并祝福他们，如同\Person{梅瑟}吩咐司祭祝福百姓的话：\BibleQuote{愿主祝福你，保护你；愿主使祂的面容光照你，赐你平安。}让主教也为百姓祈祷说：\BibleQuote{主，请拯救祢的百姓，祝福祢的产业，就是祢以基督的宝血所获得的，并呼召为君尊的司祭、圣洁的子民。}之后，举行祭献，百姓站立，默祷；奉献完成后，各等级依次领受主的身体和宝血，怀着敬畏和神圣的恐惧，如同到他们君王的身体前。妇女应蒙头前来，合乎妇女的秩序；但要看守门口，免得任何不信者或尚未领受入门圣事的人进入。\par

% structure: p0131 source=h3 target=subsubsection reason=relative-heading-rank; base=h2

\subsubsection{论为旅居者、平信徒、圣职人员和主教的推荐信；以及来到教会聚会的人，无论其身份地位均应被接纳。}

LVIII. 若有弟兄，或男或女，从别的堂区来，携带推荐信，执事应处理此事，查问他们是否属于信徒，是否属于教会？是否曾受异端玷污？此外，该人是已婚妇女还是寡妇？当他在这些问题上得到满意答复，确认他们确实是信徒，在主的事上持相同意见，就应将每人领到适合他的位置。若有司铎从别的堂区来，应由司铎们接纳他共融；若有执事来，由执事们接纳；若有主教来，应与主教同坐，并获准与他同等的尊荣。你，主教啊，应请他向百姓宣讲训导的话；因为外来的劝勉和告诫极为可悦，且大有裨益。正如圣经所说：\Quote{没有先知在自己家乡是被悦纳的。}你还应准许他奉献圣体圣事；但若他出于对你的敬意，且作为明智之人，为保全属于你的尊荣而不愿奉献，你至少应强迫他给百姓祝福。但是，若在会众坐下之后，另有世俗中仪态端庄、品行良好的人进来，无论是陌生人还是本地人，你，主教啊，若你正在宣讲天主的圣言，或正在听那歌唱或诵读的人，不可因接受人的情面而离开圣言的职务，去为他安排上座；而应保持安静，不打断你的讲论，也不中断你的专注。但弟兄们应通过执事接纳他；若没有空位，执事应说话（但不可发怒），叫较年轻的起来，让那陌生人坐在那里。那爱弟兄的人按理应自愿如此行；但他若拒绝，就应强行拉他起来，让他坐在最后面，好使其余人学会给更有尊荣的人让座。不仅如此，若有穷人、或微贱家庭的人、或陌生人进来，无论老少，又没有空位，执事也应尽心尽力为他们找个位置；这样，他不仅不按外貌待人，而且他对天主的服务必蒙悦纳。女执事对待来见她们的妇女，无论贫富，也应同样行事。\par

% structure: p0133 source=h3 target=subsubsection reason=relative-heading-rank; base=h2

\subsubsection{每位基督徒应当殷勤地每日早晚前往教会。}

LIX. 当你训导百姓时，主教啊，要命令并劝勉他们每日早晚恒常前往教会，绝不可因任何理由离弃教会，而要持续聚在一起；不可因缺席而削弱教会，使基督的身体缺少肢体。因为这话不仅是对司铎说的，也让每位平信徒耳闻而以为己任，想到主曾说：\Quote{不随同我的，就是反对我；不同我收集的，就是分散。}\BibleRef{玛 12:30} 你们这些基督的肢体，既因祂的应许而有基督为元首，祂临在并通传给你们，就不要因不聚集而分散。不要疏忽自己，也不要剥夺救主祂自己的肢体，不要分裂祂的身体，不要分散祂的肢体，不要将今世的事务置于天主圣言之上；而要每天早晚在主殿中聚集，歌唱圣咏并祈祷：早晨诵念第六十二篇圣咏，晚上诵念第一百四十篇，但尤其要在安息日。在我们主复活的日子，即主日，要更勤勉地聚集，向那藉耶稣创造万有、并派遣祂到我们这里、且屈尊让祂受苦、又从死者中复活祂的天主献上赞美。不然，那在这日不聚集聆听关于复活的救恩之言的人，将向天主作何推诿呢？在这日，我们站着祈祷三次，以纪念那三日复活的祂；在这日，有读先知书、宣讲福音、奉献祭献、圣食的恩赐？\par

% structure: p0135 source=h3 target=subsubsection reason=relative-heading-rank; base=h2

\subsubsection{外邦人和犹太人在频繁前往他们的庙宇和会堂上所表现的空虚热忱，是激励基督徒前往教会的恰当榜样和动机。}

六十．他昼夜为暂时之事操劳，却不关心永恒之事，怎能不与天主为敌？他每日关心洗涤和暂时的食物，却不关心那存留永远的事物？这样的人如今怎能避免听到主的那句话：\BibleQuote{外邦人比你们更正义？}\BibleRef{则 16:52} 正如祂责备耶路撒冷时所说的：\Quote{索多玛比你们更正义。} 因为外邦人每日睡醒后，就跑到他们的偶像前敬拜，在所有工作和劳苦之前，首先向它们祈祷；在他们的节庆和隆重集会中，他们并不缺席，反而参与；不仅在场的人，连远方的人也是如此；在他们的公共表演中，所有人都聚集在一起，如同进到会堂里。同样，那些徒有虚名的犹太人，工作了六天之后，在第七天安息，聚集到他们的会堂，从未离开或疏忽安息和聚会。然而，他们因不信而被剥夺了圣言的能力，甚至被剥夺了他们自称的\OriginalTerm{犹大}{Judah}这一名称的力量——因为\OriginalTerm{犹大}{Judah}被解释为\OriginalTerm{宣认}{Confession}——但他们并没有向天主宣认（他们不公正地导致了十字架的苦难），以致于无法在悔改中得救。因此，如果那些不得救的人尚且经常聚集为无益之事，那么，你离弃天主的教会，不效法外邦人，反而因缺席而变得懒惰、背教或作恶，你将向主天主作何辩解？对此，主藉\Person{耶肋米亚}说：\Quote{你们没有遵守我的法令；不但如此，你们也没有按照外邦人的法令行事，反而在某种程度上超过了他们。} 又说：\BibleQuote{以色列比背信的犹大更显正义。}\BibleRef{耶 3:11} 随后又说：\Quote{外邦人岂能更换他们的神祇，那些不是真神的？因此，你们过海到\Place{基提}海岛观看，派人到\Place{刻达尔}去，仔细观察是否发生过这样的事。因为这些民族没有更改他们的法令，} 但祂说：\Quote{我的百姓却将自己的光荣换成了无益之物。} 因此，凡是藐视或缺席天主教会的人，将如何为自己辩护呢？\par

% structure: p0137 source=h3 target=subsubsection reason=relative-heading-rank; base=h2

\subsubsection{我们不可将今世的事务置于天主敬礼之上。}

但是，若有人以工作为借口而轻视聚会，\Quote{为自己的罪过寻找托辞，} 这样的人要知道：信友的职业是副业，天主的敬礼才是他们的大业。因此，你们要以职业为副业，以维持生计，但要把天主的敬礼当作主业，正如我们的主所说：\BibleQuote{不要为那能够朽坏的食粮劳碌，而要为那存留直到永生的食粮操劳。}\BibleRef{若 6:27} 又说：\BibleQuote{天主的工程就是让你们相信祂所派遣的那一位。}\BibleRef{若 6:29} 因此，努力不要离开天主的教会；但若有人忽视教会，而去到外邦人污秽的庙宇，或犹太人或异端者的会堂，这样的人在审判之日将向天主作何辩解？他离弃了永生天主的神谕，那些生活且赐生命的神谕，能够拯救人脱离永罚，却进入了魔鬼的房屋，或到杀害基督者的会堂，或恶人的集会？——不听从那说：\Quote{我恼恨恶人的集会，绝不与不敬虔的人同席；我没有与虚妄的集会同坐，也决不与非虔诚者同坐。} 又说：\Quote{不随从不敬虔者的计谋、不站在罪人的道路、不坐在轻蔑者的座位，而是专心喜爱上主的法律，日夜默思祂法律的人是有福的。} 但你却离弃信友的聚集、天主的教会和祂的法律，而尊重那些\Quote{贼窝，} 称祂称为不洁之物为圣，使祂所圣化的东西成为不洁。不仅如此，你还追逐外邦人的虚浮，奔向他们的剧院，渴望被认为是进到其中的人，参与不雅、甚至可憎的言语，不听从\Person{耶肋米亚}所说：\BibleQuote{上主，我没有坐在他们的集会中，因为他们是轻慢者；但我因祢的手而惧怕。}\BibleRef{耶 15:17} 也不听从\Person{约伯}类似地说：\BibleQuote{我若曾与轻慢者同行，就让我被公义的天平称量。}\BibleRef{约 31:5} 但你为何要参与外邦人的神谕，它们不过是死人藉魔鬼的灵感说出致命的事，趋向颠覆信德，并引诱那些听从的人走向多神论？因此，凡听从天主法律的人，你们要珍视这些法律过于今世的需要，更加尊重它们，并一同跑到主的教会，\Quote{那是祂用祂所爱的、首生的基督的血所购买的。} 因为教会是至高者的女儿，她藉恩宠之道为你们孕育，并\Quote{在你们内形成了基督，} 你们成为有分于祂的人，因而成为祂圣洁、特选的肢体，\BibleQuote{没有瑕疵或皱纹，或其他任何缺点；但你们在信德上是圣洁无瑕的，在祂内得以完全，按照那创造你们的天主的肖像。}\BibleRef{弗 5:27}\par

% structure: p0139 source=h3 target=subsubsection reason=relative-heading-rank; base=h2

\subsubsection{基督徒必须远离外邦人的一切不虔敬行为。}

第六十二章。所以，你们要注意，不要在敬拜中与那些丧亡的人联合，也就是外邦人的集会，这会导致你们的欺骗和毁灭。因为天主与魔鬼没有相通；因为那与喜爱魔鬼之事的人聚集的，将被视为他们中的一员，并将继承灾祸。也要避免不雅的表演：我指的是戏院和异教徒的排场；他们的邪术、观察预兆、占卜、净化、算命、观鸟；他们的招魂术和呼召。因为经上记载：\Quote{在雅各伯中没有占卜，在以色列中也没有卜卦。} \BibleRef{户 23:23} 又说：\Quote{占卜是不义。} 另外：\Quote{你们不可占卜，也不可追随观察预兆的人，不可行卜卦，不可与交鬼的人来往。不可让巫师存活。} \BibleRef{肋 19:26} 因此耶肋米亚劝诫说：\Quote{不要随从外邦人的道路，也不要畏惧天上的异象。} 所以，信徒的责任是避免不敬神者的集会，包括外邦人、犹太人以及其他异端者，免得我们与他们联合，给自己的灵魂设下陷阱；免得我们参加他们的宴席（这些宴席是为敬拜魔鬼而举行的），从而与他们一同在亵渎中。你们也要避免他们的公共集会，以及其中举行的运动。因为信徒不应当参加这些公共集会，除非是为了购买一个奴隶，拯救一个灵魂，同时也为了购买其他生活必需品。因此，要远离所有偶像崇拜的排场和盛况，所有他们的公共集会、宴会、决斗，以及一切属于魔鬼的表演。\par

% structure: p0141 source=h2 target=subsection reason=relative-heading-rank; base=h2

\subsection{第八节 论劳动谋生的责任}

% structure: p0142 source=h3 target=subsubsection reason=relative-heading-rank; base=h2

\subsubsection{不劳动的基督徒不得吃饭——正如伯多禄及其他宗徒是渔夫，保禄和阿桂拉是制帐篷的，雅各伯的儿子犹达是农夫。}

第六十三章。让教会的年轻人努力殷勤地服事一切所需；以相称的庄重态度处理你们的事务，这样你们就能常常有足够的来养活自己和那些有需要的人，而不加重天主的教会。因为我们自己，除了注意福音的话，也并不忽略我们卑微的工作。因为我们中有些人是渔夫，有些是制帐篷的，有些是农夫，这样我们就从不闲散。正如撒罗满在某处说：\Quote{懒惰的人，去看看蚂蚁，关注它的行径，你就比它更有智慧。因为它没有田地、监督或统治者，在夏天预备食物，在收割时积蓄大量。或者去看看蜜蜂，学习它多么勤劳，它的工作多么有价值，它的劳苦连君王和百姓都用于健康。它虽力量微弱，却因尊重智慧而得到提升等等。懒惰的人，你要在床上躺到几时？你何时才从睡眠中醒来？你睡一会儿，躺一会儿，打盹一会儿，把手交叉在胸前睡一会儿。然后贫穷像凶恶的旅客一样临到你，匮乏像快跑的赛手。但如果你勤勉，你的收获将如泉涌，匮乏将像凶恶的逃兵远离你。} 又说：\Quote{耕种自己田地的，必得饱食。} \BibleRef{箴 12:11} 此外，他处说：\Quote{懒惰人把手交叉在一起，吃自己的肉。} \BibleRef{训 4:5} 之后：\Quote{懒惰人把手藏在怀里，甚至不能送进口里。} \BibleRef{箴 19:24} 又说：\Quote{由于手的懒惰，房屋会倒塌。} \BibleRef{训 10:18} 所以要不断劳动；因为懒惰的污点不被治愈。但\Quote{若有人不愿工作，就不得吃饭} \BibleRef{得后 3:10} 在你们中间。因为主我们的天主憎恶懒惰。献身于天主的人没有一个应当闲散。\par

\SourceRecord{Translated by James Donaldson.；Ante-Nicene Fathers；Vol. 7.；Edited by Alexander Roberts, James Donaldson, and A. Cleveland Coxe.；Buffalo, NY: Christian Literature Publishing Co.,；1886.；<http://www.newadvent.org/fathers/07152.htm>.}

% structure: p0001 source=h1 target=section reason=page-title

\section{\WorkTitle{宗徒宪章}（第三卷）}

% structure: p0002 source=h2 target=subsection reason=relative-heading-rank; base=h2

\subsection{第一节 论寡妇}

% structure: p0003 source=h3 target=subsubsection reason=relative-heading-rank; base=h2

\subsubsection{应选立之寡妇的年龄}

一、应选立\Quote{年纪不低于六十岁的寡妇}，这样可借其年龄在某种程度上防止再婚的嫌疑。但若准许较年轻的女子进入寡妇之列，而她因年轻不能忍受寡居而再婚，她将给寡妇之列的光荣带来不光彩的批评，并要向天主交账；并非因为她再婚，而是因为她\BibleQuote{背离了基督而放纵情欲}\BibleRef{弟前 5:11}，且没有遵守她的承诺，因为她没有怀着信德和敬畏天主的心来守约。因此，这样的承诺不应轻率作出，而应非常谨慎：\BibleQuote{她若不发愿，比发愿而不还愿更好}\BibleRef{训 5:5}。但若有较年轻的妇女，与丈夫生活不久，因死亡或其它原因失去了丈夫，独自生活，拥有寡居的恩赐，她将被视为有福的，如同属漆冬的匝尔法特的寡妇，天主的圣先知厄里亚曾借宿在她那里\BibleRef{列上 17:9}。这样的人也可比作\Quote{亚纳，是法奴耳的女儿，属阿协尔支派，她不离开圣殿，昼夜以禁食和祈祷事奉天主，现已八十四岁，从童贞出嫁后同丈夫生活了七年，她赞美基督的来临，感谢主，并向所有期待以色列得救赎的人讲论祂。}这样的寡妇将有好名声，并被尊荣，既在地上在人前有荣耀，也在天上在天主前有永久的赞美。\par

% structure: p0005 source=h3 target=subsubsection reason=relative-heading-rank; base=h2

\subsubsection{必须避免选立较年轻的寡妇，因有嫌疑。}

二、不要让较年轻的寡妇列入寡妇之列，免得她们在青春年华，借口无力自守而再婚，并受到指责。但应帮助和支持她们，免得她们借口被遗弃而再婚，以致陷入不当的责难。因为你们应当知道，依法结婚一次是义的，因为是符合天主旨意；但在许诺之后再婚则是邪恶的，并非因婚姻本身，而是因失信。第三次婚姻是不节制的迹象。超过三次的婚姻则是明显的邪淫和无疑的不洁。因为天主在创造中赐予一男一女；因为\BibleQuote{两人成为一体}\BibleRef{创 2:24}。但对于年轻妇女，在第一个丈夫去世后，允许再婚，免得她们陷入魔鬼的谴责和许多罗网及愚昧的欲望，这些欲望伤害灵魂，给她们带来惩罚而非安息。\par

% structure: p0007 source=h3 target=subsubsection reason=relative-heading-rank; base=h2

\subsubsection{寡妇应具有的品性，以及主教应如何支持她们。}

三、真正的寡妇是那些只有一位丈夫，在众人中有好名声的，因善行受到称赞；她们确实是寡妇，清醒、贞洁、忠信、虔诚，好好教养了子女，无可指责地款待了旅客，应被支持作为奉献于天主的人。此外，主教啊，要记念贫困者，伸出援手，供应他们所需，作天主的管家，按时将祭品分给他们每人：寡妇、孤儿、无人照顾者、以及受痛苦考验的人。\par

% structure: p0009 source=h3 target=subsubsection reason=relative-heading-rank; base=h2

\subsubsection{我们应当对各类有需要的人施行爱德。}

四、倘若有人既非寡妇也非鳏夫，却因贫穷、疾病或抚养众多子女而需要援助，那该怎么办？你们有责任监督所有人，并照顾他们所有人。因为施舍的人并非自行决定赠予寡妇，而只是带来礼物，称之为自由奉献，这样，你作为知道谁在患难中的好管家，可将礼物分给他们当得的分。因为天主知道施予者，即使在你不在时你分发给需要的人。他获得善行的赏报，而你则因凭良心施与而有福。但你要告诉他们施予者是谁，使他们可以指名为他祈祷。因为我们的责任是对所有人行善，不可有偏爱，无论他们是谁。因为主说：\Quote{凡求你的，就给他。}显然这是指真正有需要的人，无论是友是敌，是亲属是外人，单身或已婚。因为在整个圣经中，主都给我们关于穷人的劝诫，首先通过依撒意亚说：\Quote{把你的饼分给饥饿的人，将无家可归的穷人接到你家中。你若见赤身露体的人，就遮盖他；不可忽略你本族和同种的穷人。}然后通过达尼尔对掌权者说：\Quote{为此，大王啊！请你接纳我的劝告：用施舍补赎你的罪过，用怜悯穷人弥补你的过犯。}又通过撒罗满说：\Quote{仁慈和忠信补赎罪过。}又通过达味说：\Quote{眷顾贫穷和困苦的人是有福的；上主在患难之日必拯救他。}又说：\Quote{他施舍钱财，周济穷人；他的仁义永存。}撒罗满又说：\BibleQuote{怜悯穷人的，就是借给上主；他的慷慨，上主必偿还。}\BibleRef{箴 19:17}随后又说：\BibleQuote{塞耳不听穷人哀呼的，他将来呼求也不蒙应允。}\BibleRef{箴 21:13}\par

% structure: p0011 source=h3 target=subsubsection reason=relative-heading-rank; base=h2

\subsubsection{寡妇应非常谨慎自己的行为。}

五、愿每位寡妇温良、安静、温和、真诚、不轻易发怒、不多言、不喧哗、不言语急躁、不诽谤、不吹毛求疵、不一口两舌、不多管闲事。若她看见或听见任何不当之事，愿她如同不见、如同不闻。愿寡妇只专心为施舍者并为整个教会祈祷；若有人向她询问任何事，愿她不轻易回答，除非是有关信德、正义和对天主的望德的问题；将那些渴望受教于虔诚教义的人交托给管理者。愿她只作可推翻\OriginalTerm{多神论}{polytheism}错误的回答，并阐明有关\OriginalTerm{天主的独一统御}{monarchy of God}的主张。至于其余教义，愿她不轻率回答任何事，以免因无知而说出什么，使话语被人亵渎。因为主教导我们，话语就像\BibleQuote{芥菜种子}\BibleRef{玛 13:31}，具有火热的性质，若有人不善用，必觉其苦。因为在奥秘的事上，我们不可轻率，而应谨慎；因为主劝诫我们说：\BibleQuote{不要把你们的珍珠丢在猪前，恐怕它们用脚践踏，又转过来撕咬你们。}\BibleRef{玛 7:6}因为不信的人听到关于基督的教义没有被恰当解释，而是有缺陷时，尤其关于祂的降生成人或是祂的受难，他们宁可以轻蔑拒绝、讥笑为虚假，也不因此赞美天主。如此，老年妇女将因轻率和引起亵渎而获罪，并承继祸哉。因为祂说：\BibleQuote{因我的名在外邦人中受亵渎的人有祸了。}\BibleRef{依 52:5}\par

% structure: p0013 source=h3 target=subsubsection reason=relative-heading-rank; base=h2

\subsubsection{妇女不可施教，因其不宜；以及哪些妇女跟随了我们的主。}

六、我们不允许我们的\BibleQuote{妇女在教会中施教}\BibleRef{格前 14:34}，只允许她们祈祷并听受教者；因为我们的师傅和主耶稣自己，当祂派遣我们十二人去使民众和万邦成为门徒时，祂并没有差遣妇女宣讲，虽然祂并不缺这样的人员。因为我们中间有我们主的母亲和祂的姐妹们；还有玛利亚玛达肋纳、雅各伯的母亲玛利亚、拉匝禄的姐妹玛尔大和玛利亚、撒罗默，以及其他一些人。因为，若有必要让妇女施教，祂自己必会先命令这些妇女与我们一同教导百姓。因为\BibleQuote{男人是女人的头}\BibleRef{格前 11:3}，身体的其余部分治理头是不合理的。因此，愿寡妇承认自己是\Quote{天主的祭台}，并愿她坐在自己家中，不以任何借口进入信徒的家中收取东西；因为天主的祭台从不四处奔跑，而是固定在一处。因此，愿童贞女和寡妇如此：不四处奔跑，不游荡到那些不信者家中。因为这样的人是游荡者、厚颜无耻者；她们不让自己的脚安息在一处，因为她们不是寡妇，而是随时准备收受的钱囊、搬弄是非者、诽谤者、争吵的鼓动者、无耻、厚颜，这样的人不配召叫她们的那一位。因为她们不在主日来到会众的共同聚集地，如同警醒的人；而是要么打盹，要么闲谈，要么诱惑人，要么乞讨，要么陷害别人，将人带到恶者那里；不让她们在主内警醒，却操心使她们出去时如进来时一样空虚，因为她们不听天主的圣言被教导或被诵读。因为先知依撒意亚论到这样的人说：\BibleQuote{你们听了又听，却不明白；看了又看，却不领悟；因为这百姓的心变得迟钝，耳重听。}\BibleRef{依 6:9}\par

% structure: p0015 source=h3 target=subsubsection reason=relative-heading-rank; base=h2

\subsubsection{所谓虚假寡妇的特征。}

VII. 同样，这些寡妇的心耳被堵塞，以至于她们不坐在屋内与主交谈，反而四处奔走以求获得，并通过愚昧的闲谈来满足敌人的欲望。因此，这样的寡妇并不依附于基督的祭台：因为有些寡妇视获利为她们的营生；她们无耻地索取，无度地接受，反而使众人更不愿施舍。因为她们本应以教会供给的必需品为满足，节制欲望，但相反，她们从邻舍家跑到另一家，打扰他们，为自己积攒大量钱财，并以苛刻的高利贷放债，只关心\OriginalTerm{玛门}{mammon}，她们的金钱袋就是她们的神；她们将吃喝置于一切美德之上，说：\Quote{我们吃喝吧，因为明天我们就要死了。}她们视这些事物为持久而不朽的。因为那习惯于只谈论金钱的妇人，敬拜玛门而非天主——即，她是获利的奴仆，不能中悦天主，也不能顺服于祂的敬拜；无法持续地向祂代祷，因为她的心思意念追逐金钱：因为\Quote{你的财宝在哪里，你的心也必在那里。}因为她心中思想：她要去哪里接受施舍，或某个女友忘了她，她有些话要对她说。思虑这些事的妇人，将不再专注于祈祷，而是专注于那浮现的念头；因此，即使她有时为某人祈祷，也不会被垂听，因为她并非全心向主呈上恳求，而是心思分散。但那专注于天主的妇人将坐在屋内，昼夜思念主的事，以不停预备的口唇献上真诚的恳求。因此，正如以智慧闻名、以谦逊著称的\Person{友弟德}，\Quote{日夜为以色列向天主祈祷}，同样，与她相似的寡妇也将不停为教会向天主献上代祷。主必垂听她，因为她心意专一，不贪得无厌，不贪婪，不浪费；她的眼睛纯洁，耳朵清洁，双手无污，双足安静，口唇既不贪食也不轻率，而是说合宜的话，只取维持生计的必要之物。如此，她端庄稳重，不给人麻烦，必中悦天主；她一旦祈求，恩赐必临到她：正如祂说：\Quote{你还在说话时，我要说：我在这里。}愿这样的寡妇也不贪爱钱财，不傲慢，不贪不义之财，不贪得无厌，不贪食，而是自制、温良、不惊扰任何人、虔诚、谦逊、坐在家中、歌唱、祈祷、诵读、警醒、禁食；以诗歌和赞歌不断与天主交谈。愿她取羊毛，并帮助别人，而非自己需要他人的帮助；铭记那在福音中蒙主亲作见证的寡妇，她进入圣殿，\Quote{投了两个小钱，即一文铜钱的四分之一。我们的主基督、主和人心检察者看见她，便说：我实在告诉你们：这穷寡妇投入圣库的，比众人投的都多；因为众人都拿他们所余的来投，但这寡妇却由自己的不足中，把她所有的一切生活费都投进去了。}\par

因此，寡妇应当端庄、顺服她们的主教、司铎和执事，此外还有女执事，以虔诚、恭敬和敬畏之心；不擅自专权，也不愿在未得执事同意的情况下，做任何超越规定的事，例如去某人那里吃喝，或从任何人那里接受任何东西。但若她未经指示而做了其中任何一件事，就应以禁食受罚，否则因她的鲁莽而被隔离。\par

% structure: p0018 source=h3 target=subsubsection reason=relative-heading-rank; base=h2

\subsubsection{寡妇不应从不配的人那里接受施舍，如同主教或任何其他信徒不应如此。}

VIII. 因为这样的寡妇如何知道那施舍她的人品性如何？或他赖以供给食物的服务来自何种来源，是否来自劫掠或其他不良生活方式？寡妇却不记得，若她以不配天主的方式接受，就必须为每一件事交账。因为司铎们也绝不应从这样的人——例如，从劫掠者或娼妓——那里接受自愿奉献。因为经上记载：\Quote{不可贪恋你邻人的财物。}以及\Quote{不可将娼妓的所得献于上主你的天主。}这样的人的奉献，不应被接受，甚至那些被逐出教会的人的奉献也不该接受。寡妇也应随时服从上级给他们的命令，并按照主教的安排行事，服从他如同服从天主；因为那从应受责备或被逐出教会的人那里接受施舍，并为他祈祷，而他仍执意行恶，不愿悔改，那人就在祈祷中与他共融，使基督忧伤，因为基督弃绝不义之人，并藉那不配的奉献坚固他们，又与他们同玷污，不让他们悔改，以至在神面前哀哭祈祷。\par

% structure: p0020 source=h3 target=subsubsection reason=relative-heading-rank; base=h2

\subsubsection{妇女不应付洗，因为这是不虔敬的，且违背基督的道理。}

IX. 至於婦女施洗之事，我們讓你們知道，凡從事此事的，沒有小危險。因此我們不勸你們這樣做；因為這很危險，甚至可說是邪惡而不敬神的。因為若\Quote{男人是女人的頭}，而他原被立於司祭職，則廢除受造之秩序，離棄首要而趨向身體的極端部分，便不合理。因為女人是男人的身體，取自他的肋旁，順服於他，為生育子女而從他分離。因為祂說：\BibleQuote{他要管轄你。}\BibleRef{創 3:16} 因為女人的首要部分是男人，作為她的頭。但若在前面的憲章中，我們不許她們教導，又怎能有人准許她們違反本性，執行司祭的職務呢？因為這是外邦無神論的愚昧作法之一，即為女性神祇任命女司祭，而非基督的憲章之一。因為若聖洗要由婦女施行，我們的主當然會由自己的母親施洗，而非由若翰；或當祂派遣我們去施洗時，祂也會為此目的派遣婦女與我們同去。但如今，祂既未藉憲章亦未藉書寫，在任何地方將這樣的事交付給我們，因為祂知道受造的本性與行動的端莊；祂是受造的本性與憲章的立法者。\par

% structure: p0022 source=h3 target=subsubsection reason=relative-heading-rank; base=h2

\subsubsection{論平信徒不可執行任何司祭職務：他既不可施洗，不可獻祭，不可按手，也不可祝福。}

X. 我們也不准許平信徒執行任何屬於司祭職的職務；例如，不可獻祭，不可施洗，不可按手，不可祝福，無論是小祝福或大祝福。因為\BibleQuote{沒有人自己取得這尊位，而是蒙天主召叫的。}\BibleRef{希 5:4} 因為這些神聖的職務是由主教的按手授予的。但若有人未受託這樣的職務，卻擅自奪取，他將受烏齊雅的懲罰。\par

% structure: p0024 source=h3 target=subsubsection reason=relative-heading-rank; base=h2

\subsubsection{只有主教與長老，甚至下級神職人員中的任何人，都不許執行司祭的職務；祝聖權完全屬於主教，不屬於任何人。}

XI. 不但如此，我們也不准許其餘的神職人員施洗——例如，讀經者、詠經者、門衛、輔禮員——而只准許主教和長老，但執事要在其中輔助他們。凡膽敢這樣做的，將受科辣黑同夥的懲罰。\EditorialNote{戶籍紀第十六章} 我們不許長老祝聖執事、或女執事、或讀經者、或輔禮員、或詠經者、或門衛，只許主教這樣做；因為這是教會的秩序與和諧。\par

% structure: p0026 source=h3 target=subsubsection reason=relative-heading-rank; base=h2

\subsubsection{摒棄一切不愛德的行為。}

XII. 至於嫉妒、忌恨、誹謗、爭吵、好爭論，我們已經對你們說過，這些與基督徒格格不入，尤其在寡婦身上。但因為在人身內工作的魔鬼，其行為狡猾而充滿各種詭計，他走向那些並非真正寡婦的人，如同從前走向加音（因為有些人自稱是寡婦，卻不履行合於寡婦身份的誡命；如同加音未盡兄弟當盡之責：因為他們不考慮，並非寡婦之名將他們帶入天主的國，而是真正的信德與聖善的行為）。但若有人擁有寡婦之名，卻做敵對者的行為，她的寡婦身份將不被記算，她將被逐出天國，交付於永罰。因為我們聽說有些寡婦是嫉妒的、好忌恨的誹謗者，嫉妒別人的安寧。這樣的寡婦不是基督的門徒，也不是祂教義的門徒；因為當她們中有一位寡婦被某人穿衣，或收到金錢、或食物、或飲料、或鞋時，她們看見姐妹的安慰，應當說：——\par

% structure: p0028 source=h3 target=subsubsection reason=relative-heading-rank; base=h2

\subsubsection{寡婦應如何為供應她們需要的人祈禱。}

XIII. 天主，你是可稱頌的，因為你安慰了我的同伴寡婦。主啊，求你祝福並光榮那賜予她這些東西的人，願他的善行真實地上升到你面前，並在他的眷顧之日，求你記念他，賜他恩惠。至於我的主教——他如此妥善地履行了你託付的職責，並安排了一次合時的施捨給我的同伴，那位赤身的寡婦——求你增加他的光榮，並在揭示你眷顧之日賜他喜樂的冠冕。同樣，那接受施捨的寡婦應與另一位寡婦一同為那供應她的人祈禱。\par

% structure: p0030 source=h3 target=subsubsection reason=relative-heading-rank; base=h2

\subsubsection{那善待窮人者，不應喧嚷、宣揚自己的名字，應按主的憲章而行。}

十四、但若有任何妇女行善，她应如明智之人，隐藏自己的名字，不在前面吹号，好使她的施舍在天主面前是隐藏的，正如主所说：\Quote{你施舍时，不要让左手知道右手所行的，好使你的施舍在隐密中。} \BibleRef{玛 6:3} 并让寡妇为那给她施舍的人祈祷，不论他是谁，因为她如同基督的圣祭台；而\Quote{在隐密中看见的父，必要公开报答那行善的人。} 但那些不愿按天主的命令生活的寡妇，却忧虑好奇，打听是哪位女执事施舍，以及哪些寡妇接受了施舍。当她得知这些事后，她便抱怨分施慈善的女执事，说：难道你没有看见我更加困苦，更需要你的慈善吗？你为什么优先于她呢？她说这些话是愚昧的，不明白这事不在于人的意志，而在于天主的安排。因为如果她本人见证自己更亲近，经过查询，确实比别人更加贫乏和赤身，她就该明白是谁制定了这规矩，并保持沉默，不要抱怨分施慈善的女执事，而要进入自己的房屋，俯伏在地，向天主祈求，好使她的罪得赦免。因为天主吩咐携带慈善的女执事不可宣扬此事，而这个寡妇却因她没有公布她的名字以致她能知道并跑去接受而抱怨；不仅如此，她还咒骂她，忘记了那说过的话：\Quote{祝福你的，必受祝福；咒骂你的，必受咒骂。} \BibleRef{创 27:29} 但主说：\Quote{你们进一家时，要说：愿这家平安。那里如有平安之子，你们的平安就必临到他；若不然，你们的平安仍归于你们。} \BibleRef{路 10:5}\par

% structure: p0032 source=h3 target=subsubsection reason=relative-heading-rank; base=h2

\subsubsection{我们不应辱骂邻人，因为咒骂有违基督徒身份。}

十五、因此，如果平安回到那些派遣它的人，甚至回到那些先前实际给予它的人，因为它没有找到合适接受的人，那么咒骂更会不义地回到那发出它的人的头上，因为那接受它的人不配。因为凡无端辱骂他人者，实是咒骂自己，如撒罗满所说：\Quote{如同雀鸟飞来飞去，无故的咒骂也不会落在任何人身上。} \BibleRef{箴 26:2} 他又说：\Quote{出言毁谤的，极其愚昧。} \BibleRef{箴 10:18} 但如同蜜蜂，一种力量微弱的受造物，若螫了人，就失去毒刺，变成雄蜂；同样，你们对别人所行的任何不义，也必临到你们自己身上。\Quote{他掘了坑，挖了井，必要掉在自己所挖的坑里。} 又说：\Quote{为邻人挖坑的，必掉进坑里。} \BibleRef{箴 26:27} 因此，凡避免咒骂的，就不要咒骂别人；因为\Quote{你所憎恶的事，不要对别人做。} \BibleRef{多 4:16} 所以，要劝慰那些软弱的寡妇，坚固那些软弱的，并赞扬那些在圣洁中行走的。愿她们祝福，而非毁谤；愿她们缔造和平，而非挑起争端。\par

% structure: p0034 source=h2 target=subsection reason=relative-heading-rank; base=h2

\subsection{第二节  论执事与女执事、其他圣职人员及圣洗圣事}

因此，无论是主教、\OriginalTerm{司铎}{presbyter}、执事，或任何其他司祭名册上的人，都不要以毁谤玷污自己的舌头，免得继承咒骂而非祝福；而主教也有责任和关心，确保没有平信徒发出任何咒骂：因为他应当关心所有的人——圣职人员、贞女、寡妇、平信徒。为此，主教啊，要任命你的同工，那些为生命和正义劳苦的人，选出令天主喜悦的执事，那些你在全体百姓中所证明是配得的，并随时准备好服务所需的人。也要任命一位忠实而圣洁的\OriginalTerm{女执事}{deaconess}，为了对妇女的服务。因为有时他不能派遣一位男执事到妇女那里去，因着不信者的缘故。因此，你要派遣一位妇女，即女执事，因着恶人的幻想。因为我们在许多事上都需要一位妇女，即女执事；首先在妇女领受圣洗时，执事只以圣油傅抹她们的前额，之后女执事傅抹她们：因为妇女们没有必要被男人们看见；只有在覆手时，主教傅抹她的头，如同从前傅抹司祭和君王一样，不是因为现在受洗的人被祝圣为司祭，而是作为基督徒，即受傅者，源于基督——受傅者，\BibleQuote{是君王司祭、圣洁的邦国、天主的教会、婚姻之殿的柱石和根基} \BibleRef{伯前 2:9}，从前不是百姓，如今却是蒙爱和蒙选的，他们身上呼求了祂的新名，正如依撒意亚先知所见证的，说：\Quote{他们要用主为他们命名的新名称呼这百姓。}\par

% structure: p0036 source=h3 target=subsubsection reason=relative-heading-rank; base=h2

\subsubsection{论圣洗圣事的入门礼。}

十六、因此，主教啊，你应按此样式，以圣油傅抹那些即将受洗者的头，无论男女，作为属灵洗礼的预像。之后，或是你——主教，或是你下属的\OriginalTerm{司铎}{presbyter}，要以庄严的格式呼父、子及圣神之名，将他们浸入水中；并让一位执事接住男人，一位\OriginalTerm{女执事}{deaconess}接住女人，以使这不可侵犯的印记的赋予得以合宜地进行。之后，让主教以香膏傅抹那些受洗者。\par

% structure: p0038 source=h3 target=subsubsection reason=relative-heading-rank; base=h2

\subsubsection{论受洗归于基督的意义，以及其中所说所行的一切缘由。}

XVII. 此圣洗因此是归于耶稣之死：水代表埋葬，油代表圣神；印记代表十字架；傅油是宣认的确认；提到圣父，作为创始者和差遣者；同时提到圣神，作为见证；下入水中是与基督同死；从水中上来是与祂同复活。圣父是万有的天主；基督是独生子天主、爱子、荣耀之主；圣神是护慰者，由基督所派遣，受祂教导，并宣讲祂。\par

% structure: p0040 source=h3 target=subsubsection reason=relative-heading-rank; base=h2

\subsubsection{领受圣洗者应具备何种品格}

XVIII. 但领受圣洗者应远离一切不义；是已停止作罪的人，天主的朋友，魔鬼的仇敌，天主圣父的嗣子，祂圣子的同嗣子；是已弃绝撒殚、魔鬼及撒殚的诡计的人；应当是贞洁、纯洁、神圣、蒙天主所爱、天主的子女，如儿子向父亲祈祷，并从全体信众的共同集会说：\Quote{我们在天的父！愿你的名被尊为圣；愿你的国来临；愿你的旨意承行于地，如在天上一样；求你今天赏给我们日用的食粮；求你宽免我们的债，如同我们也宽免欠我们债的人；不要让我们陷入诱惑，但救我们免于凶恶：因为国度、权能、荣耀，都是你的，直到永远。阿们。}\par

% structure: p0042 source=h3 target=subsubsection reason=relative-heading-rank; base=h2

\subsubsection{执事应具备何种品格}

XIX. 执事应在一切事上无玷污，如同主教本人一样，只是更为活跃；人数根据教会的大小，以便作为不羞愧的工人服事软弱者。女执事应殷勤照顾妇女；但两者都应随时准备传递信息、旅行、服事和伺候，正如依撒意亚论到主所说：\Quote{使义人得以称义，他忠信地服事许多人。}因此，每个人应知道自己的适当位置，并一心一意、同心合意地勤勉履行，知道他们职务的报酬；但不应羞于服事那些有需要的人，因为我们的\Quote{主耶稣基督来不是为受服事，而是为服事，并交出祂的生命为大众作赎价。}\BibleRef{玛 20:28}同样，他们也应当这样做，如果不得不为弟兄舍命，也不应犹豫。因为主和救主耶稣基督并不犹豫\Quote{交出祂的生命}，正如祂自己所说，\Quote{为朋友。}\BibleRef{若 15:13}因此，如果天地的主宰为我们承受了所有苦难，那么，你这位应当效法祂为我们忍受了奴役、贫乏、鞭打和十字架的人，怎能为服事有需要的人而为难呢？我们也应当效法基督，服事弟兄。因为祂说：\Quote{你们中谁愿为大，就当作你们的仆役；谁愿为首，就当作你们的奴仆。}\BibleRef{玛 20:26}因为祂确实，而不只是言语，成就了\Quote{忠信地服事许多人}的预言。\BibleRef{依 53:11}因为\Quote{祂拿了一条毛巾束腰，随后把水倒在盆里，当我们坐席时，祂来洗了我们众人的脚，并用毛巾擦干。}\BibleRef{若 13:4}祂以此行动向我们显示了祂的仁慈和手足之爱，使我们也能彼此同样去做。因此，如果我们的主和师傅如此谦卑自己，你们这些真理的工人和虔诚的管理者，又怎能羞于对弟兄中那些软弱有病的人做同样的事呢？因此，要以善良的心服事，不要抱怨或反抗；因为你们不是为人做的，而是为天主做的，并将在你们被眷顾的日子从祂那里得到你们职务的报酬。你们执事有责任探望所有需要探望的人。并将所有在困苦中的人告诉你们的主教；因为你们应当像他的灵魂和感官——在一切事上对他积极而留心，如同对你们的主教、父亲和师傅。\par

% structure: p0044 source=h3 target=subsubsection reason=relative-heading-rank; base=h2

\subsubsection{主教应由三位或两位主教祝圣，而非一位；否则无效}

XX. 我们命令，主教应由三位主教祝圣，或至少两位；但不可由一位主教任命你们的主教；因为两三个见证人的证词更为坚固可靠。但长老和执事应由一位主教及其余圣职人员祝圣。无论是长老还是执事，都不得将平信徒擢升为圣职人员；长老只可教导、奉献祭品、施行圣洗、祝福百姓，执事则服事主教和长老，即履行辅助执事的职务，不得干预其他职务。\par

\SourceRecord{Translated by James Donaldson.；Ante-Nicene Fathers；Vol. 7.；Edited by Alexander Roberts, James Donaldson, and A. Cleveland Coxe.；Buffalo, NY: Christian Literature Publishing Co.,；1886.；<http://www.newadvent.org/fathers/07153.htm>.}

% structure: p0001 source=h1 target=section reason=page-title

\section{宗徒宪章（卷四）}

% structure: p0002 source=h2 target=subsection reason=relative-heading-rank; base=h2

\subsection{第一节  论赒济穷人}

% structure: p0003 source=h3 target=subsubsection reason=relative-heading-rank; base=h2

\subsubsection{无子女者应收养孤儿，待之如己出。}

一、当任何基督徒成为孤儿，无论是青年或少女，最好由一位无子女的弟兄收养该青年，视之如子；若有儿子年龄相仿且已届婚龄，则应娶该少女为妻。行此事者，成就了一桩大功，成为孤儿的父亲，必要从主天主领受这爱德的赏报。但若有人为取悦于人而富有，因而羞于接纳孤儿，那么孤儿的父亲和寡妇的裁判者必会照顾孤儿，而那人自己却要得到一个败家的继承人；此事必如经上所说：\Quote{圣民未曾吃过的，亚述人必要吃尽。}依撒意亚也说：\BibleQuote{在你们眼前，外邦人吞吃你们的土地。}\BibleRef{依 1:7}\par

% structure: p0005 source=h3 target=subsubsection reason=relative-heading-rank; base=h2

\subsubsection{主教应如何照顾孤儿。}

二、所以，你们这些主教，应当关心他们的生计，一无所缺；对孤儿要显出父母的关怀，对寡妇显出丈夫的关怀，对适龄者予以婚配，对工匠给予工作，对无能者给予怜悯，对旅客供给房舍，对饥饿者供给食物，对口渴者供给饮料，对裸者供给衣服，对病者予以探访，对被囚者予以援助。除此以外，更当加倍关心孤儿，使他们一无缺欠；至于少女，直到她达婚龄，你们将她嫁给一位弟兄；对青年则予以帮助，使他学一门手艺，并能从中获利维持生计；这样，当他熟练掌握这门手艺时，便能为自己购买工具，不再拖累任何弟兄或他们真诚的爱，而能自食其力。因为那能自食其力、不占取孤儿、旅客和寡妇之位的人，确实是有福的。\par

% structure: p0007 source=h3 target=subsubsection reason=relative-heading-rank; base=h2

\subsubsection{按主的定规，谁应受供养。}

三、诚如主曾说：\BibleQuote{施予比领受更为有福。}\BibleRef{宗 20:35}因为祂又说过：\Quote{凡有所有却假冒伪善地领受，或能自养却仍从他人领受的人有祸了；因为他们两人都要在审判之日向主天主交账。}但一个因年幼、或因年老体弱、或因疾病缠身、或因养育众多子女而接受施舍的孤儿，非但不受责备，反受称赞；他必被视为天主的祭台，得蒙天主的尊荣，因他热切而恒常地为施舍者祈祷；他不是白白领受，而是尽己所能以祈祷回报所领受的。这样的人必蒙天主祝福，得享永生。但那有所有却假冒伪善或懒惰地领受，而不去工作并帮助他人的人，在天主面前必受惩罚，因为他抢夺了穷人的食粮。\par

% structure: p0009 source=h3 target=subsubsection reason=relative-heading-rank; base=h2

\subsubsection{论贪爱金钱。}

四、因为那有钱财而不施与他人，也不为自己使用的人，就像一条蛇，据说它守护着宝藏而沉睡；关于他，经上的这句话是真实的：\Quote{他积聚财富，却不得尝用}；当他合理地灭亡时，那些财富对他毫无用处。因为经上说：\Quote{在震怒之日，财富毫无益处。}这样的人不信天主，只信自己的黄金；他以此为神，并信赖其中。此人是真理的伪装者，看人情面，不忠信，欺诈，胆怯，无丈夫气，轻浮，无价值，好抱怨，常痛苦，是自己的仇敌，也非任何人的朋友。他的钱财必要丧失，一个外人必要消耗它，或是他生前被盗，或是死后被继承。\Quote{因为不义之财必被吐出。}\par

% structure: p0011 source=h3 target=subsubsection reason=relative-heading-rank; base=h2

\subsubsection{人应怀着怎样的敬畏领受主的奉献物。}

五、因此，我们劝勉寡妇和孤儿，当怀着完全的敬畏和虔诚，领受施予他们的物品，并感谢那赐食物给穷人的天主，举目仰望祂。因为祂说：\Quote{你们中间谁不吃不喝能离得开祂？因为祂张开手，使一切生物饱享祂的慈恩；赐谷物给青年，赐酒给少女，赐油给活人的喜乐，赐草给牲畜，赐青蔬给人享用，赐肉给野兽，赐种子给飞鸟，赐合宜的食物给一切受造物。}因此主说：\BibleQuote{你们看那天上的飞鸟，不播种，也不收割，也不积蓄在仓里，你们的天父尚且养活它们。你们不比它们更贵重吗？所以不要忧虑说：我们吃什么？喝什么？因为你们的天父知道你们需要这一切。}既然你们享受祂如此眷顾，并分享从祂而来的美物，就应当赞美那接纳孤儿和寡妇的天主，全能的天主，借着祂的爱子、我们的主耶稣基督；愿荣耀藉着祂在精神和真理中归于天主，直到永远。阿们。\par

% structure: p0013 source=h3 target=subsubsection reason=relative-heading-rank; base=h2

\subsubsection{谁的奉献物应被接纳，谁的奉献物不应被接纳。}

六、现在主教应当知道，该接受谁的祭品，不该接受谁的。因为他要躲避奸商，不接受他们的礼物。因为\Quote{奸商必不能因罪而成义。}为此依撒意亚斥责以色列说：\Quote{你们的奸商用酒搀水。}他也要躲避犯奸淫者，因为\Quote{你不可将娼妓的酬金奉献给主。}\EditorialNote{申 23:18}他也要躲避勒索者、贪图他人财物者、通奸者，因为这样的人的\OriginalTerm{祭献}{sacrifices}在天主面前是可憎的。还有那些欺压寡妇、凌辱孤儿、以无辜者充满监狱、恶待自己仆人——我是指用鞭打、饥饿、苦役——甚至毁坏整座城市的人。主教啊，要躲避这样的人及他们那可憎的祭品。你也要拒绝恶棍、不义诉讼的讼师、造偶像者、窃贼、不义的税吏、用假秤和诈量欺哄人的、作伪证且不满足于薪饷、反而向穷人施暴的兵士、杀人者、刽子手、不义的法官、颠倒是非者、埋伏害人者、行可憎恶之事者、酒徒、亵渎者、男色者、放高利贷者，以及一切邪恶、违背天主旨意的人。因为圣经说，所有这样的人在天主面前都是可憎的。因为从这样的人那里接受礼物，并以此供养寡妇和孤儿的人，将在天主的审判台前受罚；正如\WorkTitle{列王纪}中\Person{阿多尼雅}先知，因\Person{雅洛贝罕}的邪恶而违背天主，\Quote{在主管禁止他的地方吃饼喝水}\EditorialNote{列上 13}就被狮子咬死。因为从劳苦中分给寡妇的饼，虽然少而小，也比从\OriginalTerm{不义}{injustice}和诬告中得来的，虽多而精更好。因为圣经说：\Quote{义人拥有的虽少，却胜过恶人的大量财富。}现在，即使有一个寡妇吃着从恶人那里得来的食物而饱足，为他们祈祷，她也不会被俯听。因为洞察人心的天主，已对恶人宣判说：\BibleQuote{纵然\Person{梅瑟}和\Person{撒慕尔}站在我面前为他们求情，我也不会俯听他们；}\BibleRef{耶 15:1}又说：\BibleQuote{你不要为这人民祈祷，不要为他们求怜悯，也不要为他们向我恳求，因为我不会俯听你。}\BibleRef{耶 7:16}\par

% structure: p0015 source=h3 target=subsubsection reason=relative-heading-rank; base=h2

\subsubsection{不配者的祭品，当他们如此时，不仅不能平息天主的愤怒，反而激起祂的义怒。}

七、并且不只是这些，那些在罪中而未悔改的人，他们祈祷时不仅不被俯听，反而会激起天主的怒气，因为使他们想起自己的邪恶。因此要避免这样的服务，如同避讳狗的价格和娼妓的酬金，因为两者都被法律所禁止。因为\Person{厄里叟}没有接受\Person{哈匝耳}带来的礼物，\Person{阿希雅}也没有接受\Person{雅洛贝罕}的礼物；\EditorialNote{列上 14}但如果天主的先知都不接受恶人的礼物，那么主教们，你们也不应该接受，这是合理的。不仅如此，当\Person{西满}术士向我和\Person{伯多禄}、\Person{若望}献上金钱，企图用钱购买那无法估量的\OriginalTerm{恩宠}{grace}时，我们拒绝了他，并以永远的诅咒束缚他，因为他以为天主的恩赐不是以对天主的虔诚之心，而是以金钱的代价来拥有。\EditorialNote{宗 8}因此，要避免那些不是出于良好良心的祭品献到天主的祭台上。因为他说：\BibleQuote{你要远离一切不义，你就不会惧怕，战栗也不会临近你。}\BibleRef{依 54:14}\par

% structure: p0017 source=h3 target=subsubsection reason=relative-heading-rank; base=h2

\subsubsection{用自己的劳力供给寡妇，即使少而微，也胜过从恶人那里得来的多而大；因为宁可饿死，也不接受恶人的祭品。}

八、但如果你说那些施舍的人是这样的人，如果我们不接受他们的，我们拿什么来管理寡妇？人民中的穷人靠什么维持？你要听我们说：正因此你们接受了肋未人的恩赐，即你们人民的祭品，使你们自己和那些有需要的人得以足够；使你们不至于如此窘迫以致要从恶人那里接受。但如果教会如此窘迫，宁可饿死也不要从天主的敌人那里接受任何东西，以免给天主的朋友带来羞辱和谩骂。关于这样的人，先知说：\Quote{不要让罪人的油湿润我的头。}因此要审查这样的人，只从那些行走在圣洁中的人接受，供应受苦的人。但不要从被逐出教会者那里接受，直到他们被认为配得成为教会的成员。但如果缺少礼物，就通知弟兄们，从他们那里收集，然后以此\OriginalTerm{本着正义}{in righteousness}照顾孤儿寡妇。\par

% structure: p0019 source=h3 target=subsubsection reason=relative-heading-rank; base=h2

\subsubsection{司铎应劝勉人民向穷人行善，正如智慧者\Person{撒罗满}所说。}

九、要对你管辖的人民说智慧者\Person{撒罗满}所说的：\Quote{要以你的正义劳动光荣主，从你的正义果实中向他献上\OriginalTerm{初果}{first-fruits}，使你的仓廪充满五谷，你的榨酒池溢出新酒。}因此要用信友正义劳动所得来供养和穿衣那些有需要的人。并且要按照上述方式从他们那里收集来的款项，指定用于赎回圣人、拯救奴隶、俘虏、囚犯、受虐待者，以及因基督的名被暴君判处决斗和死亡的人。因为圣经说：\Quote{要解救被带去死的人，赎回准备被杀的人，不要吝惜。}\par

% structure: p0021 source=h3 target=subsubsection reason=relative-heading-rank; base=h2

\subsubsection{一条宪令：若不敬者强行把钱掷给司铎，他们应将之用于木柴和煤炭，而不可用于食物。}

X．若你有时被迫从任何不虔敬者手中接受金钱，就当将这笔钱用于购买木柴与煤炭，以免寡妇或孤儿收受分毫，或被迫以此购买不宜食用的肉或饮品。因为不虔敬者的馈赠理应作为火的燃料，而非虔敬者的食物。律法明确规定了这一方法，它称存放过久的祭献不宜食用，并命令以火焚毁。这样的供献本质上并非邪恶，而是因奉献者的心意而成为不洁。我们如此规定，是为了不拒绝那些来到我们面前的人，因为我们知道，虔敬者日常的交往常常对不虔敬者大有裨益，但与他们的宗教共融则唯独有害。亲爱的诸位，为你们的益处，话已说够。\par

% structure: p0023 source=h2 target=subsection reason=relative-heading-rank; base=h2

\subsection{第二节 论家庭与社交生活}

% structure: p0024 source=h3 target=subsubsection reason=relative-heading-rank; base=h2

\subsubsection{论父母与子女}

XI. 你们作父亲的，要在主内教育你们的子女，用主的管教和劝诫养育他们；教导他们合乎圣言的正当手艺，免得他们因这些机会而放纵无度，且因未受父母的惩罚而提前松懈，偏离良善。因此，不要害怕责备他们，以严厉管教教训他们智慧。因为你们的纠正不会杀死他们，反而会保全他们。正如撒罗满在\WorkTitle{智慧篇}某处所说：\Quote{管教你儿子，他必使你安息；这样你对他必大有指望。你实在要用杖击打他，救他的灵魂免下死亡。} 同一位撒罗满又说：\Quote{不忍用杖打儿子的，是恨恶他；} 随后又说：\Quote{趁他年幼时击打他的肋旁，免得他变硬而违抗你。} 因此，懈怠于劝诫和教导自己儿子的，是恨恶自己的儿子。所以你们要教导子女上主的话语。用严厉的鞭打降服他们，使他们自幼顺从，教导他们圣经——基督信仰与神圣的经文，并将一切圣书传授给他们，\Quote{不给他们放纵的自由而使他们反客为主，} 违背你们的意愿。不准他们与同伴聚饮。因为如此他们会转向放荡的行径，陷于淫乱。倘若因父母的疏忽而致此，生养他们的人将为他们灵魂负责。因为如果失足的子女因父母的懈怠而与放荡之人结交，他们不仅自己受罚，其父母也将因他们而被定罪。为此，当子女达到适婚年龄时，务必尽力为他们成婚并安置妥当，免得在年少热血之时生活放荡，致使你们在审判之日需向上主天主交账。\par

% structure: p0026 source=h3 target=subsubsection reason=relative-heading-rank; base=h2

\subsubsection{论仆人与主人}

XII. 至于仆人，我们还能说什么呢？无非是：奴仆要以敬畏天主的心善待主人，即便主人不虔不义，但在敬拜上绝不可丝毫顺从。主人应爱其仆人，虽然他是他的下属。他要想到他们是平等的，因为同是凡人。若有信徒作主人，仆人应爱他如主、如同道、如父亲，但仍需保存主人权威：\Quote{不是在眼前事奉，而是真心爱主，因知道天主必为他的服从报偿他。} 同样，若有信徒作仆人，主人应爱他如子、如兄弟，因信仰相通，但仍须保持主仆之别。\par

% structure: p0028 source=h3 target=subsubsection reason=relative-heading-rank; base=h2

\subsubsection{论在何事上应服从现世统治者}

XIII. 在一切悦乐天主的事上，要服从一切王权和统治权，视他们为天主的仆役和惩罚不虔敬者的人。要向他们献上应得的敬畏、供物、一切税赋、尊敬、礼物和捐税。因为这是天主的命令：你们不要欠任何人什么，只欠彼此相爱的保证，这爱是天主藉基督所命令的。\par

% structure: p0030 source=h3 target=subsubsection reason=relative-heading-rank; base=h2

\subsubsection{论童贞}

XIV. 关于童贞，我们没有领受诫命，而是将之留给愿意者的权力，作为一项誓愿；我们劝勉他们在这事上不可轻率许诺，因为撒罗满说：\Quote{你许愿不还，不如不许。} \BibleRef{训 5:5} 因此，这样的童贞女应在身体和灵魂上都成为圣洁，作为天主的圣殿，\BibleRef{格前 7:34} 作为基督的居所，作为圣神的宫殿。因为许愿者应当行与誓愿相称的事，表明她的誓愿是真实的，且是为虔敬而得的闲暇，而非对婚姻的羞辱。她不可好出游荡，不可不合时宜地闲逛；不可三心二意，而是庄重、有节制、清醒、纯洁，避免与多人来往，尤其要避免与声名狼藉的人交往。\par

\SourceRecord{Translated by James Donaldson.；Ante-Nicene Fathers；Vol. 7.；Edited by Alexander Roberts, James Donaldson, and A. Cleveland Coxe.；Buffalo, NY: Christian Literature Publishing Co.,；1886.；<http://www.newadvent.org/fathers/07154.htm>.}

% structure: p0001 source=h1 target=section reason=page-title

\section{\WorkTitle{Apostolic Constitutions (Book V)}}

% structure: p0002 source=h2 target=subsection reason=relative-heading-rank; base=h2

\subsection{第一节　论殉道者}

% structure: p0003 source=h3 target=subsubsection reason=relative-heading-rank; base=h2

\subsubsection{一、按主的训令，信友理应供应因基督而为不信者所迫害之人的需用}

一、若有基督徒因基督之名及对天主的爱与信德，被不敬虔的人判决到竞技场、野兽或矿坑，不可疏忽他；当以你的劳作和血汗所得，供给他的生计，并酬谢兵士，使他得蒙缓解与照顾；如此，在你们能力所及之内，你们有福的弟兄不致受苦。因为为主天主之名受判决的人，是圣洁的殉道者，是主的弟兄，至高者的儿子，圣神的器皿，每位信友借着他，获得了圣福音荣耀的光照，蒙赐不朽的冠冕，并参与基督的苦难与祂的血的共融，为成为子女而肖似基督的死。为此，你们所有信友，应通过你们的主教，从你们的财物和劳作中供给圣人们。若有人没有财物，可以禁食一日，将那日省下的归出，为圣人们安排。若人有富余，应按照自己的能力多行供给。若有人能变卖所有生计，将那些人从监狱中赎回，他是有福的，是基督的朋友。因为若有人将财物施舍穷人便为成全——假定他拥有天主的知识——那为殉道者而行此事的人更是如此。这样的人堪当天主，他供给那些在万民、君王及以色列子孙面前承认主的人，以此承行天主的旨意。关于这些人，我们的主曾宣告说：\BibleQuote{凡在人面前承认我的，我在我父面前也必承认他。}\BibleRef{玛 10:32} 这些人既蒙基督在父面前作证，你们就不该羞于到监狱中探望他们。因为你们这样做，将被视为你们的见证；实际的考验是他们的见证，而你们的乐意也将成为你们的见证，如同参与了他们的战斗。因为主在某处对这样的人说：\BibleQuote{你们蒙我父降福的，来承受自创世以来为你们预备的国度吧。因为我饿了，你们给了我吃的；我渴了，你们给了我喝的；我作客，你们收留了我；我赤身露体，你们给了我穿的；我患病，你们看顾了我；我在监里，你们来探望了我。那时义人要回答说：主啊，我们什么时候见过祢饿了，给了祢吃的？或渴了，给了祢喝的？什么时候见过祢赤身露体，给了祢穿的？或患病，或在监里，来探望过祢？祂要回答他们说：我实在告诉你们，你们为我这些最小弟兄中的一个所做的，就是为我做的。这些人要进入永生。然后祂要对左边的人说：离开我，受咒骂的，进入那为魔鬼及其使者预备的永火里吧。因为我饿了，你们没有给我吃的；我渴了，你们没有给我喝的；我作客，你们没有收留我；我赤身露体，你们没有给我穿的；我患病，在监里，你们没有看顾我。那时他们也要回答说：主啊，我们什么时候见过祢饿了，或渴了，或作客，或赤身露体，或患病，或在监里，而没有服侍祢？那时祂要回答他们说：我实在告诉你们，你们没有为我这些最小中的一个所做的，就是没有为我做。这些人要进入永罚里去。}\par

% structure: p0005 source=h3 target=subsubsection reason=relative-heading-rank; base=h2

\subsubsection{二、当假弟兄仍行恶时，我们应避免与他们交往}

二、若有人自称弟兄，却被恶者诱惑，行邪恶之事，并被定罪判处死刑——如为奸淫者或杀人者——要离开他，使你们安全，免得你们中有人被怀疑与如此可憎的行为有份；也免得有恶名传出，仿佛所有基督徒都喜爱不法之事。所以，要远离他们。但那些因基督被不敬虔之人凌辱、关押在狱中，或被交付死亡、锁链或流放的人，你们当尽一切努力援助他们，为从恶人手中解救你们的肢体。若有人与这些人交往而被捕，陷于灾祸，他是有福的，因为他与殉道者同份，是效法基督苦难的人。因为我们自己也曾多次因基督的缘故，受盖法、亚历山大和亚纳斯的鞭打，却\BibleQuote{欢喜地离开，因为堪当为救主受这样的苦。}你们在受这样的苦时也要欢喜，因为那一天你们是有福的。\par

% structure: p0007 source=h3 target=subsubsection reason=relative-heading-rank; base=h2

\subsubsection{三、即使自己陷于危险，也当伸手援助为基督而被剥夺一切的人}

三、也要接待那些因信德受迫害、并因主的诫命\BibleQuote{从一城逃往另一城}的人，\BibleRef{玛 10:23} 要协助他们如同殉道者，因你们有份于他们的迫害而喜乐，因为知道主称他们为有福的。主亲自说：\BibleQuote{几时人为了我的缘故辱骂你们，迫害你们，捏造一切坏话毁谤你们，你们是有福的。你们欢喜踊跃吧！因为你们在天上的赏报是丰厚的，因为你们以前的先知也这样迫害过。}\BibleRef{玛 5:11} 又说：\BibleQuote{他们若迫害了我，也要迫害你们。}\BibleRef{若 15:20} 随后又说：\BibleQuote{若有人在这城迫害你们，就逃往另一城去。因为在世界上你们有苦难；他们要把你们交到会堂，并要为我的缘故被带到官长和君王面前，对他们作证。}又说：\BibleQuote{那坚持到底的，才可得救。}\BibleRef{玛 10:22} 因为为信德受迫害，并为基督作证而忍耐的人，真是天主的人。\par

% structure: p0009 source=h3 target=subsubsection reason=relative-heading-rank; base=h2

\subsubsection{否认基督乃可怖且毁灭之事}

IV. 但凡有人否认自己是基督徒，以免被人憎恨，如此爱自己的生命胜过爱主（生命气息在其手中）的，是可悲而可怜的，因他成为可憎可厌的，渴望作人的朋友，却是天主的仇敌，不再与圣徒同分，而与受咒诅者同列；他选择的是永恒之火，那为魔鬼及其使者预备的，而非蒙福的国度：他不再被人憎恨，却被天主弃绝，逐出祂的临在。关于这样的人，我们的主曾宣告说：\BibleQuote{凡在人前否认我，并以我的名为耻的，我在天之父面前也必否认他，并以他为耻。}\BibleRef{玛 10:33} 祂又亲自向我们，祂的门徒，这样说：\BibleQuote{凡爱父亲或母亲胜过我的，不配是我的；凡爱儿子或女儿胜过我的，不配是我的；凡不背起自己的十字架跟随我的，不配是我的。凡寻得生命的，必要丧失生命；凡为我的缘故丧失生命的，必要寻得生命。人若赚得全世界，却赔上自己的灵魂，有什么益处呢？或者人能拿什么换回自己的灵魂呢？} 随后又说：\BibleQuote{不要害怕那些杀害身体，却不能杀害灵魂的；但要害怕那能在地狱里毁灭灵魂和身体的那一位。}\BibleRef{玛 10:28}\par

% structure: p0011 source=h3 target=subsubsection reason=relative-heading-rank; base=h2

\subsubsection{论我们应在受苦中效法基督，并热切追随祂的忍耐}

V. 凡学习任何技艺的人，当他看见他的师傅以勤勉和技巧完善他的技艺时，自己也竭力使手中的工作与之相似。若他不能，他的工作便不能完善。我们既有师傅，即我们的主耶稣基督，为何不跟随祂的教导？——因为祂为着对圣父的虔诚和对我们人类族群的爱，舍弃了安息、享乐、荣耀、财富、骄傲、报复之权、祂的母亲和兄弟，甚至祂自己的生命；祂不仅遭受了迫害和鞭打、凌辱和嘲笑，而且遭受了十字架之苦，为要拯救悔改的犹太人和外邦人。因此，若祂为我们舍弃了安息，不以十字架为耻，也不视死亡为可耻，我们为何不以祂所赐的忍耐，效法祂的苦难，并为了祂而舍弃我们自己的生命呢？因为祂所做的一切都是为了我们，而我们做这一切是为了我们自己：祂并不需要我们，但我们却需要祂的怜悯。祂只要求我们信德的真诚和准备，正如圣经所说：\BibleQuote{你若行义，你能给祂什么？或祂从你手中能接受什么？你的邪恶只害你这样的人，你的公义只益人子。}\par

% structure: p0013 source=h3 target=subsubsection reason=relative-heading-rank; base=h2

\subsubsection{论信友既不应因安稳而轻率冒险，也不应因懦弱而过度畏怯，却当因惧怕而逃遁；但若落入敌手，当为那为他存留的冠冕而竭力奋斗}

六、因此，当父母、亲属、朋友、妻子、儿女、财产及一切生活的享乐成为虔敬的阻碍时，让我们舍弃这一切。因为我们应当祈祷，免得陷于诱惑；但若我们被召叫去殉道，要恒心宣认祂宝贵的圣名，并因此受罚，让我们喜乐，如同奔向永生。当我们受迫害时，不要觉得奇怪；不要爱现世，也不要求人的称赞，也不要爱统治者的荣耀和尊荣，如同有些犹太人因害怕大司祭和其他统治者而惊奇于我们主的奇事，却不信祂：\Quote{因为他们喜爱人的称赞胜过天主的称赞。}\BibleRef{若 12:43}但现在，透过宣认美好的信仰，我们不仅拯救自己，也坚固那些新受光照的人，并增强慕道者的信德。但若我们稍减我们的宣认，因我们信念的软弱和害怕短暂的惩罚而否认虔敬，我们不仅剥夺自己的永远荣耀，也会成为他人丧亡的原因；并要受双重的惩罚，因为我们的否认使人怀疑我们先前如此夸耀的那真理是谬误的教训。因此，我们不可鲁莽仓促地投入危险，因为主说：\Quote{祈祷免得陷于诱惑：心神固然愿意，肉体却软弱。}也不可在陷入危险时，因我们的信仰而感到惧怕或羞愧。因为若有人因否认自己的望德——即耶稣，天主子——而免于一时的死亡，次日却病倒在床上，患上肠、胃或头部的不治之症，如肺痨、坏疽、痢疾、肠梗阻、水肿、腹绞痛等，突发灾祸而离世；他岂非失去了现世，也失去了永世？更确切地说，他已在永罚的边缘，\Quote{被丢在外面的黑暗里，那里要有哀号和切齿。}\BibleRef{玛 8:12}但那蒙受殉道荣耀的人，要因主而喜乐，因为由此获得如此大的冠冕，并借着宣认离开此世。即使他只是慕道者，也让他无扰地离去；因为为基督受苦对他将是更真实的圣洗，因为他真正与基督同死，而其他人只是预像。因此，让他效法他的师傅而喜乐，因为如此规定：\Quote{愿每人都完全，如同他的师傅一样。}\BibleRef{路 6:40}如今，祂和我们的师傅，主耶稣，为我们受苦：祂忍受辱骂和毁谤，恒久忍耐。祂被唾弃，脸被打，受拳击；被鞭打后，被钉在十字架上。祂喝醋和苦胆；当完成了所有经上所记载的事后，祂对祂的天主和圣父说：\Quote{父啊！我把我的灵魂交在祢手中。}\BibleRef{路 23:46}因此，凡愿做祂门徒的人，要热切跟随祂的战争：要效法祂的忍耐，知道即使被人用火烧，他也毫无损伤，如同那三少年；\EditorialNote{达尼尔书第三章}或者即使受苦，也必从主那里获得赏报，因为他信从唯一真实的天主和圣父，借着耶稣基督，伟大的大司祭、我们灵魂的救赎者，又是我等受苦的赏报者。愿光荣归于祂，直到永远。阿们。\par

% structure: p0015 source=h3 target=subsubsection reason=relative-heading-rank; base=h2

\subsubsection{关于复活、西彼拉以及斯多葛派关于凤凰鸟的几点论证}

VII. 因为全能的天主自己将藉我们的主耶稣基督，按照祂那不能落空的应许，使我们复活，并赐给我们与从起初以来所有长眠者一同复活；那时我们将如同现在我们的形体一样，没有任何缺陷或朽坏。因为我们将会不朽地复活：无论我们死于海上，或是散于地上，或被野兽和飞鸟撕裂，祂都会以祂的大能复活我们；因为整个世界都在天主手中。现在祂说：\BibleQuote{你们的头发一根也不会失落。} \BibleRef{路 21:18} 因此祂劝告我们说：\BibleQuote{你们要凭着坚忍，保全你们的灵魂。} \BibleRef{路 21:19} 但关于死人的复活和殉道者的报偿，加俾额尔向达尼尔说：\BibleQuote{许多长眠于尘土中的人要醒来，有的得到永生，有的得到永久的羞辱和永远的轻蔑。明悟的人要发光有如穹苍的光辉，有如星辰。} \BibleRef{达 12:2} 因此最神圣的加俾额尔预言圣人们要像星辰发光：因为祂的圣名给他们作证，使他们明白真理。复活不仅是为殉道者宣告的，也是为所有的人，义人和不义的人，虔敬和不虔敬的人，好使每人按自己的功过受报。因为圣经说，天主\BibleQuote{将审问一切所为的事，包括一切隐秘，无论是善是恶。} \BibleRef{训 12:14} 这复活从前犹太人并不相信，他们曾说：\Quote{我们的骨头枯干了，我们完了。} 上主回答他们说：\Quote{看，我要打开你们的坟墓，把你们从坟墓中领出来；我要把我的神放在你们内，你们就要活起来；你们要知道我上主说了，也必实行。} 祂又借依撒意亚说：\BibleQuote{死人要复活，墓中的人要起来。安息在尘土中的人要喜乐，因为从你来的甘露是医治他们的。} \EditorialNote{依 26:19} 关于复活，关于义人在荣耀中的存留，关于不虔敬者的惩罚、堕落、弃绝、定罪、羞辱，\BibleQuote{永火和不死的虫} \BibleRef{依 66:24}，确实有许多不同的说法。如果祂愿意所有人都能不死，祂在哈诺客和厄里亚的例子中显示了祂的能力，没有容许他们经历死亡。或者如果祂愿意在每一代中复活那些死去的人，这祂也能做到，祂曾借着厄里亚复活了寡妇的儿子\EditorialNote{列上 17}，通过厄里叟复活了叔能妇人的儿子\EditorialNote{列下 4}。但我们相信死亡不是惩罚的报应，因为连圣人也经历了它；甚至圣人的主、耶稣基督，信者的生命和死者的复活，也经历了它。因此，按照古时的做法，对于居住在大城中的人，在战斗之后，祂暂时带来解体，以便当祂复活每个人时，可以弃绝或加冕他。因为用尘土造了亚当身体的那位，也会在解体之后复活其余人的身体和最初那个人的身体（以偿还人的理性本性所欠的债，我们指的是历代的存续。因此，带来解体的那一位，也会实现复活。）因为祂曾说：\BibleQuote{上主从地上取尘土，造了人，将生命之气吹在他脸上，人就成了一个有灵的生物} \BibleRef{创 2:7}，在违命之后又说：\BibleQuote{你原是土，还要归于土} \BibleRef{创 3:19}，同一位后来应许了我们复活。因为祂说：\BibleQuote{凡在坟墓里的，都要听见天主子的声音，而听见的人必要生活。} \BibleRef{若 5:25} 除了这些论证，我们也从我们主的复活相信将有复活。因为祂叫拉匝禄在坟墓四天后复活\EditorialNote{若 11}，叫雅依洛的女儿复活\EditorialNote{谷 5}，又叫寡妇的儿子复活\EditorialNote{路 7}。祂自己在父的命令下，三天内复活了自己，祂是我们复活的保证。因为祂说：\BibleQuote{我是复活，是生命。} \BibleRef{若 11:25} 那位曾在三天内活生生地、毫无损伤地将约纳从鱼腹中带出\EditorialNote{约 2}，将三个青年从巴比伦的火窑中带出，将达尼尔从狮子口中带出的那一位，也不缺乏使我们复活的能力。但如果外邦人嘲笑我们，不相信我们的圣经，至少让他们自己的女先知西比拉迫使她们相信，她明确地向他们说：\par

\Quote{但当一切变为尘土和灰烬，\newline{}不灭的天主熄灭了所燃的火，\newline{}天主将再次用那些骨头和灰烬造人，\newline{}并让终有一死的人恢复原状。\newline{}那时将有审判，天主将施行公义，\newline{}并再次审判世界。但所有因不敬虔而犯罪的凡人\newline{}将再次被覆盖在尘土下；\newline{}但所有敬虔的人将再次在世上生活。\newline{}当天主将祂的神放在他们内，并同时赐给敬虔者生命和恩惠，\newline{}那时所有人都将看见自己。}\par

因此，如果这位女先知宣认复活，不否认万物的恢复，并区分虔敬者与不虔敬者，那么他们否认我们的教义便是徒然的。不仅如此，他们还说他们能展示复活的一个类似物，而他们自己并不相信他们所宣称的事：因为他们说有一种独一无二的鸟，为复活提供了一个充足的证明，他们说这种鸟没有配偶，是受造界中唯一的。他们称它为\OriginalTerm{凤凰}{phœnix}，并讲述说，每五百年它来到\Place{埃及}，来到那称为太阳祭台的地方，带来大量的肉桂、桂皮和香脂木，面向东方，正如他们所说，向太阳祈祷，自行燃烧，化为灰烬；但从那些灰烬中又生出一条虫，当虫受热时，便成形为新生的凤凰；当它能飞时，便飞往埃及诸国之外的\Place{阿拉伯}。因此，如果正如他们自己所说，藉着一只无理性的鸟就展示了一个复活，那么他们为何徒然贬低我们的叙述呢？我们宣认，那凭祂的权能使先前不存在之物得以存在的那一位，有能力恢复这个身体，并在它分解后再次复活它。因为基于这望德的完全确证，我们才忍受鞭打、迫害和死亡。否则，如果我们没有对这些诺言的完全确证——我们自称是这些诺言的宣讲者——我们便是无端忍受这些事情了。因此，正如我们相信\Person{梅瑟}所说的：\BibleQuote{在起初天主创造了天地。}\BibleRef{创 1:1}并且我们知道祂并不需要质料，而只凭祂的意愿便使那些基督受命创造之物得以存在：即天、地、海、光、夜、昼、光体、星辰、飞鸟、鱼类、四足兽、爬行动物、植物和草木；同样地，祂也将凭祂的意愿使众人复活，不需要任何帮助。因为创造世界和复活死者是同一种大能的工作。然后，祂从不同的部分造了人——在此之前人并不存在，赐予他一个从无中受造的灵魂。但现在，祂要将已分解的身体归还给仍然存在的灵魂：因为复活属于已存之物，不属于无有之物。因此，那从无中造出原始身体并赋予它们各种形态的，也必再次复活并兴起那些死者。因为那在母腹中用微小种子形成人，并在人内创造了一个先前不存在的灵魂的那一位——正如祂自己某处对\Person{耶肋米亚}所说：\BibleQuote{我在母腹中形成你以前，我已认识了你；}\BibleRef{耶 1:5}又另一处说：\BibleQuote{我是上主，祂建立了诸天，奠定了大地根基，并在人内形成了人的精神，}\BibleRef{匝 12:1}——也必兴起众人，作为祂的化工；正如神圣的圣经也证明天主对祂的独生子基督说：\BibleQuote{让我们按我们的肖像，按我们的模样造人。天主于是造了人：按天主的肖像造了祂；造了他们，男和女。}\BibleRef{创 1:26}而极其神圣而坚忍的\Person{约伯}——圣经记载说：\BibleQuote{他必与主所复活的人一同复活}——向天主这样说：\BibleQuote{祢不是曾像倒奶一样倒出我，使我凝结如同干酪吗？祢以皮肤和肌肉覆盖了我，以骨和筋编织了我。祢赐给了我生命和恩宠，祢的眷顾保全了我的精神。我心中有这些，我知道祢能做一切事，没有什么是祢不可能的。}\BibleRef{约 10:10}因此，我们的救主和主\Person{耶稣基督}也说：\BibleQuote{在人不可能的，在天主却可能。}\BibleRef{路 18:27}\Person{达味}，天主所爱的，说：\BibleQuote{祢的手造了我，形成了了我。}又说：\BibleQuote{祢知道我的结构。}之后又说：\BibleQuote{祢形成了了我，将祢的手放在我身上。祢的知识被宣称为对我太奇妙；它非常伟大，我不能达到。}\BibleQuote{祢的眼睛看见我尚未成形的体质；众人都要写在祢的书上。}不仅如此，\Person{依撒意亚}在向祂的祈祷中说：\BibleQuote{我们是泥土，祢是我们的陶工。}\BibleRef{依 64:8}因此，如果人是祂的化工，由\Person{基督}所造，那么在他死后，必定由祂再次复活，目的是或为他的善行加冕，或为他的过犯受罚。但如果祂，作为立法者，以公义审判；正如祂惩罚不虔敬者，祂也施恩并拯救信者。那些为祂的缘故被人杀害的圣人们，\BibleQuote{祂要使其中一些如星辰发光，使另一些如光体明亮，}\BibleRef{达 12:3}正如\Person{加俾额尔}对\Person{达尼尔}所说的。所以我们所有信者，即\Person{基督}的门徒，相信祂的许诺。因为那许诺者不会说谎；正如蒙福的先知\Person{达味}说：\BibleQuote{主在祂一切言语上是信实的，在祂一切作为上是圣洁的。}因为那为自己从童贞女取得身体的，也是其他众人的创造者。那使自己从死者中复活的，也必复活所有已倒下的人。那从一粒麦子使地上长出多茎的麦穗的，那使砍下的树发出新枝的，那使\Person{亚郎}的枯杖发芽的，祂必使我们荣耀地复活；那使瘫痪者完全康复并医治枯手者的，那用泥土和唾沫为生来盲者补足缺陷的，祂必使我们复活；那用五个饼两条鱼饱食五千人并剩下十二筐，从水变酒，藉我\Person{伯多禄}从鱼口中送出钱币给收税者的，祂必复活死者。因为我们为祂作证这一切事，先知们也证明其他事。我们这些与祂一同吃喝、亲眼目睹祂的奇事、祂的生活、祂的品行、祂的言语、祂的苦难、祂的死亡、祂从死者中的复活，并在祂复活后与祂来往四十天的人，\BibleRef{宗 1:3}并从祂领受了命令，向普世传福音，使万民成为门徒，\BibleRef{玛 28:19}并藉着宇宙的天主——祂的圣父的权柄，和圣神——祂的护慰者的见证，因祂的死亡给他们授洗——我们教导你们这一切，是祂藉着祂的宪章所吩咐我们的，在\BibleQuote{祂被接升天}\BibleRef{宗 1:9}回到那派遣祂者之前。如果你们相信，你们便有福；如果你们不相信，我们将被视为无罪，与你们的不信无关。\par

% structure: p0019 source=h3 target=subsubsection reason=relative-heading-rank; base=h2

\subsubsection{論主之弟兄雅各伯及首位殉道者斯德望}

八、如今論及殉道者，我們告訴你們，他們應當在你們中間備受尊敬，如同我們尊敬有福的雅各伯主教及聖斯德望我們的同僕。因為這些人為天主所視為有福，並為聖潔之人所尊敬，他們純潔無過，面對罪的誘惑堅定不移，未被說服偏離善工，無可爭議地堪受讚頌。達味也論及他們說：\Quote{在上主眼中，聖者的逝世是寶貴的。} 撒羅滿也說：\Quote{義人的紀念備受讚揚。} \BibleRef{箴 10:7} 先知也論及他們說：\Quote{義人被接去。}\par

% structure: p0021 source=h3 target=subsubsection reason=relative-heading-rank; base=h2

\subsubsection{論假殉道者}

九、以上我們所說的，是論及那些確實為基督殉道的人，而非論及假殉道者。關於後者，神諭說：\Quote{惡人的名聲必朽壞。} \BibleRef{箴 10:7} 因為\Quote{忠實的證人不說謊，但不義的證人惹動謊言。} \BibleRef{箴 14:5} 因為那在作證時離世、不為真理而說謊的人，乃是忠信的殉道者，值得相信他在那些為敬虔之言而流血奮鬥的事上。\par

% structure: p0023 source=h2 target=subsection reason=relative-heading-rank; base=h2

\subsection{第二節 應避免一切與偶像的往來}

% structure: p0024 source=h3 target=subsubsection reason=relative-heading-rank; base=h2

\subsubsection{勸勉道德：我們應戒絕虛浮言談、穢語、戲謔、醉酒、淫蕩、奢侈}

十、現在我們勸誡你們，弟兄及同僕們，要避免虛浮言談、穢語、戲謔、醉酒、淫蕩、奢侈、無節制的私慾及愚妄的言論。因為我們甚至不允許你們在喜樂的主日裡說或做任何不合宜的事。因為聖經某處說：\Quote{當以敬畏事奉上主，以戰兢而歡樂。} 因此，連你們的喜樂也當以敬畏和戰兢來行。因為一個忠信的基督徒，既不應重複異教徒的讚美詩，也不應唱誦穢歌，否則他將因那讚美詩而被迫提及那令人憎惡的魔鬼之名；並且，代替聖神，那惡者將進入他內。\par

% structure: p0026 source=h3 target=subsubsection reason=relative-heading-rank; base=h2

\subsubsection{勸導人避免可憎的偶像崇拜之罪}

十一、你們也被禁止指著它們起誓，或從口中說出它們可憎的名字，也不可敬拜它們或懼怕它們如同神。因為它們不是神，而是邪惡的魔鬼或人荒謬的發明。因為天主關於以色列人曾說：\Quote{他們離棄了我，指著非神起誓。} \BibleRef{耶 5:7} 後來又說：\Quote{我必從他們的口中除去你們偶像的名號。} \BibleRef{匝 13:2} 另一處說：\Quote{他們以非神激起我的妒火，以偶像惹動我的怒氣。} 在全部聖經中，這些事皆為上主天主所禁止。\par

% structure: p0028 source=h3 target=subsubsection reason=relative-heading-rank; base=h2

\subsubsection{我們不應唱異教或穢褻之歌，也不應指著偶像起誓；因為這是不敬之事，違反對天主的認識}

十二、立法者不僅禁止我們關於偶像的事，也警告我們關於眾光體，不可指著它們起誓，也不可事奉它們。因為他們說：\Quote{免得你觀看太陽、月亮和星辰，就被引誘去敬拜它們。} 另一處說：\Quote{不要學習隨從異邦人的行徑，也不要懼怕天上的徵兆。} 因為星辰和眾光體是賜給人為照亮他們，而非為敬拜；雖然以色列人因他們乖謬的性情，\Quote{敬拜受造之物而不敬拜造物主，} 侮辱了他們的製造者，過分崇拜受造之物。有時他們製造牛犢，如在曠野中；有時敬拜\OriginalTerm{巴耳培敖爾}{Baalpeor}；另一時拜\OriginalTerm{巴耳}{Baal}，以及\OriginalTerm{塔慕次}{Thamuz}\EditorialNote{厄則克耳 8:14}，和漆冬的\OriginalTerm{阿斯塔特}{Astarte} \BibleRef{列上 11:5}；又拜\OriginalTerm{摩肋客}{Moloch}和\OriginalTerm{革摩士}{Chamos} \BibleRef{列上 11:7}；另一時拜太陽，\BibleRef{則 8:16} 如厄則克耳所記載；甚至拜走獸，如埃及的\OriginalTerm{阿皮斯}{Apis}和\OriginalTerm{孟底斯的山羊}{Mendesian goat}，以及金銀的神像，如在猶大。因這一切，上主警告他們，並藉先知說：\Quote{猶大家行這些可憎的事，難道是小事嗎？因為他們使這地充滿強暴，惹我發怒。看，他們竟像那些嘲笑的人。我也必發烈怒。我的眼必不顧惜，也不憐憫；他們必大聲向我耳中呼求，我也不垂聽。} \BibleRef{則 8:17} 親愛的，請看，上主針對拜偶像者及拜太陽月亮者宣告了多少事！因此，一個屬天主的人，作為基督徒，不應指著太陽、月亮或星辰起誓，也不應指著天、地、或任何元素，無論大小。因為如果我們的主命令我們不可指著真天主起誓，為使我們的話語比誓言更堅定；也不可指著天本身——因為那是異教徒的邪惡——也不可指著耶路撒冷、天主的聖所、祭台、禮物、祭台的飾金、或自己的頭——因為這習俗是猶太人的敗壞，因此被禁止——；並且祂勸勉忠信的人，他們的\Quote{是}就是\Quote{是}，他們的\Quote{非}就是\Quote{非}，並說\Quote{多餘的便是出於邪惡者}；那麼，那些在起誓時引證虛假神明，並榮耀想像之物而非真實者，天主因他們的乖謬\Quote{任憑他們陷於妄謬，去行那些不合理的事} \BibleRef{羅 1:28}，他們的過失豈不更大嗎？\par

% structure: p0030 source=h2 target=subsection reason=relative-heading-rank; base=h2

\subsection{第三節 論慶節與齋戒日}

% structure: p0031 source=h3 target=subsubsection reason=relative-heading-rank; base=h2

\subsubsection{應守之主的慶節列表，以及各當何時舉行}

XIII. 弟兄们，要遵守庆节；首先是生日，你们要在第九个月的第二十五日庆祝；此后，主显节应受到你们最大的尊崇，因为主在这一天向你们显示了他的天主性，定在第十个月的第六日；此后，应遵守四旬期的斋戒，因为它包含了对我们的主的生活方式与立法的纪念。但这庄严的节期应在逾越节斋戒之前，从一周的第二天开始，直到预备日结束。在这些庄严的节期之后，中断你们的斋戒，开始圣逾越周，你们全体都要以恐惧和战栗守斋，在其中为那些将要灭亡的人祈祷。\par

% structure: p0033 source=h3 target=subsubsection reason=relative-heading-rank; base=h2

\subsubsection{论我们的主的受难，以及他苦难中的每一天所发生的事；并论犹达斯，以及主交付奥迹给门徒时犹达斯不在场。}

第十四章. 因为他们在一周的第二日，即正月，就是\OriginalTerm{山提古月}{Xanthicus}，开始召开会议反对主；商议持续到一周的第三日；但在第四日，他们决定以钉十字架的方式夺去祂的性命。\Person{犹达斯}知道这事，他长久以来已经堕落，但那时被魔鬼本身以贪财之心所击打，虽然他长久受托管理钱囊，\BibleRef{若 12:6} 且惯常偷取为穷人预留的款项，但主因极大的耐心并没有弃绝他；不仅如此，当我们一次与祂同席时，祂愿意既使他归正，又以其预知教导我们，便说：\Quote{我实实在在告诉你们，你们中间有一个人要出卖我。} 我们每个人都说：\Quote{是我吗？} 主保持沉默，我——十二位宗徒之一，且比其余人更受祂钟爱——从靠在祂怀里起来，恳求祂告诉谁将是出卖祂的人。但即使那时，我们善主也没有说出他的名字，却给了两个关于出卖者的记号：一个说：\Quote{同我蘸手在盘子里的}；第二个说：\Quote{我蘸了\OriginalTerm{蘸饼}{sop}递给谁。} 甚至当他自己说：\Quote{师傅，是我吗？} 主并没有说是，却说：\Quote{你说的是。} 为要在这事上惊吓他，祂说：\Quote{出卖人子的人有祸了！那人若没有生下来，为他倒好。他听了这话，就出去对司祭们说：\InnerQuote{你们愿意给我什么，我就把祂交给你们？} 他们便同他议定了三十块钱。} \BibleRef{玛 26:15} 经上所载的话就应验了：\BibleQuote{他们取了那三十块钱，就是以色列子民为祂所估定的价钱，买了窑户的一块田。} \BibleRef{玛 27:9} 在一周的第五日，我们与祂同吃了逾越节，\Person{犹达斯}将手蘸在盘子里，接过了蘸饼，夜间出去以后，主对我们说：\Quote{时候到了，你们将要分散，撇下我独自一人。} \BibleRef{若 16:32} 每人都激烈地声明不会离弃祂，我\Person{伯多禄}还加上这个许诺，说愿意与祂同死，祂便说：\Quote{我实在告诉你，鸡叫以前，你要三次否认认识我。} \BibleRef{路 22:34} 当祂把祂宝贵体血的\OriginalTerm{象征奥迹}{representative mysteries}交付我们时，\Person{犹达斯}不在我们中间，祂就出去，到了\Place{橄榄山}附近，\Place{克德龙溪}旁的一个园子；\BibleRef{若 18:1} 我们与祂同在，并照常规唱了赞歌。\BibleRef{玛 26:30} 祂离开我们不远，向父祈祷说：\Quote{父啊，把这杯撤去；但不要随我的意愿，而要随你的意愿。} 祂这样做三次后，我们因沮丧而睡着了，祂来对我们说：\Quote{时候到了，人子被交在罪人手中。看，\Person{犹达斯}，以及与他同来的一大群不虔敬的人。} \BibleRef{路 22:47} 他们给他看那出卖祂的暗号——一个欺骗的亲吻。他们接到约定的暗号后，就抓住主，捆绑了祂，带到\Person{盖法}大司祭的家里，那里聚集了许多人，不是百姓，而是一大群暴徒，不是圣洁的公议会，而是恶人的集会和不可敬者的议会；他们向祂行了许多事，没有留下任何一种伤害：唾沫吐在祂身上，对祂吹毛求疵，打祂，搧祂的脸，辱骂祂，试探祂，从祂那里寻求虚假的占卜以代替真实预言，称祂为骗子、亵渎者、违反\Person{梅瑟}者、毁坏圣殿者、废止祭献者、罗马人的敌人、\Person{凯撒}的对头。这些公牛和犬类疯狂地向祂投掷这些辱骂，直到清晨很早的时候；然后他们把祂带到\Person{亚纳斯}那里，他是\Person{盖法}的岳父；他们在那里也做了同样的事，那天是预备日，他们就把祂交给罗马总督\Person{比拉多}，控告祂许多重大的事，却一件也不能证实。于是总督不耐烦地对他们说：\Quote{我查不出祂有什么罪。} \BibleRef{路 23:14} 但他们带来两个说谎的见证人，想要诬告主；但他们被发现有矛盾，证词不一致，他们就改口控告祂谋反，说：\Quote{这个人说祂是王，又禁止纳税给\Person{凯撒}。} \BibleRef{路 23:2} 他们自己成了控告者、见证人、审判官和判决的制定者，喊着说：\Quote{钉祂在十字架上！钉祂在十字架上！} \BibleRef{路 23:21} 这就应验了先知论祂所写的：\BibleQuote{不义的见证人聚集起来攻击我，不义向自己说谎。} 又说：\BibleQuote{许多犬类围绕我，恶人的集会议围攻我。} 又在别处说：\BibleQuote{我的产业对我好像林中的狮子，发声攻击我。} \BibleRef{耶 12:8} 因此\Person{比拉多}因懦弱而玷污了他的权柄，更重视群众而不重视这义人，自己定了恶行，因为他一面作证祂是无辜的，一面又如同有罪一般把祂交出去受十字架的刑罚，尽管罗马人制定了法律：没有被定罪的人不得处死。但行刑者拿住荣耀的主，把祂钉在十字架上，在第六时辰钉了祂，但祂被判刑的判决是在第三时辰接到的。之后，他们给祂喝醋调着苦胆。然后拈阄分祂的衣服。然后他们把两个凶犯与祂同钉，一边一个，以应验所记载的：\BibleQuote{他们给我苦胆吃，我渴了，他们给我醋喝。} 又说：\BibleQuote{他们分了祂的外衣，又为祂的内衣拈阄。} 在另一处：\BibleQuote{我被列于罪犯之中。} \BibleRef{依 53:12} 然后有黑暗三个时辰，从第六时辰到第九时辰，傍晚又有光，正如所记载的：\BibleQuote{不是白昼，也不是黑夜，到晚上必有光。} 所有这一切，当那些与祂同钉的凶犯看见时，其中一个指责祂，好像祂软弱不能自救；但另一个斥责同伴的无知，转向主，因被祂光照，并承认那受苦的是谁，就祈求祂在将来的国里记念他。祂随即赦免了他以前的罪，领他进入乐园享受神秘的福乐；祂也在第九时辰大声喊叫，对父说：\Quote{我的天主！我的天主！你为什么舍弃了我？} 稍后，祂大声喊着说：\Quote{父啊，宽恕他们，因为他们不知道他们做什么。} 又加上：\Quote{父啊！我把我的灵魂交托在你手中。} 就断了气，在日落前被安放在一个新坟墓里。在一周的第一日破晓时，祂从死人中复活了，完成了祂在受难前向我们预言的事，说：\Quote{人子必须在地心里三天三夜。} 祂从死人中复活以后，首先显现给\Person{玛利亚玛达肋纳}和\Person{雅各伯}的母亲\Person{玛利亚}，然后显现给路上的\Person{克罗帕}，之后显现给我们这些因怕犹太人而逃散的门徒，但我们私下却非常关切祂。但这些事也写在福音书上。\par

% structure: p0035 source=h3 target=subsubsection reason=relative-heading-rank; base=h2

\subsubsection{关于圣周，以及为何命令我们在星期三和星期五守斋}

十五、祂因此亲自命令我们因犹太人的不敬和过犯而守这六天斋，同时也吩咐我们为他们哀悼，为他们的丧亡而悲伤。因为主自己也曾\BibleQuote{为他们哭泣，因为他们不知道眷顾他们的时候。}但祂命令我们在每周的第四日和第六日守斋：前者因祂被出卖，后者因祂的受难。但祂指定我们在第七日鸡叫时开斋，却要在安息日守斋。这并不是因为安息日本是守斋之日，因为它是创造后的安息，而是因为我们只应在这一个安息日守斋，因为在这一天造物主在阴间。因为正是在他们的庆节上，他们逮捕了主，为使那预言得以应验：\BibleQuote{他们在庆节中央放置了他们的标记，却不认识。}因此你们应当为他们哀悼，因为主来到时，他们不信祂，反而弃绝祂的教导，自认为不配得救恩。你们却是幸福的，你们从前不是子民，如今却是神圣的民族，从偶像的欺骗、无知和不敬中得救，从前未蒙怜悯，如今因你们衷心的顺从而蒙了怜悯；因为对你们这些归化的外邦人，生命之门已经敞开，你们从前不被爱，如今却被爱；是特为天主拥有的子民，为彰显祂的美德。关于你们，我们的救主曾说：\BibleQuote{我不寻求我的人找到了我，我不询问我的人向我显现。我说：看我！向那并不呼求我名的民族。}因为当你们不寻求祂时，你们被祂所寻求；你们这些信靠祂的人听从了祂的召唤，离开了多神论的疯狂，逃向真实的独一统治，即全能的天主，藉着基督耶稣，并成为得救者数目的圆满——\BibleQuote{万万千千和千万的群众；}\BibleRef{达 7:10} 正如在达味书中所写的：\BibleQuote{在你旁边倒下一千，在你右边一万；}又：\BibleQuote{天主的战车是万万千千，胜利者成千上万。}但对于不信的以色列，祂说：\BibleQuote{我整天向悖逆顶嘴的百姓伸出双手，他们行走不善之道，却随从自己的罪恶，是常在我面前惹我发怒的百姓。}\BibleRef{依 65:2}\par

% structure: p0037 source=h3 target=subsubsection reason=relative-heading-rank; base=h2

\subsubsection{列举预告基督的先知预言，虽然犹太人看见应验，却因心性邪恶不信祂是天主的基督，并将荣耀的主定罪于十字架}

XVI. 看，百姓因不信主而如何激怒了主！因此祂说：\Quote{他们激怒了圣神，祂就变成了他们的仇敌。}\BibleRef{依 63:10} 因为他们的心思邪恶，瞎眼就落在了他们身上，因为他们看见耶稣时，不相信祂就是天主的基督，祂在万世以前由祂所生，祂的独生子，天主圣言，他们因不信而不承认祂，既不因祂的奇能异事，也不因关于祂所写的预言。因为祂要由童贞女诞生，他们读过这预言：\Quote{看，童贞女将怀孕，生一个儿子，他们将称祂的名字为厄玛奴耳。}\Quote{因为有一个婴孩为我们诞生了，有一个儿子赐给了我们，祂的权柄在祂的肩上；祂的名字称为伟大会谋天使、奇妙谋士、全能的天主、宰制者、和平之王、未来时代之父。} 因为他们极其邪恶而不相信祂，主用这些话表明：\Quote{谁相信了我们的报告？主的臂膀向谁显示了呢？} 之后：\Quote{你们听是要听见，却不明白；看是要看见，却不领悟；因为这百姓的心变迟钝了。} 因此知识从他们那里被夺走，因为看见却忽略，听见却听不见。但你们，归化的外邦人，国度赐给了你们，因为你们这些不认识天主的，因宣讲而相信，\Quote{认识了祂，或者更恰当地说，被祂所认识，} 借着耶稣，那拯救和救赎凡希望于祂之人的救主和救赎者。因为你们从从前虚浮而烦冗的生活方式被迁移，藐视了无生命的偶像，轻视了在黑暗中的魔鬼，投奔了\Quote{真光，}\BibleRef{若 1:9} 并借着它\Quote{认识了独一的真天主和父，}\BibleRef{若 17:3} 因而被承认为祂国度的继承人。既然你们\Quote{受洗归于主的死亡}\BibleRef{罗 6:3} 和祂的复活，如同\Quote{新生婴儿，}\BibleRef{伯前 2:2} 就应当完全脱离一切罪恶的行为；\Quote{因为你们不是自己的人，而是用祂自己的血买了你们的那位的人。}\BibleRef{格前 6:19} 关于先前的以色列，主因他们的不信这样说：\Quote{天主的国必从他们夺去，赐给那结果子的民族；}\BibleRef{玛 21:43} 这就是说，将国度赐给你们这些从前远离祂的人，祂期待你们感恩和正直的果实。因为你们是那被派往葡萄园却未听从的人，而这些人却是听从的；但你们为否认而悔改了，现在在里面工作。他们却因自己的盟约不安，不仅让葡萄园荒芜，还杀死了葡萄园主的管家——一个被石头打死，一个被刀杀；一个被锯成两半，一个在圣所被杀，\Quote{在圣殿和祭台之间；} 最后甚至\Quote{把继承人自己扔出葡萄园，杀死了祂。} 他们弃绝祂如同无用的石头，\BibleRef{玛 21:42} 但你们接纳祂为基石。因此关于你们祂说：\Quote{我所不认识的百姓事奉了我，一听见就服从了我。}\par

% structure: p0039 source=h3 target=subsubsection reason=relative-heading-rank; base=h2

\subsubsection{如何庆祝逾越节。}

XVII. 因此，弟兄们，你们是被基督宝血所赎的，应当精确地、小心地遵守逾越节的日子，在春分之后，免得一年中两次记念同一苦难。只为那一次死亡而一年一次守节。\par

你们不要自己计算，而要在你们受割礼的弟兄计算时一同守节；若他们计算有误，你们不必担忧。在无酵饼日中间守住夜。当犹太人欢宴时，你们要禁食，为他们哀号，因为在他们的节日他们钉死了基督；当他们哀哭并吃苦酵饼时，你们要欢宴。但不再与犹太人一起守节，因为我们与他们已无共融；因为他们在这计算本身也入了迷途，他们以为计算完美，其实处处被误导，远离了真理。但你们要仔细注意春分，它在第十二月（即狄斯特罗月，三月）的二十二日发生，小心观察到月亮的二十一日，免得月十四落在另一周，因错误而一年庆祝两次逾越节，或在一周以外的日子庆祝我们主的复活日。\par

% structure: p0042 source=h3 target=subsubsection reason=relative-heading-rank; base=h2

\subsubsection{关于大逾越周的一条宪令。}

XVIII. 你们在逾越节的日子禁食，从第二日（星期一）开始直到预备日（星期五）和安息日，共六天，只吃面包、盐、蔬菜，喝清水；这些天禁酒和肉，因为这是哀悼的日子，不是欢宴的日子。你们中能者应在预备日和安息日完全禁食，直到夜里鸡啼为止；若有人不能连续两天禁食，至少应遵守安息日；因为主在某处论及自己说：\Quote{当新郎从他们中被带走时，那些日子他们就要禁食。} 在这些日子，祂被那些假称犹太人的从我们中带走，钉在十字架上，\Quote{被列于罪犯之中。}\BibleRef{依 53:12}\par

% structure: p0044 source=h3 target=subsubsection reason=relative-heading-rank; base=h2

\subsubsection{关于大安息日通宵守夜，以及复活日。}

十九、因此，我们劝你们在这些日子守斋，如同我们被祂带走时一直斋戒到晚上一样；但在预备日之前的其余日子，让每个人在第九时辰或晚上，或按各人的能力进食。但从第五日的傍晚直到鸡鸣时，当一周的第一日，即主日破晓时，你们要开斋。从傍晚直到鸡鸣时，你们要警醒，在教会中聚集，守望、祈祷并恳求天主；当你们彻夜不眠时，诵读法律、先知书和圣咏集，直到鸡鸣时，并为你们的慕道者施洗，以敬畏和战栗之心诵读福音，并向百姓宣讲有益于他们得救的事：结束你们的悲伤，恳求天主使以色列归化，并赐给他们悔改的机会，赦免他们的不敬。因为那外邦的审判官\Quote{洗手说：\InnerQuote{我流这义人的血，是清白无辜的，你们看着办吧。}但以色列人喊道：\InnerQuote{祂的血归在我们和我们的子孙身上。}}\BibleRef{玛 27:24}当比拉多说：\Quote{我该把你们的王钉在十字架上吗？}他们喊道：\Quote{我们没有王，只有凯撒；钉祂在十字架上，钉祂在十字架上！因为凡自称为王的，就是反对凯撒。}又说：\Quote{你若释放这个人，就不是凯撒的朋友。}总督比拉多和黑落德王便下令将祂钉在十字架上；于是那预言应验了，说：\Quote{外邦人为什么愤怒，万民为什么妄图虚妄的事？世上的君王起来，首领聚集在一起，敌对主和祂的基督。}又说：\Quote{他们抛弃了所爱的，如同一个可憎的死人。}\BibleRef{依 14:19}祂既在预备日被钉死在十字架上，又在主日破晓时复活，经上的话就应验了，说：\Quote{天主，求祢起来，审判大地，因为祢要在万民中承受产业。}又说：\Quote{我要起来，主说；我要将祂安置在安全中，我因祂而勇敢。}又说：\Quote{但是，主，求祢怜悯我，使我再起来，我好报复他们。}因此，现在主既已复活，你们也要献上你们的祭献，关于这祭献，祂曾藉我们订立诫命说：\Quote{你们要这样做，为纪念我。}\BibleRef{路 22:19}从此以后，你们要停止守斋，欢喜快乐，守庆节，因为耶稣基督，我们复活的保证，已从死人中复活。这要成为永久的法令，直到世界的终结，直到主的来临。因为对于犹太人，主仍然是死的；但对于基督徒，祂已复活：前者是由于他们的不信，后者是由于他们充分确定的信德。因为对祂的望德就是不死的永生。八天后，要举行另一个应受尊敬的庆节，即第八天本身，在那天祂使我多默——难以相信的人——完全确信，让我看钉痕和枪刺在祂肋旁的伤痕。\BibleRef{若 20:25}再者，从第一个主日起计算四十天，从主日到一周的第五天，庆祝主升天节，在那一天祂完成了祂全部的管理和规定，回到那派遣祂的天主和父那里，并坐在大能的右边，直到祂的仇敌被放在祂脚下；祂也将在世界的终结时，带着大能和极大的光荣来临，审判活人与死人，按照各人的行为报应各人。那时，他们必看见他们所刺透的天主的爱子；\BibleRef{匝 12:10}当他们认识祂时，他们要为自己哀悼，支派一队一队，妻子们各在一处。\par

% structure: p0046 source=h3 target=subsubsection reason=relative-heading-rank; base=h2

\subsubsection{关于基督耶稣的先知性预言}

二十. 因为即使如今，在\OriginalTerm{果尔皮月}{Gorpiæus}的第十日，当他们聚集时，他们诵读\BibleBook{耶肋米亚}{哀歌}，其中写着：\BibleQuote{我们面前的灵，基督、主被捉在他们的毁灭中；}\BibleRef{哀 4:20} 以及\BibleBook{}{巴路克}，其中写着：\Quote{这是我们的天主；没有别的能与祂相提并论。祂发现了知识的一切途径，并将它显示给祂的儿子\Person{雅各伯}和祂所爱的\Person{以色列}。后来祂出现在地上，与人类交往。} 当他们诵读这些时，他们哀伤哭泣，自以为是哀叹那由\Person{拿步高}所造成的荒凉；但事实上，他们不情愿地为那将要临到他们的哀叹作了前奏。但在升天后十天，即从第一个主日算起的第五十天，要举行盛大庆节：因为在那天的第三时辰，\Person{主耶稣}赐给我们\OriginalTerm{圣神}{Holy Ghost}的恩赐，我们充满了祂的能力，我们\BibleQuote{以新的语言说话，正如那圣神感动我们；}\BibleRef{宗 2:4} 并向\Person{犹太人}和\OriginalTerm{外邦人}{Gentiles}宣讲，祂是天主的基督，是祂\BibleQuote{所预定为审判活人死人的那位。}\BibleRef{宗 10:42} 关于祂，\Person{梅瑟}曾作证说：\BibleQuote{主从主那里取火，并将它降下。}\BibleRef{创 19:24} \Person{雅各伯}曾看见祂像一个人，说：\Quote{我面对面看见了\Person{天主}，我的灵魂得保存。} \Person{亚巴辣罕}曾款待祂，并认出祂是审判者，他的主。\Person{梅瑟}曾在荆棘中看见祂；关于祂，他在\WorkTitle{申命纪}中说过：\BibleQuote{主你的天主要从你弟兄中给你兴起一位先知像我；你们要听从他，凡他对你们所说的一切。凡不听从那先知的灵魂，必从民间消灭。}\EditorialNote{申命纪 18:15} \Person{农的儿子若苏厄}看见祂，祂是主军队的元帅，身着盔甲，来帮助他们攻打\Place{耶里哥}；若苏厄向祂俯伏朝拜，如同仆人向主人一样。\BibleRef{苏 5:14} \Person{撒慕尔}知道祂是\Quote{\Person{天主}的\OriginalTerm{受傅者}{Anointed}，}\BibleRef{撒上 12:3} 并因此称司祭和君王为受傅者。\Person{达味}认识祂，并唱了一首关于祂的诗歌：\Quote{一首关于那可爱者的诗；} 并以祂本人的身份加上说：\Quote{请将你的剑佩在腰间，你是威武英俊的：为了真理、温良和正义，出发，凯旋，为王；你的右手必奇妙地引导你。你的箭是锐利的，大王！万民必仆倒在你脚下，在君王仇敌的心中。因此，\Person{天主}，你的天主，以喜乐之油膏你，胜过你的同伴。} \Person{撒罗满}也谈及祂，以祂本人的身份说：\BibleQuote{主在祂道路的起始就造了我，为祂的作为：在世界之前，在起初，在祂造地之前，在水泉涌出之前，在群山奠定之前，在一切山陵之前，祂就生了我。}\BibleRef{箴 8:22} 又说：\BibleQuote{智慧建造了自己的房屋。}\BibleRef{箴 9:1} \Person{依撒意亚}也论及祂说：\Quote{从\Person{叶瑟}的根须要生出一根枝条，从他的根须要开出一朵花。} 又说：\Quote{必有\Person{叶瑟}的根；那要兴起统治外邦人的，外邦人将信赖祂。} \Person{匝加利亚}说：\Quote{请看，你的君王来到你这里，正义的，带着救恩；温良的，骑在驴上，骑在驴驹上，就是母驴的崽子。} \Person{达尼尔}描述祂是\Quote{\OriginalTerm{人子}{Son of man}来到\Person{父}面前，}并受到父的一切审判和尊荣；又描述祂是\Quote{一块非用手凿的石头，从山中凿出，变成一座大山，充满全地，}粉碎列国的许多政权和众神的多神崇拜，却宣讲唯一的天主，并建立了\Place{罗马人}的帝国。\Person{耶肋米亚}也曾预言论及祂说：\Quote{祂面前的灵，主基督，陷入了他们的网罗：我们曾说，我们将在祂的荫庇下寄居在外邦人中。} \Person{厄则克耳}以及后来的先知们，处处确认祂就是基督、主、君王、审判者、立法者、圣父的使者、\OriginalTerm{独生的天主}{only-begotten God}。因此我们也向你们宣讲祂，并宣告祂是\OriginalTerm{天主圣言}{God the Word}，祂为创造宇宙服务了祂的天主和父。藉着相信祂，你们将得生命，但若不信，你们将受惩罚。因为\BibleQuote{那不信从子的人不得见生命，而天主的义怒常在他身上。}\BibleRef{若 3:36} 因此，你们在庆祝了\OriginalTerm{五旬节}{Pentecost}后，再庆祝一周，然后禁食；因为为天主的恩赐而喜乐，并在那放松之后禁食是合理的：因为\Person{梅瑟}和\Person{厄里亚}都禁食了四十天，\Person{达尼尔}\Quote{三周的日子里没有吃美味的饼，肉和酒没有进入他的口。} 有福的\Person{亚纳}，当她祈求\Person{撒慕尔}时，说：\BibleQuote{我没有喝淡酒和烈酒，而是向主倾吐我的灵魂。}\BibleRef{撒上 1:15} \Place{尼尼微人}禁食三天三夜，\BibleRef{纳 3:5} 逃脱了怒气的执行。\Person{艾斯德尔}、\Person{摩尔德开}和\Person{友弟德}，藉着禁食，逃脱了不虔敬的\Person{敖罗斐乃}和\Person{哈曼}的叛乱。\Person{达味}说：\Quote{我的膝盖因禁食而虚弱，我的肉因缺少油而枯槁。} 因此，你们也要禁食，并向天主求你们的愿望。我们命令你们每周第四日和每个\OriginalTerm{预备日}{the day of the preparation}禁食，并将你们禁食剩余的分施给贫困者；除一个安息日外，每个安息日和每个主日，你们要举行庄严的集会，并要喜乐：因为凡在主日——复活之日，或在五旬节期间，或一般地在主的庆节上禁食的，便是犯罪。因为在那些日子我们应当喜乐，而不当哀伤。\par

\SourceRecord{Translated by James Donaldson.；Ante-Nicene Fathers；Vol. 7.；Edited by Alexander Roberts, James Donaldson, and A. Cleveland Coxe.；Buffalo, NY: Christian Literature Publishing Co.,；1886.；<http://www.newadvent.org/fathers/07155.htm>.}

% structure: p0001 source=h1 target=section reason=page-title

\section{\WorkTitle{宗徒宪章（第六卷）}}

% structure: p0002 source=h2 target=subsection reason=relative-heading-rank; base=h2

\subsection{第一节 论异端}

% structure: p0003 source=h3 target=subsubsection reason=relative-heading-rank; base=h2

\subsubsection{那些胆敢制造分裂且未逃脱惩罚的人}

I. 首先，主教啊，你要避开那些可悲、危险且最无神的异端，像避火一样，这火会烧伤靠近它的人。也要避免分裂，因为将心思转向邪恶的异端，或出于野心而与同道分离，都是不许可的。因为古时那些胆敢如此行的人并未逃脱惩罚。达堂和阿彼兰曾对抗梅瑟，被地吞没。科辣黑和那二百五十个与他一同反叛亚郎的人，被火烧死。米黎盎曾指责梅瑟，被逐出营外七天，因为她说梅瑟娶了位雇士女子为妻。至于阿匝黎雅和乌齐雅\EditorialNote{编年纪下 26}，后者是犹大王，他胆敢篡夺司铎职，想要献香，这是他不该做的，被大司铎阿匝黎雅和八十位司铎阻止；当他不肯听从时，额上便起了癞病，急忙出去，因为主已谴责了他。\par

% structure: p0005 source=h3 target=subsubsection reason=relative-heading-rank; base=h2

\subsubsection{论不可反抗君王或司铎职务}

II. 因此，亲爱的，让我们想一想那些造反者有何等荣耀，有何等惩罚。因为如果反抗君王的人应受惩罚，即使他是儿子或朋友，何况反抗司铎的人呢！因为司铎职比王权更高贵，因为它关心灵魂，所以胆敢反抗司铎职的人所受的惩罚，比胆敢反抗王权的人更大，尽管两者都难逃惩罚。因为阿贝沙隆和阿布达丹\EditorialNote{撒慕尔纪下 18-20}都未逃脱惩罚；科辣黑和达堂\EditorialNote{户籍纪 16}也未逃脱。前者反抗达味，争夺王位；后者反抗梅瑟，争夺高位。他们都说恶言：阿贝沙隆指责他的父亲达味不公义，对每个人说：\BibleQuote{你的话是好的，但没有人听你，为你伸冤。谁立我为统治者呢？}\BibleRef{撒下 15:3} 但阿布达丹说：\BibleQuote{我在达味中没有分，在叶瑟的儿子中没有产业。}\BibleRef{撒下 20:1} 显然他不愿受达味统治，而天主曾说：\BibleQuote{我找到了叶瑟的儿子达味，一个合乎我心的人，他要遵行我的一切命令。}\BibleRef{宗 13:22} 但达堂、阿彼兰和科辣黑的同伙对梅瑟说：\Quote{你把我们从埃及地、从流奶流蜜之地领出来，这算小事吗？为什么你遮住了我们的眼睛？你要统治我们吗？}他们就聚集了大队会众攻击他。科辣黑的同伙说：\Quote{天主只对梅瑟说话吗？为何祂将大司铎职只给了亚郎？主的全体会众不是圣洁的吗？为何只有亚郎拥有司铎职？}在此之先，有人说：\BibleQuote{谁立你作我们的首领和审判官呢？}\BibleRef{出 2:14}\par

% structure: p0007 source=h3 target=subsubsection reason=relative-heading-rank; base=h2

\subsubsection{论梅瑟的德性、犹太民族的不信，以及天主在他们中所行的奇妙作为}

III. 他们起来反抗天主的仆人梅瑟，他是世上最谦和的人\BibleRef{户 12:3}，并且忠信，却以最大的忘恩负义侮辱这样一位伟人。这位伟人是他们的立法者、守护者、大司铎和君王，是神圣事务的管理者；他像造物主一样显示造物主的伟大作为；他是最谦和的人，最无骄傲，充满坚韧，性情最温和；他曾以他的圣德拯救他们脱离许多危险和死亡；他曾在百姓面前奉天主行了那么多征兆与奇事，为他们施行了荣耀而奇妙的作为；他曾在埃及降下十灾；他曾分开红海，使水左右立起如墙，带领百姓从干地走过，而将法郎和埃及人以及所有跟随他们的人淹死；他曾用木头使苦水变甜，曾在他们口渴时从磐石中引水；他曾降下玛纳从天给他们吃，从空中分赐肉给他们；他曾赐给他们夜间以火柱照明引导，白日以云柱遮蔽，因日头酷热；他曾将天主的法律展示给他们，这法律是由天主的口、手和书写刻在石版上，共十条诫命完美无缺。\BibleQuote{天主与他面对面说话，如同人对朋友说话一样；} 祂曾说：\BibleQuote{再没有兴起一位像梅瑟一样的先知。}科辣黑的同伙和勒乌本人起来攻击他，向梅瑟投石；梅瑟祈祷说：\Quote{不要收纳他们的祭献。}于是天主的荣耀显现，将一些人投入地中，另一些人用火烧灭。至于那些说\Quote{让我们为自己立一个首领}的分裂欺骗的首领们，地开口吞没了他们和他们的帐棚及一切所有，他们活活下到了地狱；但祂用火烧灭了科辣黑的同伙。\par

% structure: p0009 source=h2 target=subsection reason=relative-heading-rank; base=h2

\subsection{第二节 异端的历史与教义}

% structure: p0010 source=h3 target=subsubsection reason=relative-heading-rank; base=h2

\subsubsection{分裂不是由那些从不虔敬者中分离出来的人造成的，而是由那些离开虔敬者的人造成的}

四、因此，如果天主对那些因野心制造分裂的人立即施加惩罚，那么对那些不虔敬异端的首领，岂不更要如此？祂岂不要对那些亵渎祂的眷顾或祂的创造的人施加更严厉的惩罚？但是，你们这些受圣经教导的弟兄们，要谨慎，不要在意见上制造分裂，也不要在合一中制造分裂。因为那些建立非法意见的人，是百姓丧亡的标记。同样，你们平信徒不要靠近那些宣扬与天主意旨相悖的教义的人，也不要参与他们的不虔敬。因为天主说：\Quote{你们从这些人中间分别出来，免得与他们一同灭亡。}又说：\Quote{你们从他们中间出来，分别自己，主说，不要触摸不洁之物，我就接纳你们。}\par

% structure: p0012 source=h3 target=subsubsection reason=relative-heading-rank; base=h2

\subsubsection{从先知预言证明：所谓以色列因何被天主弃绝。}

五、因为那些亵渎天主的人必定要躲避。大部分不虔敬的人确实不认识天主；但这些人，作为对抗天主的人，被一种邪恶的意志所占据，如同疾病一般。因为从这些异端者的邪恶中，\Quote{污秽散布在全地上}，正如耶肋米亚先知所说。因为邪恶的会堂如今被主抛弃，祂的房屋被祂弃绝，正如祂在某处所说：\Quote{我撇弃了我的房屋，我丢弃了我的产业。}又说，依撒意亚说：\BibleQuote{我要忽略我的葡萄园，不再修剪，不再挖掘，荆棘要如同在荒野中长出来；我要命令云彩不再降雨在上面。}\BibleRef{依 5:6} 因此，\BibleQuote{祂遗弃了祂的子民，如同葡萄园中的棚屋，如同无花果或橄榄园中的禾堆，如同被围困的城。}\BibleRef{依 1:8} 祂在他们身上取走了圣神和先知之雨，并以灵性恩宠充满祂的教会，如同\Quote{初熟时节的埃及河}；并高举教会如\Quote{山上的房屋，或如一座高山；如同产奶和脂油丰饶的山，天主乐意居住其中。因为主要居住其中直到永远。}祂在耶肋米亚中说：\BibleQuote{我们的圣所是荣耀的崇高宝座。}\BibleRef{耶 17:12} 祂在依撒意亚中说：\BibleQuote{在末日，上主的山必光荣，上主的殿必矗立在众山之巅，超越众山。}\BibleRef{依 2:2} 因此，既然祂遗弃了祂的子民，祂也留下了祂的圣殿荒凉，撕裂了圣殿的幔子，并从他们那里取走了圣神；因为祂说：\BibleQuote{看，你们的房屋被留给你们荒凉。}\BibleRef{玛 23:38} 祂却将灵性恩宠赐予你们这些归化的外邦人，正如祂通过岳厄尔所说：\BibleQuote{此后，天主说，我要将我的神倾注在一切有血肉的人身上；你们的儿子和女儿要说预言，你们的老年人要作梦。}\BibleRef{岳 2:28} 因为天主从那个民族那里取走了祂话语的一切权能和效力，以及类似的眷顾，并将它们转移到你们这些归化的外邦人身上。正因如此，魔鬼对天主的圣教会极其愤怒；他被转移到你们中间，在你们身上引起了敌意、叛乱、侮辱、分裂和异端。因为此前他通过杀害基督已经征服了那个民族。但你们这些离开他虚妄的人，他以不同的方式试探你们，如同他试探有福的约伯一样。确实，他反对大司祭\Person{若苏厄}（\Person{约匝达克}的儿子）\BibleRef{匝 3:1}，他曾多次试图筛我们，使我们的信德失败。但我们的主和师傅，将他带到审判前，对他说：\Quote{愿上主责斥你，撒殚！愿上主责斥你，祂拣选了耶路撒冷。这不是从火中抽出来的一根柴吗？}那时祂对站在大司祭旁边的人说：\Quote{脱下他污秽的衣服，}并加上：\Quote{看，我免除了你的罪恶。}现在祂要说，正如祂从前在我们聚集时所说的：\BibleQuote{我已为你祈祷，使你的信德不致缺失。}\BibleRef{路 22:32}\par

% structure: p0014 source=h3 target=subsubsection reason=relative-heading-rank; base=h2

\subsubsection{即使在犹太人中也出现了多种为天主所憎恶的异端教义。}

六、因为甚至犹太民族也有邪恶的异端：其中有\OriginalTerm{撒杜塞人}{Sadducees}，他们不承认死人复活；\OriginalTerm{法利塞人}{Pharisees}，他们将罪人的行为归因于命运和定命；\OriginalTerm{巴司摩特人}{Basmotheans}，他们否认天主眷顾，说世界是自发运动形成的，并且否定了灵魂的不死不灭；\OriginalTerm{每日洗者}{Hemerobaptists}，他们每日若不洗就不吃饭——不仅如此，若不洁净他们的床、桌子、盘子、杯子、座位，他们就不使用任何东西；还有在我们中间新崛起的\OriginalTerm{厄彼翁派}{Ebionites}，他们坚持天主子只是一个凡人，由人意而出生，由若瑟和玛利亚的结合而生。也有那些从所有这些分离出来、遵守他们祖先律法的人，这就是\OriginalTerm{厄色尼人}{Essenes}。这些出现于前一个民族中。如今，那善于作恶、对善事一无所知的邪恶者，从我们中间驱赶了一些人，并通过他们制造了异端和分裂。\par

% structure: p0016 source=h3 target=subsubsection reason=relative-heading-rank; base=h2

\subsubsection{异端从何而起，谁是不虔敬的罪魁祸首。}

七. 新异端的起源是这样开始的：魔鬼进入一个名叫西满的人里面，他来自一个叫吉特的村庄，是撒玛黎雅人，以行魔术为业，并使他成为其邪恶计划的工具。\EditorialNote{宗 8} 因为当我们的同宗徒斐理伯，藉主的恩赐和祂圣神的德能，在撒玛黎雅施行医治的奇迹，以致撒玛黎雅人深受感动，接受了宇宙的天主和主耶稣的信仰，并奉祂的名受洗；不仅如此，连西满本人，当他看到那些不藉任何魔术仪式而行的征兆和奇事时，也惊叹、相信、受洗，并持续禁食祈祷——我们听说了因斐理伯而在撒玛黎雅人中的天主的恩宠，就下到他们那里；并大大扩充了教义的言语，给所有受洗的人覆手，使他们领受圣神的分享。但当西满看到圣神是因我们的覆手而赐予信徒时，他就拿钱来献给我们，说：\Quote{也赐给我这权柄，使我无论给谁覆手，谁就能领受圣神；}\BibleRef{宗 8:19} 他渴望如同魔鬼因亚当尝了那树而剥夺了他所应许的永生一样，西满也因收钱而引诱我们，从而切断我们与天主的恩赐的联系，以便通过交换，我们将圣神的无价恩赐卖给他。但当我们都为此献礼感到困扰时，我伯多禄，定睛注视他内在的那条恶蛇，对西满说：\Quote{让你的银子与你一同丧亡，因为你以为天主的恩赐可用金钱买到。你在这件事上没有份，在这信仰上没有基业；因为你的心在天主面前不正。因此，你当悔改这邪恶，并祈求主，或许你心中的念头可得赦免。因为我看出你正处于苦胆和不义的束缚中。}但随后西满惊慌地说：\Quote{我恳求你，为我祈求主，使你所说的这一切不致临到我身上。}\EditorialNote{宗 8:24}\par

% structure: p0018 source=h3 target=subsubsection reason=relative-heading-rank; base=h2

\subsubsection{八. 谁是西满不敬神的后继者，以及他们设立了哪些异端}

八. 但是当我们出去到外邦人中宣讲生命之言时，魔鬼就在人民中作工，打发假宗徒跟在我们后面，以败坏圣言；他们打发了一个叫克莱俄彼的，与西满联合，这两人成了多西托的门徒，但他们藐视他，把他从首领的位置上推下来。后来，另一些人成了荒谬教义的创始者：克林托、马尔谷、默南德尔、巴西理得、撒都尼鲁。其中一些人承认多神论，一些人只承认三神，但彼此矛盾，无始无终，永远共存；还有一些人承认无数神，包括那些不可知者。有些人拒绝婚姻，他们的教义是婚姻并非天主的安排；另一些人厌恶某些食物；有些人在不洁中无耻，如那些被错误称为尼苛拉党人的人。当西满在斯特拉托尼斯的凯撒勒雅遇到我伯多禄时（在那里，那有信德的外邦人科尔乃略藉着我信了主耶稣），他就试图曲解天主的话；与我同在的有圣洁的子弟匝凯，他曾是税吏，还有巴尔纳伯；以及克莱孟主教的兄弟尼凯塔斯和阿桂拉，克莱孟是罗马的公民和主教，他是保禄的门徒，保禄是我们的同宗徒和福音的同工。我在他们面前与他三次辩论，论及真先知和天主的独一统治；当我藉主的德能战胜他，使他哑口无言后，我就将他驱往意大利。\par

% structure: p0020 source=h3 target=subsubsection reason=relative-heading-rank; base=h2

\subsubsection{九. 西满如何渴望藉某些魔术飞翔，却因伯多禄的祈祷从高空中头朝下坠落，摔碎了脚、手和踝骨}

九. 当他在罗马时，他大大扰乱了教会，倾覆了许多人，将他们拉拢到自己一方，并以他的魔术技巧令外邦人惊讶，以至于有一次正午时分，他进入他们的剧场，命令人们必须强行把我带入剧场，并应许要飞上天；当所有人民为此屏息以待时，我独自祈祷。他果真被邪魔带到空中，高高飞翔，说他正返回天上，要从那里赐给他们好处。人民向他欢呼如神，我向天伸出手，用心恳求天主藉主耶稣打倒这个祸害，毁掉那些邪魔的势力，它们利用他诱惑和败坏人类；使他摔在地上，受伤但不致死。然后，我定睛看着西满，对他说：\Quote{如果我是天主的人，是耶稣基督的真宗徒，是虔诚的导师，而非欺骗者，像你西满一样，我命令那背弃虔诚的邪恶权势，就是载着行魔术者西满的，放开他们，让他从高处头朝下坠落，让他被那些被他诱惑的人嘲笑。}我说完这些话，西满就失去了能力，头朝下轰然坠落，猛烈撞在地上，摔碎了髋骨和踝骨；人民喊叫说：\Quote{只有一位天主，就是伯多禄真实宣讲的那位。}许多人离开了他；但一些该丧亡的人仍坚持他的邪恶教义。就这样，最无神的西满异端首先在罗马建立；魔鬼也藉着其余假宗徒运作。\par

% structure: p0022 source=h3 target=subsubsection reason=relative-heading-rank; base=h2

\subsubsection{异端彼此之间以及它们与真理之间如何区分}

X. 这一切都有同一个无神论的意图：亵渎全能的天主，宣扬祂是不可知的存在，不是基督之父，也不是世界的创造者；而是不可言说、不可名状、自存的；我们不应使用法律和先知；没有天主眷顾，也没有复活可信；没有审判，也没有报应；灵魂并非不朽；我们只应纵情享乐，不加区分地转向各种敬拜。他们中有人说有许多神，有人说有三位无始之神，有人说有两位无始之神，有人说有无数永世。此外，他们中有人教导人不应当结婚，应当戒绝肉食和酒，声称婚姻、生育和食用某些食物是可憎的；这样，他们作为清醒之人，就能使自己的邪恶见解被当作可信的。他们中有些人绝对禁止吃肉，认为那不是畜生的肉而是具有理性灵魂的受造物的肉，仿佛敢于宰杀它们的人会被指控谋杀之罪。但其他人中有人说只需戒绝猪肉，但可以吃法律上洁净的；并且应当按法律行割损，相信耶稣如同一位圣人和先知。但另有些人教导人应当在污秽中无耻，滥用肉体，行一切不圣之事，仿佛这是灵魂躲避这世界统治者的唯一途径。所有这些都是魔鬼的工具，是义怒之子。\par

% structure: p0024 source=h2 target=subsection reason=relative-heading-rank; base=h2

\subsection{第三节 宗徒所攻击的异端}

% structure: p0025 source=h3 target=subsubsection reason=relative-heading-rank; base=h2

\subsubsection{宗徒宣讲的阐述。}

十一、但我们，作为天主的子女和和平之子，宣讲虔敬的神圣正确言语，宣告唯一的天主，法律和先知的主，世界的创造者，基督之父；不是自存的或自生的，如他们所想，而是永恒的、无始的、居住在不可接近的光明中；不是两位或三位或多位，而是永远唯一；不是不可知或不可言说的，而是由法律和先知所宣讲的；全能者，万有的至高治理者，全能的存在；独生子和万物首生者的天主和父；唯一的天主，一位儿子的父，不是许多的；借由基督造了一位护慰者，造了其他阶层，借由基督是各种受造物的唯一创造者，借由祂也是同一位保存者和立法者；复活的起因，以及由祂所做的审判和报应：这同一位基督甘愿成为人，一生无罪，受苦，从死者中复活，回到派遣祂的那一位。我们也说天主所造的万物都是好的，没有一样是可憎的；凡是用于维持生命的，当正义地分享时，都是极好的：因为按照圣经，\BibleQuote{一切都很好。}\BibleRef{创 1:31} 我们相信合法的婚姻和生育是尊贵无玷的；因为性别的差异是在亚当和厄娃中为增加人类而形成的。我们承认我们拥有一个无形体和不朽的灵魂——不像身体那样可朽，而是不朽的，因为是理性的和自由的。我们憎恶一切非法的结合，以及某些人所行的违反本性的行为，视之为邪恶和不虔敬。我们宣认将有义人和不义之人的复活，以及报应。我们宣认基督不仅是一个纯粹的人，而是天主圣言，也是天主与人之间的中保，父的大司祭；我们不与犹太人一同受割损，因为我们知道\BibleQuote{那应得存留产业的那一位}\BibleRef{创 49:10}已经来了，为祂的缘故各族类得以区分——\BibleQuote{外邦人的期望}，耶稣基督，由犹大而生，出于枝条的儿子，叶瑟的花，政权担在祂的肩上。\par

% structure: p0027 source=h3 target=subsubsection reason=relative-heading-rank; base=h2

\subsubsection{对那宣认基督却渴望犹太化者}

第十二章。但因在当时，这异端似乎更有能力诱骗人，整个教会面临危险，我们十二位宗徒一同聚集在\Place{耶路撒冷}（因为\Person{玛弟亚}被选为宗徒，取代那负卖者的位置，继承了\Person{犹达斯}的职分；正如经上记载：\Quote{他的职分让别人承继}）。\newline{}ewline{}我们与主的兄弟\Person{雅各}一同商议当行何事；他认为和长老们应当向民众宣讲教义的言语。因为也有一些人从\Place{犹太}下到\Place{安提约基雅}，教导那里的弟兄们说：\Quote{你们若不按\Person{梅瑟}的规矩受割损，并遵行他订立的其他规条，就不能得救。}\BibleRef{宗 15:1}\newline{}ewline{}当下起了不小的争论和辩论，\Place{安提约基雅}的弟兄们得知我们都在讨论这个问题，就派遣他们中间那些在圣经上忠实且有见识的人来我们这里询问此事。他们到了\Place{耶路撒冷}，向我们报告了\Place{安提约基雅}教会中出现的疑问——就是说，有些人认为应当受割损，并遵守其他的洁礼。当众人各执一词时，我\Person{伯多禄}站起来对他们说：\newline{}ewline{} \Quote{诸位仁人弟兄，你们知道，早在古代，\Person{天主}已在你们中间拣选了我，使外邦人从我口中听到福音的言语而相信；那洞察人心的\Person{天主}为他们作了证。\BibleRef{宗 15:7}因为有一次，一位上主的天使曾向\Person{科尔乃略}显现，\EditorialNote{宗徒大事录第十章}他是罗马军队的百夫长，对他谈到我，要他派人去请我，听我从口中传的生命之言。于是他派人从\Place{若培}到\Place{凯撒勒雅（斯特拉托）}请我；我正准备去他那里时，想先吃些东西。他们预备的时候，我在楼上祈祷；看见天开了，有一个像一块大布的东西，四角系着，缒到地上，里面各种四足走兽、地上的爬虫和天上的飞鸟。有声音从天上传给我说：伯多禄，起来，宰了吃！我说：主啊，绝对不行！因为我从来没有吃过任何凡俗或不洁之物。第二次有声音说：天主所洁净的，你不可当作俗物。这样一连三次，那器皿就被收回天上去了。但我正猜疑这异象是什么意思，圣神对我说：看，有人找你；起来，跟他们下去，不要疑惑，因为是我派他们来的。那些人正是从百夫长那里来的，于是我领悟了经上记载的主的话：\InnerQuote{凡呼求上主之名的，必获得救恩。}又说：\InnerQuote{大地四极之民都要记起，归向上主；万民的家族都要朝拜祂；因为国度属于上主，祂是万民的治理者。}我见处处都有关于外邦人蒙召的言词，就起来同他们去，进了那人的家。当我传讲圣言时，圣神降在他和同他一起的人身上，正如起初降临在我们身上一样；祂在我们与他们之间没有分别，因信德洁净了他们的心。我明白了天主是不看情面的；凡在敬畏祂、行义的人，无论哪一国，都蒙祂悦纳。连那些信主的受割损的人，也对此惊讶。那么，你们为什么试探天主，把难以负的轭放在门徒的颈上，连我们和我们的祖先都不能承受？但我们相信，借着主的恩宠，我们将如他们一样得救。因为主解开了我们的捆绑，使我们的担子轻松，以祂的仁慈解除了我们沉重的轭。}\newline{}ewline{}我说这些话时，全体群众静默不语。\newline{}ewline{}但主的兄弟\Person{雅各}回答说：\newline{}ewline{} \Quote{诸位仁人弟兄，请听我说！\Person{西默盎}已经述说了天主当初如何眷顾外邦人，从他们中选取一个民族归于自己的名下。先知的话也与此相符，正如经上记载：\InnerQuote{此后我要回来，重建\Person{达味}倒塌的帐幕，修复其废墟，重新树立它，使其余的人寻求上主，所有被称为我名下的民族——上主说，做这些事的天主。}天主的一切作为从创世以来都是祂所知道的。因此我的意见是：不要烦扰那些从外邦人归向天主的人，但要吩咐他们远避外邦人的污秽、祭偶像之物、血、勒死的牲畜和奸淫；这些律法在律法之前，在自然律之下，已赐给古时的先祖，如\Person{厄诺士}、\Person{哈诺客}、\Person{诺厄}、\Person{默基瑟德}、\Person{约伯}以及任何其他类似的人。}\newline{}ewline{}于是我们宗徒、\Person{雅各}主教和长老们，同整个教会，认为应当从我们中间选派几个人，同\Person{巴尔纳伯}和\Person{塔尔索的保禄}（外邦人的宗徒）以及号称\Person{巴尔撒巴的犹达}和\Person{息拉}——弟兄们中的首领——差往你们那里，并亲笔写信如下：\newline{}ewline{} \Quote{宗徒们、长老们和弟兄们，给在\Place{安提约基雅}、\Place{叙利亚}和\Place{基里基雅}的外邦人中的弟兄们请安。我们听说有些人从我们这里出去，用言语搅扰你们，颠覆你们的心灵，其实我们并没有吩咐他们。所以我们同心合意地决定，选派几位我们亲爱的人\Person{巴尔纳伯}和\Person{保禄}同你们一起，他们是曾为我们主\Person{耶稣基督}而冒生命危险的人，正是你们派他们到我们这里来的。我们也派\Person{犹达}和\Person{息拉}同去，他们也会亲口说明同样的事。因为圣神和我们决定，不把别的重担加在你们身上，只有几件必要的事：就是禁戒祭偶像之物、血、勒死的牲畜和奸淫；你们若保守自己远离这些事，就必行得好。愿你们平安。}\newline{}ewline{}我们随即送出了这封信；但我们在\Place{耶路撒冷}又住了多日，一同商议为公共益处，整顿一切事理。\par

% structure: p0029 source=h3 target=subsubsection reason=relative-heading-rank; base=h2

\subsubsection{我们必须与异端分离。}

十三、但是过了很久，我们探访了弟兄们，用虔诚的话坚固了他们，并嘱咐他们躲避那些以基督和梅瑟的名义反对基督和梅瑟，披着羊皮隐藏狼性的人。因为这些人就是假基督、假先知、假宗徒，是欺骗者和败坏者，是狐狸的碎片，是葡萄园中草药的毁灭者：\Quote{因此，许多人的爱会冷淡。但坚忍到底的，必要得救。} \BibleRef{玛 24:12} 关于他们，主为保护我们，曾宣告说：\Quote{将有披着羊皮的人到你们这里来，里面却是凶暴的狼。你们凭他们的果实可以认出他们；要提防他们。因为将有假基督和假先知起来，欺骗许多人。}\par

% structure: p0031 source=h3 target=subsubsection reason=relative-heading-rank; base=h2

\subsubsection{哪些人是天主教教义的宣讲者，以及他们给了什么诫命？}

十四、为此，我们如今聚集在一处的——\Person{伯多禄}和\Person{安德肋}；\Person{载伯德}的儿子\Person{雅各伯}和\Person{若望}；\Person{斐理伯}和\Person{巴尔多禄茂}；\Person{多默}和\Person{玛窦}；\Person{阿尔斐}的儿子\Person{雅各伯}，和号称\Person{塔德奥}的\Person{勒巴}；以及热诚者\Person{西满}，和代替犹达斯与我们同数的\Person{玛弟亚}\BibleRef{玛 10:2}，以及主的兄弟、耶路撒冷的主教\Person{雅各伯}，以及外邦人的导师、蒙选的器皿\Person{保禄}——全都聚在一起，为了坚固你们，将这份天主教教义写给你们，你们被委托管理普世教会。其中我们向你们声明：只有一位全能的天主，除祂以外没有别的，你们必须唯独借着我们的主耶稣基督，在圣神内敬拜并朝拜祂；你们要使用圣经、法律和先知；要孝敬父母；要避免一切不法行为；要相信复活与审判，并期待报应；要感恩地使用祂的一切受造物，作为天主的化工，它们本身并无邪恶；要合法地结婚，因为这样的婚姻无可指摘。因为\Quote{女人是由主配合给男人；}\BibleRef{箴 19:14} 主又说：\Quote{那从起初创造的，造了他们一男一女，并且说：为此，人要离开父亲和母亲，依附自己的妻子，两人成为一体。}\BibleRef{玛 19:4} 也不可认为婚后可以休掉无可指摘的妻子。因为祂说：\Quote{你要谨慎你的心灵，不可离弃你青年时的妻子，因为她是你的终身伴侣，是你灵气的存留。我，没有别人，造了她。} 因为主说：\Quote{天主所结合的，人不可拆散。} 因为妻子是生命的伴侣，由天主将两个合为一体。但若将已合一的分开，就是反对天主的创造，与祂的眷顾为敌。同样，保留被玷污的妻子是违反自然法律，因为\Quote{保留淫妇的人是愚蠢而不虔敬的。} 因为祂说：\Quote{把她从你身上砍掉；} 因为她不是帮助，而是网罗，使她的心偏向别人。也不要在肉体上受割损，但心中由圣神施行的割损，为信者就足够了；因为祂说：\Quote{你要向你的天主受割损，要在你心包的包皮上受割损。}\par

% structure: p0033 source=h3 target=subsubsection reason=relative-heading-rank; base=h2

\subsubsection{我们不应再行圣洗，也不应接受由不虔敬者所施的圣洗，那不是圣洗，而是污秽。}

十五、同样，你们应满足于唯一的圣洗，即归入主死亡的圣洗；不是由邪恶的异端者所施的圣洗，而是由无可指摘的司铎所施的，\Quote{因父及子及圣神之名} \EditorialNote{玛窦福音28:19}；并且你们不可接受来自不虔敬者的圣洗，也不可因第二次而废除由虔敬者所行的。因为正如只有一位天主，一位基督，一位护慰者，和主在肉身上的一次死亡，同样，那归于祂的圣洗也只有一个。但是那些从不虔敬者接受污秽圣洗的人，将与他们意见相同。因为他们不是司铎。因为天主对他们说：\Quote{因为你抛弃了知识，我也必抛弃你，使你不再作我的司铎。} 的确，那些由他们受洗的人不是被入门，而是被玷污，没有领受罪过的赦免，却领受了不虔敬的束缚。此外，那些试图给已入门者施洗的人，是重新钉死主，再次杀害祂，嘲笑神圣事物，亵渎圣事，侮辱圣神，把基督的宝血当作寻常的血，不虔敬地对待派遣者、受难者和见证者。不，那出于轻蔑而不愿受洗的人，将被定罪为不信者，并被斥为忘恩负义和愚昧。因为主说：\Quote{人除非由水和圣神而生，不能进天主的国。} \BibleRef{若 3:5} 又说：\Quote{信而受洗的必要得救；但不信的必被判罪。} \BibleRef{谷 16:16} 但那些说\Quote{我要在死时受洗，免得我犯罪玷污我的圣洗}的人，是不认识天主，也忘记了自己的本性。因为\Quote{不要迟延归向上主，因为你不知道明天将带来什么。} 你们也要给你们的婴孩施洗，并在天主的训诲和告诫中养育他们。因为祂说：\Quote{让小孩子到我跟前来，不要阻止他们。} \BibleRef{玛 19:14}\par

% structure: p0035 source=h3 target=subsubsection reason=relative-heading-rank; base=h2

\subsubsection{关于带有假题名的书籍。}

XVI. 我们已将这一切寄给你们，好让你们知道我们的意见如何；并且使你们不致接受那些以我们名义流传、却由不虔者所写的书。因为你们不应留意宗徒的名号，而应留意事物的本质及他们确定的见解。因为我们知道，西满和克娄俾及其追随者曾以基督及其门徒的名义编撰了有毒的书籍，并四处传播，以欺骗你们这些爱基督和我们这些基督仆役的人。在古人中，也有人伪托梅瑟、哈诺客、亚当、依撒意亚、达味、厄里亚以及三位圣祖之名编写伪经，这些书籍有害且与真理相悖。那些邪恶的异端者如今也做了同样的事，他们毁谤受造界、婚姻、天主眷顾、生育子女、法律和先知；并题写某些野蛮的名号，按其想法是天使的名号，但实则是魔鬼的名号，是魔鬼向他们提示的。要避开他们的道理，以免你们与那些为诱惑和败坏主耶稣的忠信且无可指摘的门徒而编写此类书籍的人同受惩罚。\par

% structure: p0037 source=h3 target=subsubsection reason=relative-heading-rank; base=h2

\subsubsection{\Emph{关于神职人员的婚姻训诫}}

XVII. 我们已说过，主教、司铎和执事在被任命时，必须只结婚一次，无论其妻子健在或已故；他们若在领受圣秩时未婚，此后不得结婚；若当时已婚，则不得再婚，而应满足于他们在领受圣秩时已有的妻子。我们也规定辅礼人员、唱经员、读经员和守门者只能结婚一次。但若他们在结婚之前进入圣职，我们允许他们结婚，如果他们有此意向，以免他们犯罪而受罚。我们不允许任何圣职人员娶娼妓、奴婢、寡妇或被休弃者为妻，正如法律所说\BibleRef{肋 21:7}。女执事应为纯洁的童贞女，或至少是只结过一次婚的寡妇，忠信且受尊敬。\par

% structure: p0039 source=h3 target=subsubsection reason=relative-heading-rank; base=h2

\subsubsection{\Emph{劝勉避免与不虔的异端者共融}}

十八、接纳忏悔者，因为这是天主在基督内的旨意。以宗教的基本要素教导慕道者，然后为他们付洗。躲避那些不敬神的异端者，他们已无可悔改，将他们与信友分离，逐出天主的教会，并嘱咐信友完全远离他们，既不在讲道中也不在祈祷中与他们来往：因为这些人就是教会的敌人，为教会设下陷阱；他们败坏羊群，玷污基督的产业，不过是智慧的冒充者，最卑鄙的人；关于他们，智者撒罗满曾说：\Quote{作恶者假装行善。}因为他说：\Quote{有一条路，人以为正，但终局却指向地狱的深处。}\BibleRef{箴 14:12}这些人就是主以苦涩和严厉表明祂心意的那等人，说：\Quote{他们是假基督、假教师；}\BibleRef{玛 24:24}他们亵渎了恩宠之神，辜负了他们在圣洗恩宠之后从祂所领受的恩赐，\Quote{他们将不得赦免，无论在今生还是来世；}\BibleRef{玛 12:32}他们比犹太人更邪恶，比外邦人更不敬神；他们亵渎万有的天主，践踏祂的圣子，辜负圣神的教导；他们否认天主的话，或假冒伪善地接受，以侮辱天主，欺骗那些接近他们的人；他们滥用圣经，至于正义，他们根本不知为何物；他们败坏天主的教会，如同\Quote{小狐狸毁坏葡萄园；}我们劝你们躲避他们，免得为你们的灵魂设下陷阱。\Quote{与智慧人同行，必得智慧；与愚昧人往来，必被认识。}因为我们既不可与盗贼同跑，也不可与奸夫同分；因为圣达味说：\Quote{主，我恨那恨祢的人，我因祢的仇敌而枯槁。我以完全的恨恨他们，他们成了我的仇敌。}天主藉约纳先知责备约沙法特与阿哈布的友谊，以及他与阿哈布和阿哈齐雅的结盟：\Quote{你与罪人交朋友吗？或者你帮助主所恨的人吗？}\Quote{因此主的忿怒会忽然临到你，但你的心在主面前是完美的。因此主宽恕了你；然而你的工作被粉碎，你的船只也被击碎。}所以，要远离他们的交往，与他们的友谊断绝。因为关于他们，先知曾宣告说：\Quote{与不敬虔的人一同欢乐是不合法的，}主说。因为他们是隐藏的豺狼、不能吠的哑巴狗，目前为数虽少，但随着时间推移，当世界终结临近时，他们会更多、更麻烦，主曾论及他们说：\Quote{人子来临时，能在世上找到信德吗？}\EditorialNote{路 18:8}又说：\Quote{由于罪恶的增加，许多人的爱德会冷淡；}又说：\Quote{将有假基督、假先知出现，在天上显征兆，以致如果可能，连被选者也要被迷惑。}愿天主藉着我们的希望耶稣基督，救我们脱离他们的欺骗。因为我们自己周游列国，坚固各教会，以许多劝勉和医治的话治愈了一些人，当他们已确实走上死亡之路时，我们又将他们挽回。至于那些不可救药的，我们将其从羊群中逐出，免得那些患了疥癣的小羊被传染，而能在主天主面前保持纯洁无玷、健全无污。我们在各城各地、普世各处都如此行，并将这大公的道理，作为对信天主者的纪念或确认，值得而正义地留给你们主教及其他司铎；我们已藉同工克肋孟——我们在主内最忠信和亲密的儿子——送去，连同巴尔纳伯、我们最亲爱的儿子弟茂德以及真正的马尔谷；与此一同，我们也推荐弟铎、路加、雅松、路基约和索西帕托给你们。\BibleRef{罗 16:21}\par

% structure: p0041 source=h2 target=subsection reason=relative-heading-rank; base=h2

\subsection{第四节 论法律}

藉此我们也劝你们在主内远离旧日的言行、虚浮的束缚、分离、规条、食物的分别和每日的洗涤；因为\Quote{旧的已成过去，看，一切都成了新的。}\BibleRef{格后 5:17}\par

% structure: p0043 source=h3 target=subsubsection reason=relative-heading-rank; base=h2

\subsubsection{致诽谤法律者}

十九、因为你们已藉耶稣基督认识了天主及祂的全部安排，即从起初祂就赐下一条明确的律法来辅助自然律，这律法是纯洁、拯救人且神圣的，其上刻有祂的名号，完美无缺，永不失效，由十诫所成全，无玷，转化人灵；当希伯来人忘记时，祂藉玛拉基亚先知提醒他们，说：\Quote{你们当纪念天主的仆人梅瑟的律法，他吩咐你们诫命和典章。}这律法如此神圣正义，以至于我们的救主一次治愈了一个癞病人，后来又治愈了九个，对第一个说：\Quote{去，把你自己给司祭看，并献上梅瑟所吩咐的礼物，作为对你们的证据；}之后对那九个说：\Quote{去，把你们自己给司祭看。}\BibleRef{路 17:14}因为祂从未废掉律法，如同西满所妄称的，而是成全了它；因为祂说：\Quote{即使天地过去了，一撇或一画也决不会从律法上过去，直到一切完成。}因为祂说：\Quote{我来不是为废掉律法，而是为成全。}因为梅瑟本人——他同时是立法者、大司祭、先知和君王——以及众先知的热心追随者厄里亚，在我主显圣容时都在山上，并为祂的降生成人和苦难作了见证，作为基督的密友，而非仇敌或陌生人。由此可见，律法是良善而神圣的，先知们也是如此。\par

% structure: p0045 source=h3 target=subsubsection reason=relative-heading-rank; base=h2

\subsubsection{论何者为自然律，何者为后来引入的法律，以及为何引入}

XX. 法律就是\OriginalTerm{十诫}{decalogue}，是主以可听见的声音向人们颁布的，那时人民尚未铸造那头代表埃及阿彼斯（\OriginalTerm{Apis}{Apis}）的牛犊。这法律是正义的，因此被称为法律，因为审判乃是依照\OriginalTerm{自然法}{law of nature}而作出，而西满的追随者却滥用这法律，以为他们不会因此受审，从而逃脱惩罚。这法律是善的、圣的，并且不强制要求具体的事。因为祂说：\Quote{若你们愿为我筑一座祭台，应用土筑。}它并没有说\Quote{筑一座}，而是说\Quote{若你们愿筑}。它没有施加必须性，而是任凭他们自由选择。因为天主不需要祭献，祂的本性超越一切需要。但祂知道，正如古时，天主所爱的\Person{亚伯}、\Person{诺厄}、\Person{亚巴郎}以及后来的先人，在没有被要求的情况下，只是出于自然法的自发感动，以感恩之心向天主献祭；同样，祂现在也允许希伯来人，不是命令，而是如果他们愿意，就允许他们；如果他们怀着正确的意向献祭，就表示悦纳他们的祭献。因此祂说：\Quote{若你们愿意奉献，不要奉献给我，好像我是需要它似的，因为我什么都不需要；因为世界是我的，其中的一切都是我的。}但当这人民忘记了这事，反而呼求一只牛犊为神，取代了真天主，并把这牛犊当作他们出埃及的原因，说：\Quote{以色列，这些就是你的神，领你出埃及地的神。}这些人行了恶事，拜了那\Quote{吃草的牛犊的像}，并否认了天主——那位曾在他们困苦中藉\Person{梅瑟}访问他们，用祂的手和杖行神迹，用十灾击打埃及人；那位曾将红海的水分为两半，引导他们在水中行走，如同在旱地上；那位曾淹没了他们的仇敌和埋伏者；那位曾在\Place{玛辣}使苦泉变甜；那位曾从坚硬的磐石中流出水来，直到他们满足；那位曾用云柱为他们遮蔽炎热，又用火柱照亮并引导他们，使他们知道当行的路；那位曾从天上赐下玛纳，从海中赐下鹌鹑为肉；那位曾在山上赐给他们法律，并屈尊让他们听到祂的声音——他们竟否认了祂，并对\Person{亚郎}说：\Quote{请为我们造神，在我们前面引路。}于是他们铸了一只牛犊，向偶像献祭。那时，天主因他们忘恩负义而发怒，用无法解脱的枷锁、令人痛苦的负担和沉重的轭束缚了他们，不再说\Quote{若你们愿筑}，而是说\Quote{筑一座祭台}，并不断献祭；因为你们是健忘和忘恩负义的。所以你们要不断地献燔祭，好叫你们记念我。既然你们邪恶地滥用了你们的自由，我就为将来对你们施加必须性，并命令你们戒除某些肉类；我规定你们洁净与不洁净之物的区别，尽管每一种受造物都是善的，因为是出于我手；我指定你们各种分离、洁净、频繁的洗涤和洒水、各种净化以及各种安息时间；如果你们忽略其中任何一条，我就规定对不服从者恰当的惩罚，好让你们被自己的轭所压迫和刺痛，从而离开多神论的错谬，抛弃那句\Quote{以色列，这些就是你的神}，而记住那句\Quote{以色列，请听：上主你的天主是唯一的主}，并重回我铭刻在众人本性中的法律，\Quote{在天地之间只有一位天主，并且要全心全意全力全意爱祂}，除祂以外不怕任何神，也不让其他神的名字进入你们的心，也不让你们的舌头从口中说出。祂因他们心硬而束缚他们，好让他们藉着献祭、安息、自洁以及其他类似的规定，得以认识那为他们制定这些事的天主。\par

% structure: p0047 source=h3 target=subsubsection reason=relative-heading-rank; base=h2

\subsubsection{我们这些信基督的人是在恩宠之下，不在那附加法律的奴役之下。}

XXI. \Quote{但你们的眼睛是有福的，因为看得见；你们的耳朵是有福的，因为听得见。}\BibleRef{玛 13:16}我说，你们这些相信唯一天主的人，不是出于必须，而是出于健全的理智，听从了那召叫你们的。因为你们已从捆绑中得释放，从奴役中得自由。因为祂说：\Quote{我不再称你们为仆人，而是朋友，因为我从父那里所听见的一切，都告诉了你们。}\BibleRef{若 15:15}因为对于那些不愿看、不愿听的人——不是因为缺乏这些感官，而是因为他们极度的邪恶——\Quote{我赐给他们的律例并不美好，按照这些律例，他们不能生活。}\BibleRef{则 20:25}这些律例被视为不美好，如同烧伤和刀剑，又如病人眼中的药物和手术，被认为是仇敌；而且由于他们的顽固，这些律例无法遵守；因此它们因不被服从而给人们带来了死亡。\par

% structure: p0049 source=h3 target=subsubsection reason=relative-heading-rank; base=h2

\subsubsection{那为祭献而设的法律是附加的，基督来临时便把它废除了。}

XXII. 所以，你们这些从诅咒中被解救出来的人是有福的。因为天主子基督，祂的来临确认并完成了法律，但除去了附加的规条——虽非全部，至少是那些更为沉重的；祂确认了前者，废除了后者，并再次恢复了人的自由意志，不再使人受制于暂时的死亡，而是按照另一种律法赐予人法则。因此祂说：\Quote{若有人愿意跟随我，让他来。}\BibleRef{玛 16:24}又说：\Quote{你们也愿意走吗？}\BibleRef{若 6:67}此外，在祂来临之前，当人民常常献祭时，祂拒绝了他们的祭品，因为他们得罪了祂，以为可藉祭品而非悔改来平息祂的愤怒。因为祂这样说：\Quote{你们为什么从舍巴给我带来乳香，从远方带来肉桂？你们的全燔祭不蒙悦纳，你们的祭品我也不喜欢。}此后又说：\Quote{你们把全燔祭和祭品集合起来，吃肉吧！因为我领你们出埃及地的时候，并没有吩咐你们关于全燔祭和祭品的事。}\BibleRef{耶 7:21}祂又藉依撒意亚说：\Quote{你们向我献上那么多祭品，有什么益处呢？上主说：公羊的全燔祭我已饱足，肥犊的脂油我不再接受，公牛和山羊的血我也不要。你们不要再来见我面前；谁向你们要求这些事？不要再践踏我的庭院。如果你们给我送来细面，也是徒然；乳香是我所憎恶的；你们的月朔、安息日和盛大的日子，我无法忍受；你们的斋戒、休息和节期，我的灵魂憎恨它们；我对此已过度饱足。}祂又藉另一先知说：\Quote{离开我！你们赞歌的声音，你们乐器的诗篇，我不愿听。}撒慕尔对撒乌耳说，当他想要献祭时：\Quote{听从胜于祭献，服从胜于公羊的脂油。因为，看，上主喜爱听从祂胜于祭献。}祂又藉达味说：\Quote{我不从你家中取牛犊，也不从你羊圈中取公山羊。如果我饥饿，我不会告诉你；因为世界和其中的一切都是我的。难道我要吃公牛的肉，喝山羊的血吗？向天主献上赞颂的祭品，并向至高者偿还你的誓愿。}同样，在全部圣经中，祂因他们得罪祂而拒绝他们的祭品。因为\Quote{不虔敬者的祭品为上主所憎恶，因为他们以非法的方式献上。}又说：\Quote{他们的祭品对他们如同哀悼之饼；凡吃它们的必被玷污。}因此，在祂来临之前，祂寻求\Quote{洁净的心和痛悔的精神}胜过祭品，那么当祂来临后，祂更会废除那些祭品——我指那些流血的祭品。然而祂废除它们的方式是先成全它们。因为祂也受了割损，被洒，献了祭品和全燔祭，并遵行了其余的风俗。那位立法者自己成了法律的圆满者；祂不是取消了自然律，而是删除了后来附加的规条，尽管也不是全部。\par

% structure: p0051 source=h3 target=subsubsection reason=relative-heading-rank; base=h2

\subsubsection{基督如何成为法律的圆满者，以及祂终结、改变或转移了哪些部分}

XXIII. 因为祂并没有废除\OriginalTerm{自然法}{law of nature}，反而予以确认。因为那在法律中说过\BibleQuote{上主你的天主是唯一的主}的，同一位置也在《福音》中说\BibleQuote{使他们认识你，唯一的真天主。}\BibleRef{若 17:3} 那说过\BibleQuote{你应爱人如己}\BibleRef{肋 19:18}的，在《福音》中更新同一诫命，说\BibleQuote{我给你们一条新命令，你们要彼此相爱。}\BibleRef{若 13:34} 从前禁止杀人的，如今禁止无故发怒。\BibleRef{玛 5:22} 从前禁止奸淫的，如今禁止一切不法的情欲。从前禁止偷窃的，如今称那以自己的劳动供应缺乏之人为最有福的。\BibleRef{宗 20:35} 从前禁止仇恨的，如今称那爱仇敌的人为有福的。\BibleRef{玛 5:7} 从前禁止报复的，如今命令忍耐；\BibleRef{玛 5:43} 并非因为正当的报复是不义，而是因为忍耐更为卓越。祂也没有立法根除我们的自然情欲，只是禁止其过度。\BibleRef{玛 5:38} 那曾命令孝敬父母的，自己却顺服了他们。\BibleRef{路 2:51} 那曾命令守安息日，藉安息日休息以默想法律的，如今命令我们天天思考创造之法和眷顾，并感谢天主。祂在亲自成全了割损之后，便废除了它。因为祂就是\BibleQuote{继承产业、万民所期待的那一位。}\BibleRef{创 49:10} 那为正当起誓立法、禁止假誓的，如今命令我们完全不可起誓。\BibleRef{玛 5:33} 祂以多种方式改变了圣洗、祭献、司祭职以及局限于一处地方的敬礼：因为祂用一次——即归于祂死亡的圣洗——代替了每日的圣洗。祂指定从各民族中选出最优秀者领受司祭职，代替了一个支派；并且不是检验他们身体上的残疾，而是检验他们的信仰和生活。祂用那理性的、无血的奥迹性的祂身体和血的祭献，代替了\OriginalTerm{血腥祭献}{bloody sacrifice}，这祭献藉象征举行，以代表主的死亡。祂命令并指定应在祂统治范围内的每一处，从日出到日落受光荣，代替了局限于一处地方的敬礼。因此，祂并没有从我们身上取去法律，而是取去了\OriginalTerm{桎梏}{bonds}。因为关于法律，梅瑟说：\BibleQuote{你应坐在家中、起来行走在路上时，默想我所吩咐你的话。}达味也说：\BibleQuote{他喜爱上主的法律，昼夜默思他的法律。}因为祂愿意我们在一切事上服从祂的法律，而不成为违犯者。因为祂说：\BibleQuote{在道路上无瑕疵的、遵行上主法律的人是有福的。那寻求祂法度的人是有福的，他们全心寻求祂。}又说：\BibleQuote{以色列啊，我们有福了，因为天主所喜悦的事，我们都知道了。}\BibleRef{巴 4:4} 主说：\BibleQuote{你们既然知道这些事，若实行便是有福的。}\BibleRef{若 13:17}\par

% structure: p0053 source=h3 target=subsubsection reason=relative-heading-rank; base=h2

\subsubsection{主愿意\OriginalTerm{正义法}{law of righteousness}由罗马人显扬。}

XXIV. 祂也不愿意正义法只由我们显扬，而愿意它藉罗马人出现并照耀。因为这些罗马人信主，放弃了他们的\OriginalTerm{多神信仰}{polytheism}和不义，接纳善人，惩罚恶人。但他们使犹太人缴\OriginalTerm{贡赋}{tribute}，不许他们使用自己的\OriginalTerm{规条}{ordinances}。\par

% structure: p0055 source=h3 target=subsubsection reason=relative-heading-rank; base=h2

\subsubsection{天主如何因他们对基督的\OriginalTerm{不敬}{impiety}，使犹太人成为俘虏，并使他们受\OriginalTerm{贡赋}{tribute}。}

XXV. 确实，他们自愿把奴役加在自己身上，因为说了\BibleQuote{我们没有别的君王，只有凯撒。}\BibleRef{若 19:15} 和\BibleQuote{如果我们不杀死基督，众人都会信从祂，罗马人就会来夺去我们的地方和民族。}\BibleRef{若 11:48} 于是他们无意中说了预言。因为照样，列邦信从了祂，他们自己被罗马人剥夺了权力和合法的敬礼；他们被禁止随意杀人和随意献祭。因此他们受到诅咒，因为他们不能履行所命之事。因为祂说：\BibleQuote{凡不持续遵行写在法律书上的一切事的人，是可诅咒的。}如今在他们的分散中，在外邦人中，他们不可能履行法律中的一切。因为神圣的梅瑟禁止在耶路撒冷以外筑祭台，也禁止在犹太地边界以外宣读法律。所以，让我们跟随基督，好承受祂的祝福。让我们藉福音遵行法律和先知。让我们躲避多神崇拜者、杀害基督者、杀害先知者，以及邪恶和无神的\OriginalTerm{无神异端者}{atheistical heretics}。让我们服从基督如同我们的君王，祂有权柄改变某些\OriginalTerm{宪章}{constitutions}，并作为\OriginalTerm{立法者}{legislator}有智慧在不同的情况下制定新的宪章；然而，到处都不可更改地保存\OriginalTerm{自然法}{laws of nature}。\par

% structure: p0057 source=h2 target=subsection reason=relative-heading-rank; base=h2

\subsection{第五节 宗徒的教训：反对犹太人和外邦人的\OriginalTerm{迷信}{superstitions}，尤其是关于婚姻与丧葬}

% structure: p0058 source=h3 target=subsubsection reason=relative-heading-rank; base=h2

\subsubsection{我们应当躲避异端者，如同躲避\OriginalTerm{灵魂的败坏者}{corrupters of souls}。}

XXVI. 所以，诸位主教，以及你们平信徒，要躲避所有妄用法律和先知书的异端者。因为他们是全能天主的仇敌，违抗祂，不承认基督是天主子。他们也否认祂按肉身的出生；他们以十字架为耻；他们辱骂祂的受难和死亡；他们不知祂的复活；他们否认祂在万世之前的出生。不仅如此，其中一些人又以另一种方式不敬，以为主不过是凡人，认为祂由灵魂和肉体组成。另一些人则认为耶稣自己就是万有之上的天主，将祂作为自己的圣父来荣耀，认为祂既是圣子又是护慰者；还有什么教义比这更可憎恶呢？还有一些人禁食某些食物，认为生育子女的婚姻是邪恶的，是魔鬼的计谋；他们自己不敬虔，因自己的邪恶不愿从死者中复活。因此他们也嘲笑复活，说：\Quote{我们是圣洁的人，不愿吃喝；}他们幻想自己将从死者中复活，成为没有肉身的魔鬼，永远在永火中被定罪。所以，要逃离他们，免得你们在他们的不敬中与他们一同灭亡。\par

% structure: p0060 source=h3 target=subsubsection reason=relative-heading-rank; base=h2

\subsubsection{论某些犹太人和外邦人的惯例。}

XXVII. 如果有人遵守关于遗精和夜遗以及合法交合的犹太惯例和规定，\EditorialNote{肋未纪第十五章} 让他们告诉我们，在他们经历此类事情的那些时辰或日子，他们是否不祈祷、不触摸圣经、不领圣体？如果他们承认是这样，那么显然他们缺少始终与信者同在的圣神。因为关于圣者，撒罗满说：\BibleQuote{这样，无论你睡时，它必守护你；你醒时，它必与你交谈。}\BibleRef{箴 6:22} 因为如果你，女人，认为在你月经的七天中，你缺少圣神，那么如果你突然死去，你将缺少圣神而离开，并且对天主没有确定的希望；否则你必须认为圣神总是与你不可分离，因为它不在一个地方。但你在祈祷、圣体和圣神的临在方面有需要，因为你在这件事上并没有犯任何过错。因为合法交合、生育、月经、夜遗都不能玷污人的本性，或将圣神从他身上分离。只有不敬和非法行为才能做到。因为圣神总是与那些拥有它的人同在，只要他们相称；那些失去它的人，它使他们孤独，暴露于恶神。现在每个人要么被圣神充满，要么被不洁之神充满；人不可能避免其中之一，除非他们能够接受相反的神。因为护慰者憎恨一切谎言，魔鬼憎恨一切真理。但每一个按照真理受洗的人，与魔鬼之神分离，并在圣神之下；只要他行善，圣神就与他同在，以智慧和明悟充满他，不允许恶神靠近他，而是看顾他的脚步。所以，你，女人，如果你说，在你月经的日子里你缺少圣神，那么你那时就被不洁者充满；因为忽视祈祷和读经，你会邀请他到你这里来，即使他不愿意。因为这个神，在一切之中，喜爱忘恩、懒惰、粗心、昏睡的人，因为他自己因忘恩而心志扭曲，并被天主剥夺了尊位；他宁愿成为魔鬼，也不愿做总领天使。所以，女人，要避免这些虚妄的话，要时常记念创造你的天主，并向祂祈祷。因为祂是你的主，也是宇宙的主；要默想祂的法律，而不遵守任何这样的规条，如月经、合法交合、生育、流产、身体残缺；因为这些规条是愚妄之人的虚空发明，这些发明毫无意义。人的埋葬、死人的骨头、坟墓、某种特别的食物、夜遗，都不能玷污人的灵魂；只有对天主的不敬、违命、对邻人的不义，即抢夺、暴力，或任何违反祂公义的事，奸淫或淫乱。所以，亲爱的，要回避并远离这些规条，因为它们是外邦人的做法。因为我们不像他们那样憎恶死人，因为我们希望他将复活。我们也不恨恶合法交合；因为在这种事情上行不敬是他们的习惯。因为男人和妻子的结合，如果是公正的，就符合天主的旨意。\BibleQuote{那起初造人的，是造男造女，并且说：因此，人要离开父母，与妻子连合，二人成为一体。}\BibleRef{玛 19:4} 如果两性的区别是出于天主的旨意，为了生养众多，那么男女的结合也必须符合祂的旨意。\par

% structure: p0062 source=h3 target=subsubsection reason=relative-heading-rank; base=h2

\subsubsection{论男色、奸淫和淫乱。}

二十八．但我们并不这样说那违背本性的混合，或任何不法行为；因为这些都是与天主为敌。因为索多玛的罪是违背本性的，与兽类行淫也是如此。但奸淫和邪淫是违反法律的；其一是不敬，其二是非义，总之，无非是大罪。然而这两种罪在其自身本性上都不是没有惩罚的。因为第一种行径的人企图瓦解世界，并试图使事物的自然进程转变为不自然的；而第二种人——奸淫者——是不义的，因为他们败坏别人的婚姻，将天主所合一的分为二，使子女受到怀疑，并使真正的丈夫陷入他人的罗网。邪淫是败坏自己的肉身，不是为了生育子女，而完全是为了享乐，这标志不节制，而非德行的标记。所有这些都为法律所禁止；因为神谕如此说：\newline{}\BibleQuote{不可跟男人同寝，像跟女人同寝一样。}\newline{}\BibleQuote{因为这样的人是可诅咒的，你们要用石头打死他们；他们行了可憎的事。}\newline{}\BibleQuote{凡与兽类同寝的，要杀死他；他在自己的人民中行了恶事。}\newline{}\BibleQuote{若有人玷污了有夫之妇，要把他们二人处死；他们行了恶事，他们是咎由自取，让他们死。}\newline{}之后又说：\newline{}\BibleQuote{在以色列子民中不可有邪淫者，在以色列女儿中不可有娼妓。你不可将妓女的酬金或狗的价格带到上主你的天主的祭台上。}\newline{}\BibleQuote{因为妓女的酬金所带来的誓愿是不洁的。}这些事法律禁止了，但法律尊重婚姻，并称它为有福的，因为天主祝福了它，祂将男和女结合在一起。聪明的\Person{撒罗满}在某处说：\BibleQuote{妻子是上主为丈夫配就的。}\Person{达味}说：\BibleQuote{你的妻子像一棵茂盛的葡萄树，在你家的四壁之内；你的子女像橄榄枝，环绕着你的桌子。看哪，敬畏上主的人必如此蒙福。}因此\BibleQuote{婚姻是可敬的}，也是体面的，生育子女是纯洁的，因为善的事物中没有恶。因此，自然的洁净在上主面前也不是可憎的，祂规定妇女在三十天内发生此事，为了她们的益处和健康状态，使她们少走动，通常留在家中。不仅如此，甚至在福音中，当那患血漏的妇人触摸主的衣边，希望得到医治时，主并没有向她发怒，也没有责备她；相反，祂治好了她，说：\BibleQuote{你的信德救了你。}当妻子自然洁净出现时，丈夫不要亲近她们，出于对将要生育的子女的考虑；因为法律禁止了这事，因为法律说：\BibleQuote{月经期间，不可亲近你的妻子。}\BibleRef{肋 18:19}而且，当妻子怀孕时，丈夫也不要与她们频繁同房。因为他们这样行不是为了生育子女，而是为了享乐。一个爱天主的人不该是贪爱享乐的人。\par

% structure: p0064 source=h3 target=subsubsection reason=relative-heading-rank; base=h2

\subsubsection{妻子应如何服从自己的丈夫，丈夫应如何爱自己的妻子。}

XXXIX. 你们作妻子的，要服从自己的丈夫，尊重他们，以敬畏和爱服事他们，如同圣\Person{撒辣}尊敬\Person{亚巴郎}一样。因为她不能忍受直呼其名，而称他为主，她说：\BibleQuote{我的主已经老了。}同样，你们作丈夫的，要爱自己的妻子如同自己的肢体，如同生活的伴侣，为生育子女的助手。因为祂说：\BibleQuote{要与你幼年所娶的妻子同乐。愿她的言谈对你如同可爱的母鹿和愉快的小驹；让她独自引导你，时刻与你同在：因为如果你处处被她的友谊环绕，你将在她的陪伴中幸福。}因此要爱她们如同自己的肢体，如同自己的身体；因为经上这样写：\BibleQuote{上主在你和你幼年所娶的妻子之间作了见证；她是你的伴侣，别人并没有造她：她是你的精神所留之物；}以及：\BibleQuote{当留意你的精神，不要离弃你幼年所娶的妻子。}因此，丈夫和妻子在合法的婚姻中同房后，彼此起来，可以无需任何观察而祈祷，无需洗濯便是洁净的。但凡是败坏和玷污别人妻子的人，或是与娼妓行淫被玷污的人，当他从她起来后，即使他在整个大洋和所有河流中洗濯，也不能成为洁净。\par

% structure: p0066 source=h2 target=subsection reason=relative-heading-rank; base=h2

\subsection{第六部分 本书的结论}

% structure: p0067 source=h3 target=subsubsection reason=relative-heading-rank; base=h2

\subsubsection{犹太人和外邦人的习俗：遵守自然洁净和厌恶死尸，但这一切与基督宗教相反。}

XXX. 因此，不要因法律性和自然性的洁净而守任何这类规矩，以为自己被它们玷污。也不要求犹太人的隔绝、不断的洗涤，或触摸尸体后的洁净。但不要这些规矩，在安息所聚集，阅读圣经，并为已安息的殉道者、从创世以来的所有圣人、以及在主内安息的弟兄们歌唱，并献上可接受的圣体圣事，即基督王体之像，无论在你们的教堂还是安息所；在逝者的葬礼上，若他们在基督内是忠信的，就歌唱陪伴他们。因为\BibleQuote{主的圣者之死在主眼中是珍贵的。}又说：\BibleQuote{我的灵魂，回到你的安息，因为主已厚待了你。}又在别处：\BibleQuote{义人的记念受人称赞。}\BibleRef{箴 10:7}并且，\BibleQuote{义人的灵魂在天主手中。}\BibleRef{智 3:1}因为那些信了天主的人，虽然安息了，却没有死。因为我们的救主对撒杜塞人说：\BibleQuote{论到死人复活，你们没有读过所记载的：我是亚巴郎的天主，依撒格的天主，雅各伯的天主？天主不是死人的天主，而是活人的天主；因为众人为祂而活。}\BibleRef{出 3:6}因此，对于与天主同活的人，甚至他们的遗骸也并非没有尊荣。因为连先知厄里叟在安息后，也使一个被叙利亚海盗杀害的死人复活。他的身体接触了厄里叟的骨头，就起来并复活了。这事若不因厄里叟的身体是神圣的，就不会发生。贞洁的若瑟在雅各伯死后拥抱他在床上；梅瑟和农的儿子若苏厄带走了若瑟的遗骸，并不认为那是污秽。因此，你们这些主教和其余的人，若没有这些规矩而触摸逝者，也不该以为自己被玷污。也不要厌恶这些人的遗骸，但要避免这类规矩，因为它们是愚昧的。要以圣洁和贞洁装饰自己，好成为不朽的参与者，天主国的同伴，领受天主的应许，并永远安息，藉着我们的救主耶稣基督。\par

因此，愿那位能开启你们心灵之耳以接受天主的神谕——这神谕由福音和纳匝肋人耶稣基督的教导传给你们；祂在本丢·彼拉多和黑落德手下被钉十字架，死了，又从死人中复活，并要在世界末日带着大能和大光荣再来，使死人复活，终结这个世界，按各人的行为施行报应；祂把自己给了我们作为复活的担保；祂被天主和父的大能在我们眼前升到天上，从死里复活后四十天与我们一同吃喝；祂坐在全能天主威严宝座的右边，在革鲁宾之上；祂曾听到：\BibleQuote{坐在我的右边，等我使你的仇敌作你的脚凳；}那位蒙祝福的斯德望看见祂站在大能的右边，就呼喊说：\BibleQuote{看，我见天开了，人子站在天主的右边，}\BibleRef{宗 7:56}作为所有理性秩序的大司祭——愿崇敬、威严和光荣藉着祂归于全能的天主，从现在直到永远。阿们。\par

\SourceRecord{Translated by James Donaldson.；Ante-Nicene Fathers；Vol. 7.；Edited by Alexander Roberts, James Donaldson, and A. Cleveland Coxe.；Buffalo, NY: Christian Literature Publishing Co.,；1886.；<http://www.newadvent.org/fathers/07156.htm>.}

% structure: p0001 source=h1 target=section reason=page-title

\section{宗徒宪章（第七卷）}

\Emph{论基督徒生活、圣体圣事及进入基督。}\par

% structure: p0003 source=h2 target=subsection reason=relative-heading-rank; base=h2

\subsection{第一节　论两条道路——生命之路与死亡之路}

% structure: p0004 source=h3 target=subsubsection reason=relative-heading-rank; base=h2

\subsubsection{论有两条道路——一条是自然的生命之路，另一条是后来引入的死亡之路；前者来自天主，后者来自错谬，出于仇敌的罗网。}

一、立法者梅瑟对以色列人说：\Quote{看，我已将生命之路和死亡之路摆在你们面前。}又加上：\Quote{选择生命，好使你们活下去。}先知厄里亚也对民众说：\Quote{你们摇摆不定要到几时呢？如果上主是天主，就跟随祂。}\EditorialNote{列上 18:21}主耶稣也公正地说：\Quote{没有人能事奉两个主人：因为他要么恨这一个，而爱那一个，要么依附这一个，而轻视那一个。}\BibleRef{玛 6:24}我们也跟随我们的导师基督，\Quote{祂是众人的救主，尤其是信徒的救主}\BibleRef{弟前 4:10}，我们必须说有两条道路——一条是生命之路，另一条是死亡之路；二者无法相比，因为它们极为不同，更确切地说是完全分离的；生命之路是自然的，而死亡之路是后来引入的——它不符合天主的旨意，而是出于仇敌的罗网。\par

% structure: p0006 source=h3 target=subsubsection reason=relative-heading-rank; base=h2

\subsubsection{主之宪章的道德劝谕，与神圣法律的古代禁令一致。禁止愤怒、怨恨、腐败、奸淫及一切被禁行为。}

二、因此，第一条道路是生命之路；它也是法律所规定的：\Quote{你要全心、全灵爱上主你的天主，祂是独一的天主，除祂以外没有别的神；}并\Quote{要爱人如己。}凡你不愿意别人对你做的，也不要对别人做。\BibleRef{多 4:15}\Quote{咒骂你们的，要祝福他们；凌辱你们的，要为他们祈祷。}\BibleRef{玛 5:44}\Quote{爱你们的仇敌；如果你们只爱那爱你们的人，有什么可酬谢的呢？连外邦人也这样做。}\BibleRef{路 6:32}\Quote{要爱那恨你们的人，你们就没有仇敌。}因为祂说：\Quote{不可恨任何人，连埃及人也不可，连厄东人也不可；因为他们都是天主的化工。}不要躲避恶人本人，而要躲避他们的思想。\Quote{戒避肉体和世俗的情欲。}\Quote{如果有人打你的右脸，把另一面也转给他。}这不是说报复是恶，而是说忍耐更可敬。因为达味说：\Quote{我若向那以恶报我的人有所回报。}\Quote{如果有人强迫你走一里路，你就同他走两里。}并且，\Quote{那愿意控告你，拿你内衣的，连外衣也让他拿去。}\BibleRef{玛 5:40}\Quote{夺你东西的，不要再向他追讨。}\BibleRef{路 6:30}\Quote{凡求你的，就给他；有人向你借贷，不要转背不顾。}\BibleRef{玛 5:42}因为\Quote{义人仁慈，且肯借贷。}因为你们的父愿意你们给众人，祂自己\Quote{使太阳升起，照耀恶人也照耀善人，降雨给义人也给不义的人。}\BibleRef{玛 5:45}因此，用你自己的劳动所得给众人是合理的；因为祂说：\Quote{要以你正义的劳动光荣主}，但要优先照顾圣徒。\Quote{不可杀人；}就是说，不可毁灭一个像你一样的人，因为你拆毁了受造的精品。并非一切杀人都邪恶，只杀害无辜者是邪恶的；但正义的杀戮只保留给掌权者。\Quote{不可奸淫；}因为你把一体分成了两半。\Quote{二人成为一体；}\BibleRef{创 2:24}因为丈夫和妻子在本性、同意、结合、性情和生活行为上是一体的；只是在性别和数目上分开。\Quote{不可猥亵男童；}\BibleRef{肋 18:22}因为这恶是违反本性的，起源于索多玛，因此索多玛被天主降下的火完全烧尽。\EditorialNote{创 19}\Quote{这样的人该受诅咒；全体民众要说：阿们。}\EditorialNote{申 27}\Quote{不可行淫；}因为祂说：\Quote{在以色列子民中不可有淫乱者。}\BibleRef{申 23:17}\Quote{不可偷窃；}因为阿干在耶里哥偷窃以色列的财物，被石头砸死；\EditorialNote{苏 7}革哈齐偷窃并撒谎，得了纳阿曼的癞病；\EditorialNote{列下 5}犹达斯偷窃穷人的钱，出卖了荣耀之主给犹太人\BibleRef{若 12:6}，而后后悔，上吊自杀，肚腹破裂，肠子都流了出来；\BibleRef{玛 27:5}阿纳尼雅和妻子撒斐辣偷窃自己的财物，并\Quote{试探主的神}，立刻因我们同宗徒伯多禄的判决而倒地死亡。\EditorialNote{宗 5}\par

% structure: p0008 source=h3 target=subsubsection reason=relative-heading-rank; base=h2

\subsubsection{禁止邪术、杀婴、伪证和假见证。}

三、不可行邪术。不可行巫术；因为祂说：\Quote{不可容巫女存活。}不可堕胎杀害你的孩子，也不可杀死已受孕的；因为\Quote{凡成形并从天主领受了灵魂的，若被杀，必遭报复，因为是不义地被毁灭。}\Quote{不可贪恋你邻人的财物，如他的妻子、奴仆、牛或田地。}不可发假誓；因为经上说：\Quote{不可发誓。}\BibleRef{玛 5:34}但如果无法避免，就当诚实地发誓；因为\Quote{凡指着祂发誓的，必受称赞。}\Quote{不可作假见证；}因为\Quote{那诬告穷人的，激怒了造他的主。}\BibleRef{箴 14:31}\par

% structure: p0010 source=h3 target=subsubsection reason=relative-heading-rank; base=h2

\subsubsection{禁止毁谤、怒气、欺诈行为、闲话、谎言、贪婪和伪善。}

四、不可说恶言；因祂说：\Quote{不要爱说恶言，免得你被除去。}也不可记仇；因为\Quote{记仇之人的道路是通向死亡。}不可心怀二意，也不可一口两舌；因为\Quote{人的嘴唇是他自己的坚固罗网，}\BibleRef{箴 6:2}并且\Quote{多言的人在地上不会亨通。}你的话语不可虚妄；因为\Quote{你要为每句闲话交账。}\BibleRef{玛 12:36}不可说谎：因祂说：\Quote{你要消灭一切说谎的人。}不可贪婪也不可强夺；因祂说：\Quote{祸哉，那用邪恶的贪婪向邻舍贪婪的人。}\BibleRef{哈 2:9}\par

% structure: p0012 source=h3 target=subsubsection reason=relative-heading-rank; base=h2

\subsubsection{禁止恶意、偏待人、愤怒、怨恨和嫉妒。}

五、你不可假冒为善，免得你的\Quote{分与他们相同。}你不可心肠恶劣，也不可骄傲；因为\Quote{天主拒绝骄傲的人。}\Quote{审判时不可偏袒，因为审判属于主。}\Quote{你不可恨任何人，你应切实劝戒你的兄弟，免得因他而成为有罪；}并且，\Quote{劝戒智慧人，他必爱你。}要躲避一切邪恶及类似邪恶的事；因为祂说：\Quote{远离不义，战栗就不会临近你。}不可轻易发怒，不可心怀怨恨，不可暴躁，不可狂怒，不可鲁莽，免得你遭受\Person{加音}、\Person{撒乌耳}和\Person{约阿布}的命运：因为第一个杀了他的兄弟\Person{亚伯尔}，因为亚伯尔在天主面前被看为比他更优，且亚伯尔的祭献被悦纳；第二个因嫉妒跳舞妇女的赞美而迫害圣\Person{达味}，后者杀了培肋舍特人\Person{哥肋雅}；第三个杀了两位军队将领——以色列的\Person{阿贝乃尔}和犹大的\Person{阿玛撒}。\par

% structure: p0014 source=h3 target=subsubsection reason=relative-heading-rank; base=h2

\subsubsection{关于占卜和邪术。}

六、不可占卜，因为那导向偶像崇拜；因为\Person{撒慕尔}说：\Quote{占卜是罪；}\BibleRef{撒上 15:23}并且，\Quote{在雅各伯中不应有占卜，在以色列中不应有法术。}\BibleRef{户 23:23}你不可为你的孩子使用邪术或净化的仪式。你不可做占卜者，也不可凭大鸟或小鸟占卜。你不可学习邪恶的技艺；因为这一切都是法律所禁止的。不可心存恶意，否则你将陷入不堪的罪恶。你不可说污秽的话，不可有淫荡的眼神，不可醉酒；因为从这些事中产生淫乱和奸淫。不可贪爱钱财，免得你\Quote{事奉金钱而非天主。}\BibleRef{玛 6:24}不可虚荣，不可傲慢，不可自高自大。因为从这一切中生出骄矜。要记住那一位说：\Quote{主，我的心不傲慢，我的眼不高傲；我不从事伟大的事，也不从事超过我能力的事；但我谦卑。}\par

% structure: p0016 source=h3 target=subsubsection reason=relative-heading-rank; base=h2

\subsubsection{禁止抱怨、傲慢、骄傲和狂妄。}

七、不可抱怨，要记念那些曾抱怨\Person{梅瑟}的人所受的惩罚。不可固执己见，不可心怀恶意，不可心硬，不可暴躁，不可小气；因为这一切都引向亵渎。但要温柔，如\Person{梅瑟}和\Person{达味}，因为\Quote{温柔的人将继承大地。}\BibleRef{玛 5:5}\par

% structure: p0018 source=h3 target=subsubsection reason=relative-heading-rank; base=h2

\subsubsection{关于忍耐、纯朴、温柔和坚忍。}

八、要缓慢发怒；因为这样的人非常明智，因为\Quote{急躁的人是非常愚昧的。}要仁慈；因为\Quote{仁慈的人是有福的，因为他们要得到仁慈。}\BibleRef{玛 5:7}要真诚、安静、良善，\Quote{敬畏天主的言语。}\BibleRef{依 66:2}你不可自高，像法利塞人那样；因为\Quote{凡自高的必被降卑，}\BibleRef{路 18:14}并且\Quote{在人所尊贵的，在天主面前是可憎的。}\BibleRef{路 16:15}你不可以在你灵魂中自恃；因为\Quote{自信的人必陷入祸患。}你不可与愚昧人同行，但要与智慧者和义人同行；因为\Quote{与智慧者同行必得智慧，与愚昧人同行必被显露。}\BibleRef{箴 13:20}要以平静的心接受降临于你的苦难，不要过分悲伤地接受人生的变故，要知道天主必给你赏报，如同给了\Person{约伯}和\Person{拉匝禄}一样。\par

% structure: p0020 source=h3 target=subsubsection reason=relative-heading-rank; base=h2

\subsubsection{我们应敬重基督徒导师超过父母——前者使我们得福，后者仅使我们得生。}

九、你要敬重那对你宣讲天主圣言的人，日夜思念他；你要尊敬他，不是视之为你生身的作者，而是视之为你获得福乐的因由。因为哪里有天主的道理，那里就有天主同在。你应每日寻求圣人们的仪容，好使你心悦诚服于他们的话语。\par

% structure: p0022 source=h3 target=subsubsection reason=relative-heading-rank; base=h2

\subsubsection{我们不应与圣人们分裂，而应在争吵者之间调解，公正审判，不可偏待。}

十、你不可在圣人们中制造分裂，但要记念\Person{科辣黑}党徒的结局。你应在纷争者之间调解，如同\Person{梅瑟}劝他们成为朋友那样。你要公正审判；因为\Quote{审判属于主。}当你为罪责斥责人时，不可偏袒；但要像\Person{厄里亚}和\Person{米加雅}对待\Person{阿哈布}，\Person{厄提约丕雅人厄贝得默肋客}对待\Person{漆德克雅}，\Person{纳堂}对待\Person{达味}，\Person{若翰}对待\Person{黑落德}那样。\par

% structure: p0024 source=h3 target=subsubsection reason=relative-heading-rank; base=h2

\subsubsection{关于心怀二意和沮丧的人。}

十一、你祈祷时不要心怀疑惑，不知是否蒙应允。因为主在海面上对我\Person{伯多禄}说：\Quote{小信德的人，你为什么疑惑？}\BibleRef{玛 14:31} \BibleQuote{不要准备伸手接受，却要关闭手而不施予。}\BibleRef{德 4:31}\par

% structure: p0026 source=h3 target=subsubsection reason=relative-heading-rank; base=h2

\subsubsection{关于行善。}

XII. 若你靠双手劳作有所收获，就当施与，好为你的罪赚取救赎；因为经上说：\BibleQuote{施舍与信德的行为能清除罪恶。} \BibleRef{箴 16:6} 你不可吝啬而不济贫，也不可在施与时发怨言，因为你晓得谁将报答你的酬报。因为他说：\BibleQuote{怜悯穷人的，便是借给上主；按他的施与，必得偿还。} \BibleRef{箴 19:17} 你不可转身不理困乏的人，因为他说：\BibleQuote{掩耳不听穷人哀求的，他自己呼求时，也无人应允。} \BibleRef{箴 21:13} 你当在一切事上与你的弟兄共享，不可说你的财物是你自己的；因为生活必需品的共同分享是天主为众人制定的。你不可从你的儿子或女儿身上收回你的手，而要自幼教导他们敬畏天主；因为他说：\BibleQuote{管教你的儿子，他必使你得到美好的盼望。} \BibleRef{箴 19:18}\par

% structure: p0028 source=h3 target=subsubsection reason=relative-heading-rank; base=h2

\subsubsection{主人应如何对待仆人，仆人应如何服从}

XIII. 你不可心怀苦毒地命令你的男仆或女仆，因为他们也信靠同一位天主，免得他们向你叹息，使天主的义怒临于你。而你们作仆人的，\BibleQuote{要顺服你们的主人} \BibleRef{弗 6:5}，如同顺服天主的代表，存敬畏的心，\BibleQuote{如同顺服主，不是顺服人} \BibleRef{弗 6:7}。\par

% structure: p0030 source=h3 target=subsubsection reason=relative-heading-rank; base=h2

\subsubsection{论伪善、守法与告罪}

XIV. 你当憎恶一切伪善；凡上主所喜悦的，你当遵行。不可离弃主的诫命，却要持守你所领受的，不可增加，也不可减少。\BibleQuote{不可加添祂的话语，免得祂责备你，你就成了说谎的。} \BibleRef{箴 30:6} 你当向主你的天主认罪，不可再犯，好使你因主你的天主而蒙福，祂不喜悦罪人死亡，却喜悦他悔改。\par

% structure: p0032 source=h3 target=subsubsection reason=relative-heading-rank; base=h2

\subsubsection{论孝敬父母}

XV. 你当孝敬你的父亲和母亲，他们是你的生身之源，好使你在主你天主所赐给你的地上得享长寿。不可忽视你的弟兄或亲属，因为经上说：\BibleQuote{不可忽视与你亲近的人。} \BibleRef{依 58:7}\par

% structure: p0034 source=h3 target=subsubsection reason=relative-heading-rank; base=h2

\subsubsection{论对君王和官长的服从}

XVI. 你当敬畏君王，知道他的立定是出于主。你当尊敬他的官长，如同天主的仆役，因为他们是一切不义的复仇者。你要甘心乐意地向他们纳税、进贡和献上一切供奉。\par

% structure: p0036 source=h3 target=subsubsection reason=relative-heading-rank; base=h2

\subsubsection{论祈祷者的纯洁良心}

XVII. 在你尚未除去苦毒之前，不可在你作恶的日子前来祈祷。这是生命之路，愿你们因我们的主耶稣基督而得以行在其中。\par

% structure: p0038 source=h3 target=subsubsection reason=relative-heading-rank; base=h2

\subsubsection{后来由仇敌的网罗引入的道路，充满不敬与邪恶}

XVIII. 死亡的道路由其恶行而显明：其中有不认识天主，引入许多邪恶、紊乱和骚动；由此产生凶杀、奸淫、淫乱、伪证、非法情欲、偷窃、偶像崇拜、邪术、巫术、抢劫、假见证、伪善、心怀二意、诡诈、骄傲、恶毒、傲慢、贪婪、污秽言语、嫉妒、自信、傲慢、狂妄、无耻、迫害善人、敌视真理、喜爱谎言、不知正义。因为行这些事的人，不从善，也不行公义的判断；他们不为善守望，却为恶守望；温良和忍耐远离他们，他们喜爱虚浮之事，追逐报酬，不怜恤穷人，不为受苦之人劳力，也不认识造他们的主；杀害婴孩，毁灭天主的创造，转离困乏之人，加重受苦者的痛苦，谄媚富人，藐视穷人，满身罪恶。愿你们这些孩子，得以脱离这一切。\par

% structure: p0040 source=h3 target=subsubsection reason=relative-heading-rank; base=h2

\subsubsection{我们不可偏离虔敬之路，或左或右。立法者的劝勉}

XIX. 要当心，不可有人引诱你离开虔敬；因为祂说：\Quote{你不可偏离此路，或左或右，好使你在你所行的一切事上有智慧。} 因为如果你不偏离正路，你便不会不虔敬。\par

% structure: p0042 source=h2 target=subsection reason=relative-heading-rank; base=h2

\subsection{第二段 论信徒品格的塑造与对天主的感恩}

% structure: p0043 source=h3 target=subsubsection reason=relative-heading-rank; base=h2

\subsubsection{我们不可轻视摆在我们面前的任何食物，而要怀着感恩和有序的方式享用}

XX. 至于各种食物，主对你们说：\Quote{你们要吃地上的美物；} 又说：\Quote{凡有血肉的，你们都可以吃，如同青菜一样；} 但是，\Quote{你们要倒出血来。} 因为 \Quote{不是入口的，而是出口的，才污秽人；} 我是指亵渎、毁谤，以及其他类似的事。但要 \Quote{以正义吃地里的肥油。} 因为 \Quote{凡令人喜悦的，都是祂的；凡属美好的，都是祂的。麦子赐给少年人，酒使少女欢畅。} 因为 \Quote{离了祂，谁能吃，谁能喝呢？} 智慧者厄斯德拉也告诫你们说：\Quote{你们去，吃肥美的，喝甘甜的，不要忧愁。} \BibleRef{厄下 8:10}\par

% structure: p0045 source=h3 target=subsubsection reason=relative-heading-rank; base=h2

\subsubsection{我们当避免吃祭偶像之物}

XXI. 但要禁戒祭偶像之物；\BibleRef{格前 10:20} 因为他们献祭是为了敬拜魔鬼，即是对独一天主的不敬，免得你们与魔鬼有份。\par

% structure: p0047 source=h3 target=subsubsection reason=relative-heading-rank; base=h2

\subsubsection{我们主的规定：应如何施洗，并受洗归于祂的死亡}

XXII. 论圣洗圣事，主教或司铎，我们已给出指示，如今再说：你应照主所命令我们的方式付洗，祂说：\BibleQuote{你们要去，使万民成为门徒，因圣父、圣子、圣神之名给他们付洗（教训他们遵守我所吩咐你们的一切）。}\BibleRef{玛 28:19} 即差遣者的圣父、来临者的基督、作证者的护慰者之名为他们付洗。但你应事先以圣油傅抹那人，随后以水付洗，最后以圣膏封印；使傅油成为参与圣神，水成为基督死亡的象征，圣膏成为盟约的印记。但若既无圣油也无圣膏，则水足以用于傅油、封印及向那与基督同死或正与基督同死者的宣认。但在圣洗之前，领洗者应当守斋；因为主自己先受若翰的洗，后在旷野住了四十天四十夜。祂受洗后守斋，并非祂需要洁净、守斋或除罪——祂本性纯洁圣善——乃是为向若翰证明真理，并为我们立下榜样。因此，我们的主受洗并非归于祂的苦难、死亡或复活——那时这些事尚未发生——乃是为另一目的。故此祂在受洗后凭自己的权柄守斋，作为若翰的主。但那位将受洗归于祂死亡的人，应先守斋，然后受洗。因为与基督同葬且同复活的人，理当不在祂的复活本身显出悲伤。因为人不是我们救主性体之主，主人是一位，仆人是另一位。\par

% structure: p0049 source=h3 target=subsubsection reason=relative-heading-rank; base=h2

\subsubsection{应在一周中的哪些日子守斋，哪些日子不守，及原因。}

XXIII. 但你们的守斋不可与伪君子相同；他们在每周第二日和第五日守斋。但你们要么守满五天，要么在每周第四日和预备日守斋，因为第四日对主的判决发出，犹达斯那时答应为银钱出卖祂；预备日你们必须守斋，因为主在那日于般雀比拉多手下受十字架之死。但应将安息日与主日视为庆节；因为前者是创造之纪念，后者是复活之纪念。但在全年中你们只应守一个安息日，即主埋葬的安息日，那日人应当守斋，而非庆节。因为造物主那时在地下，为祂的哀恸强于为创造的喜乐；因为造物主按其本性与尊严，比祂的受造物更可敬。\par

% structure: p0051 source=h3 target=subsubsection reason=relative-heading-rank; base=h2

\subsubsection{什么样的人应当诵念主所传授的那段祈祷文。}

XXIV. 但\Quote{你们祈祷时，不要像伪君子一样}；而应照主在福音中给我们规定的祈祷：\Quote{我们的天父，愿你的名被尊为圣，愿你的国来临，愿你的旨意承行，如同在天上，也在地上；求你今天赐给我们我们日用的食粮；并宽免我们的债，如同我们也宽免欠我们债的人；不要让我们陷入诱惑，但救我们脱离凶恶；因为国度属于你，直到永远。阿们。}一日祈祷三次，预先预备自己，好配得圣父的义子名分；免得你们不配地称祂为父，而受祂责备，正如从前对以色列——祂的首生子——所说：\BibleQuote{我若是父亲，我的荣耀在哪里？我若是主，我的敬畏在哪里？}\BibleRef{拉 1:6} 因为父亲的荣耀是子女的圣洁，主人的尊荣是仆人的敬畏，相反则是羞辱与混乱。因为祂说：\BibleQuote{因你们，我的名字在外邦人中被亵渎。}\BibleRef{依 52:5}\par

% structure: p0053 source=h3 target=subsubsection reason=relative-heading-rank; base=h2

\subsubsection{奥秘的感恩祈祷。}

XXXV. 要常怀感恩，如同忠信诚实的仆人；关于感恩祭的祈祷，这样说：我们感谢你，我们的父，因那生命，你藉你的儿子耶稣使我们认识的生命；藉着祂，你创造了万有，并眷顾整个世界；你派遣祂为我们的救恩成为人，允许祂受苦和死亡，又使祂复活，乐意光荣祂，并让祂坐在你的右边；藉着祂，你应许我们死人复活。全能的主，永生的天主，求你将你的教会从地极聚集到你的国度里，正如这谷粒曾分散，如今成为一体。我们也感谢你，我们的父，因耶稣基督的宝血，为我们所倾流，并因祂的宝贵身体，我们举行此纪念，正如祂自己所规定的，\BibleQuote{要宣告祂的死。}\BibleRef{格前 11:26} 因为藉着祂，荣耀归于你，直到永远。阿们。未入门者不得吃这些；只有那些受洗归于主死的人才可。若未入门者隐瞒自己而分食，\Quote{他是吃永罚}，因为不属于基督的信德，他分吃了不应分吃之物，自取惩罚。若有人因无知而分食，要快快教导他，使他入门，免得他出去藐视你们。\par

% structure: p0055 source=h3 target=subsubsection reason=relative-heading-rank; base=h2

\subsubsection{参与神圣恩赐时的感恩祈祷。}

二十六．领圣体后，应这样感谢：我们感谢祢，天主，我们救主耶稣的父，为祢的圣名，祢使这圣名居住在我們中间；也为祢藉着祢的圣子耶稣赐给我们的知识、信德、爱德与不朽。祢，全能的主，宇宙的天主，藉着祂创造了世界及其中的万物，并在我们的灵魂中种植了法律，预先为人的方便准备了事物。我们圣洁无瑕的祖先亚巴郎、依撒格、雅各伯的天主，祢忠信的仆人；祢，大能、信实、真实、无欺于诺言的天主；祢曾派遣祢的基督耶稣到世上与人同住，祂本是天主圣言和天主，却成了人，要从根本上消除错误：即使现在，求祢藉着祂，记念祢的圣教会，这是祢用祢基督的宝血所赎买的，保护她脱离一切邪恶，使她在祢的爱和真理中达于完美，并领我们全体进入祢所预备的国度。愿祢的国来临。\Quote{贺三纳于达味之子！奉主名而来的当受赞美。}——主天主，曾以肉身显现于我们。若有人是圣的，让他近前来；若有人不是，让他藉着悔改成为圣的。也允许你们的长老们谢恩。\par

% structure: p0057 source=h3 target=subsubsection reason=relative-heading-rank; base=h2

\subsubsection{关于神秘香膏的感恩祈祷。}

二十七．关于香膏，应这样感谢：我们感谢祢，天主，全世界的创造者，既为香膏的芬芳，也为祢藉着祢的圣子耶稣让我们认识的不朽。因为荣耀和权能永远属于祢。阿们。凡到你们这里来并这样感恩的人，要接纳他作为基督的门徒。但如果他宣讲另一种道理，与基督藉着我们传给你们的道理不同，这样的人你们不可允许他谢恩；因为这样的人是侮辱天主，而不是光荣祂。\par

% structure: p0059 source=h3 target=subsubsection reason=relative-heading-rank; base=h2

\subsubsection{我们不可对领圣体漠不关心。}

二十八．但凡来到你们这里的人，应先审查，然后接纳；因为你们有智慧，能分辨左右，也能区分假教师和真教师。但当一个教师来到你们这里，要乐意供给他所需。即使假教师来了，你们也要给他生活所需，但不可接受他的错误。你们也不可与他一同祈祷，免得你们也和他一样被玷污。每个来到你们这里的真先知或真教师，都配得他的养生之物，因为他是正义言语的工人。\BibleRef{玛 10:41}\par

% structure: p0061 source=h3 target=subsubsection reason=relative-heading-rank; base=h2

\subsubsection{关于祭品的条例。}

二十九．凡酒榨、打谷场、牛群和羊群中的一切初果，你们应献给司铎，\EditorialNote{户籍纪第18章}，好使你们的仓库、粮仓和土地的产品蒙受祝福，你们因五谷、酒和油而变得强壮，牛群和羊群得以繁衍。你们应将收成的十分之一给孤儿、寡妇、穷人和外方人。你们的热面包初果、酒桶、油、蜂蜜、坚果、葡萄的初果，以及其他初果，应献给司铎；但银钱、衣服和所有财产的初果，应给孤儿和寡妇。\par

% structure: p0063 source=h3 target=subsubsection reason=relative-heading-rank; base=h2

\subsubsection{我们应当如何聚集，并庆祝我们救主的复活节。}

三十．在主复活的日子，即主日，你们要聚集在一起，不可缺席，感谢天主，赞美祂，因为天主藉着基督赐予你们这些慈悲，并拯救你们脱离无知、错误和束缚，好使你们的祭献无玷，蒙天主悦纳。天主曾论及祂的普世教会说：\Quote{在每一处，都有人为我焚香并献上洁净的祭品；因为我是伟大的君王，全能的上主说，我的名在万邦中受尊崇。}\par

% structure: p0065 source=h3 target=subsubsection reason=relative-heading-rank; base=h2

\subsubsection{被祝圣的人应具备什么资格。}

三十一．你们要先祝圣配得上主的主教、长老和执事，他们应是虔诚、正义、温和、不贪财、爱真理、经得起考验、圣洁、不徇私情，能教导虔敬的道理，并正确分解主的道理。\BibleRef{弟后 2:15} 要尊敬这样的人，如同你们的父亲，如同你们的主人，如同你们的恩人，如同你们福祉的根由。要互相规劝，不要动怒，而要温和，以仁爱与和平。要遵守主所命令你们的一切。要警醒自己的生命。\Quote{你们要束上腰，点着灯，如同等候自己主人的人，看他何时回来，或是傍晚，或是半夜，或是鸡叫时，或是清晨。因为在他想不到的时刻，主人会来；若他们给主人开门，那些仆人就有福了，因为他们被找到时是醒着的。主人要束上腰，令他们坐席，并前来伺候他们。} \BibleRef{路 12:35} 因此，你们要醒寤祈祷，免得沉睡至死。因为你们以前的好行为，若在生命的最后偏离了真信德，对你们毫无益处。\par

% structure: p0067 source=h3 target=subsubsection reason=relative-heading-rank; base=h2

\subsubsection{关于未来之事的预言。}

三十二。因为末世将有假先知增多，以及那些败坏圣言的人；羊将变成狼，爱将变成恨：由于罪恶增多，许多人的爱德将会冷淡。因为人要彼此仇恨、迫害、出卖。那时，世界的欺骗者、真理的仇敌、谎言之王将要出现，\EditorialNote{得撒洛尼后书第二章}主耶稣将\Quote{用祂口中的气息消灭他，用祂的嘴唇除掉恶人；许多人将因他而跌倒。但坚持到底的，必要得救。那时，人子的记号要在天上出现；}\BibleRef{依 11:4}此后将有总领天使的号角声；在那间隔中，沉睡的人将苏醒。那时，主将降临，祂的众圣人与祂一同降临，在云上有大震动，有祂大能的天使，\BibleRef{玛 16:27}坐在祂国度的宝座上，审判魔鬼——世界的欺骗者，并照各人的行为报应各人。\Quote{那时，恶人将进入永罚，义人却进入永生，}\BibleRef{玛 25:46}承受那些\Quote{眼所未见，耳所未闻，人心所未想到的，天主为爱祂的人所预备的；}\BibleRef{格前 2:9}他们要在天主的国里，在基督耶稣内喜乐。既然我们从祂那里蒙受了如此伟大的恩惠，让我们成为祂的恳求者，以不断的祈祷呼求祂，说：---\par

% structure: p0069 source=h3 target=subsubsection reason=relative-heading-rank; base=h2

\subsubsection{表明天主各种眷顾的祷词}

三十三。我们永恒救主，诸神之王，唯一全能者，主，万有的天主，我们圣洁无瑕的祖先及我们先祖的天主；亚巴郎的天主，依撒格的天主，雅各伯的天主；祢是仁慈悲悯，缓于发怒，富于仁慈的；万心在祢面前赤露，万心被祢鉴察，一切秘密思念向祢显明：义人的灵魂向祢高声呼号，虔敬者的希望寄于祢，祢是无瑕者的父，祢垂听以正直呼求祢者的祈祷，且知道未曾言明的祈求：因为祢的眷顾直达人内心的深处；凭祢的知识，祢探究每个人的思念，在全地各处，祈祷恳求的馨香向祢上达。哦，祢指定今世为正义的战场，为众人敞开了仁慈之门，并藉植入的知识、自然的判断和法律的劝诫，向每个人显明：财富的占有并非永久，美貌的装饰并非恒常，我们的力量和气力容易消散；一切都是虚幻和虚空；唯有真诚信德的良心穿越诸天，带着真理归来，把握那将来喜乐的右手。而且，在万物复兴的应许实现之前，灵魂本身因盼望而雀跃欢欣。因为从我们祖先亚巴郎身上的那真理，当他改变道路时，祢藉异象引导他，教导他今世是怎样的一种状态；知识先于他的信德，信德是知识的结果；盟约紧随他的信德。因为祢曾说：\Quote{我要使你的后裔如天上的星辰，并如海边的沙粒。}再者，当祢赐给他依撒格，并知道他在生活方式上与他相似时，祢就被称为他的天主，说：\Quote{我要作你的天主，并作你后裔的天主。}\BibleRef{创 26:3}当我们祖先雅各伯被派往美索不达米亚时，祢向他显示基督，并藉他说：\Quote{看，我与你同在，我要使你增多，极大繁衍。}祢又向祢忠信圣洁的仆人梅瑟，在荆棘显现中说：\Quote{我是自有者；这是我永远的名字，万世万代的纪念。}\BibleRef{出 3:14}哦，祢是亚巴郎后裔的伟大保护者，祢永受赞美。\par

% structure: p0071 source=h3 target=subsubsection reason=relative-heading-rank; base=h2

\subsubsection{表明天主各种创造的祷词}

祢受赞美，主啊，万代之王，祢藉基督创造了整个世界，并藉着祂在起初将混乱的部分整顿有序；祢藉穹苍将水分为上下，并在其中注入生命之气；祢奠定了大地，铺展了穹苍，并以精确的秩序安排了每一受造物。因为藉祢的权能，主啊，世界得以美化，穹苍如拱顶固定在我们之上，并以星辰照耀，为我们在黑暗中提供安慰。光与太阳受造，为区分日子和产生果实，月亮因盈亏而定季节的变换；一个称为夜，另一个称为昼。穹苍出现在深渊之中，祢命令水聚在一处，使旱地露出。至于海洋本身，谁能述说呢？它从大洋汹涌而来，却又退回，因祢的命令为沙所阻？因为祢曾说过：\Quote{你的浪涛必在此止住。}\BibleRef{约 38:11}祢也使它能承载大小生物，并使之可航行船只。然后大地披上绿装，种植了各类花草和多样的树木；照耀的光体，那些植物的滋养者，保持其不变的轨道，丝毫不偏离祢的命令。祢命令它们在哪里升起和落下，以作为季节和年岁的记号，使人的工作得以定期循环。随后，各类动物受造——属地的、属水的、属空的、以及兼属空水的；祢眷顾的巧妙智慧仍为每个生物赐予合适的眷顾。因为祂既能创造不同种类，也不轻视对每个种类施行不同的眷顾。在创造的终结，祢指示祢的智慧，并造了一个理性的受造物作为世界的公民，说：\Quote{让我们按我们的肖像，按我们的模样造人；}\BibleRef{创 1:26}并显示他为世界的装饰，用四种元素这些原始物形成他的身体，却从无中准备了灵魂，赐予他五官，并在他感官之上安置了心灵作为灵魂的引导者。此外，主天主啊，谁能相称地述说雨云的运行，闪电的照耀，雷声的轰隆，以供应合宜的食物，以及最宜人的空气温度？但当人悖逆时，祢剥夺了他应得的生命作为赏报。然而祢并未永远毁灭他，而是使他暂时安眠；祢藉着誓约召唤他复活，解除了死亡的束缚，啊，祢是复活死者的主，藉着耶稣基督，他是我们的望德。\par

% structure: p0073 source=h3 target=subsubsection reason=relative-heading-rank; base=h2

\subsubsection{一篇谢恩祈祷，宣告天主对祂所造万物的眷顾。}

XXXV. 大哉，全能之主，尔能力何其大，尔智慧无有计数。我等造物主与救主，富有恩惠，宽容忍耐，广施仁慈，不从不肖受造物收回尔之救恩：因尔本性为善，宽赦罪人，召其悔改；尔之训诫出于尔慈悲之肠。盖若我等即刻被传至审判，我等将如何存立？虽蒙尔如此长久忍耐，我等仍难摆脱悲惨境况。诸天宣扬尔统治，大地震动，悬于虚空，显扬尔之稳固不移。海洋波澜汹涌，养活无数群生，却被沙粒束缚，如敬畏尔命，迫使众生呼喊：\Quote{主，尔之化工何其多！尔以智慧成就一切：大地充满尔之受造。}明亮天使军队及理智之灵对\OriginalTerm{帕摩尼}{Palmoni}说：\Quote{唯有一圣者；}六翼革鲁宾与圣色拉芬同唱凯歌，以永不间断之声呼号：\Quote{圣、圣、圣，万军的上主！天地充满尔之荣耀；}\BibleRef{依 6:3}还有各类品级之众：天使、总领天使、座者、宰制者、率领者、掌权者、大能者，高呼说：\Quote{愿主之荣耀从祂之处受赞颂。}\BibleRef{则 3:12}但以色列——尔在地上之教会，从外邦人中选出——效法天上诸能，昼夜以全心及甘愿之灵魂歌唱：\Quote{天主的战车万万千千，主在西乃山，在圣所中。}诸天认识祂，祂将天穹如石砌成拱形，悬于虚空，结合水土，散布生命之气，并与之结合火以保暖及驱暗。星辰队伍令人惊叹，宣示数算它们者，指明称呼它们名者；动物宣示赋予生命者；树木显明使其生长者：这一切受造之物皆由尔圣言造成，显扬尔能力之伟大。因此，每个人都当从灵魂中经由基督，以其余万有之名，向尔献上赞歌，因基督按尔指派统驭万有。尔于恩惠中仁慈，于慈悲肠中施惠，惟尔全能：因尔愿时，能力即显；尔永恒能力灭火、封狮口、驯鲸、医病、统御万有之力、倾覆敌军、击倒自高自大之民。尔是在天之天主，在地之天主，在海之天主，在有限事物中之天主，尔自己不受任何限制。因尔威严无界；主，这非我等之言，乃尔仆人之神谕：\Quote{你心里当知道，主你的天主是天上地下唯一的天主，除祂以外别无他神。}\BibleRef{申 4:39}因除尔以外别无真神，除尔以外别无圣者，主，知识之神，诸圣之神，超越一切圣者者；因他们皆由尔手成圣。尔乃荣耀至高，本性不可见，判断不可测；生命无缺，永恒不变不衰，作为无劳，伟大无限，卓越永恒，居处不可接近，住所不变，知识无始，真理不变，作为无需助手，统治不可夺，主权无后继，国度无尽，力量不可挡，军队无数：因尔是智慧之父，是受造之物的创造者，通过中保作为原因；眷顾之赐予者，立法者，供应匮乏者，惩恶者，赏善者；基督之天主与父，以及虔敬于祂者之主，其许诺无误，其审判不受贿赂，其意念不变，其虔诚无间，其感恩永存，藉着祂，从一切理性与圣洁本性，当受合宜之崇拜。\par

% structure: p0075 source=h3 target=subsubsection reason=relative-heading-rank; base=h2

\subsubsection{纪念基督降生成人及其对圣人之各种眷顾的祈祷文}

XXXVI. 全能的主，祢藉着基督创造了世界，并为此设立了安息日，因为在那一天祢使我们从工作中安息，以默思祢的律法。祢也设立了节期以使我们的灵魂欢欣，使我们记起祢所创造的智慧：祂如何为了我们而甘愿由女人所生；祂显现在生命中，在祂的圣洗中证明自己；那显现的祂既是天主又是人；祂因祢的允许为我们受苦，死去，并藉祢的大能复活。为此，我们隆重集会，在主日庆祝复活节庆，并因那战胜死亡、带来生命与不朽的那位而欢欣。因为藉着祂，祢将外邦人召选为专属祢的子民，真以色列，蒙天主所爱并看见天主的子民。因为主啊，祢将我们的祖先从埃及地领出，从铁炉、泥土和制砖中拯救他们，从法老和他手下的人手中救赎他们，领他们过海如过干地，在旷野中容忍他们的行为，赐给他们各样美善。祢赐给他们律法或十诫，由祢的声音宣告并用祢的手书写。祢命令遵守安息日，不是给他们懒惰的机会，而是敬虔的机会，为使他们认识祢的大能，并禁止邪恶；祢将他们限制在如同圣洁的圈子内，为教导的缘故，在第七个时期欢欣。为此设立了一周、七周、第七个月、第七年，以及这些的循环——禧年，即第五十年为豁免，使人没有借口假装无知。为此，祂允许人在每个安息日休息，以免有人在安息日怒气冲冲地吐出一个字。因为安息日是创造的停止，世界的完成，对律法的探究，并为天主赐予人的恩惠而献上感恩的赞美。这一切，主日都超越之，并显示中保本身、供应者、立法者、复活的缘由、一切受造物的首生者、天主圣言，以及由玛利亚独生（无男人）、圣洁生活、在本丢彼拉多手下被钉十字架、死亡并死者中复活的人。因此，主日命令我们为一切向祢献上感谢。因为这是祢赐予的恩宠，因其伟大而掩盖了所有其他福分。\par

% structure: p0077 source=h3 target=subsubsection reason=relative-heading-rank; base=h2

\subsubsection{包含祂眷顾之纪念，并枚举藉基督以天主眷顾赐予圣人们的各种恩惠之祈祷。}

XXXVII. 祢，成就了祢藉先知所作的应许，并怜悯了熙雍，同情了耶路撒冷，通过在祂中间高举祢的仆人达味的宝座，藉着基督的诞生——祂按肉躯由达味后裔、由一位童贞女独生；求祢现在，主天主啊，接纳从祢外邦子民口中发出的祈祷，他们诚心呼求祢，正如祢接纳了各世代义人的祭献一样。首先，祢看中了亚伯尔的祭献，\EditorialNote{创世纪4}并接纳了它，就如祢接纳了诺厄从方舟出来后的祭献一样；\EditorialNote{创世纪8}亚巴郎的祭献，当他离开加色丁地时；\EditorialNote{创世纪12}依撒格在誓约井边的祭献；雅各伯在贝特耳的祭献；梅瑟在旷野的祭献；亚郎在死人活人之间的祭献；\EditorialNote{户籍纪16}农的儿子若苏厄在基耳加耳的祭献；\EditorialNote{若苏厄书5}基德红在磐石和羊毛前的祭献（在他犯罪之前）；玛诺亚和他妻子在田间的祭献；三松在口渴时（在过犯之前）的祭献；耶弗他在战争中的祭献（在他鲁莽的誓愿之前）；巴辣克和德波辣在息色辣时期的祭献；撒慕尔在米兹帕的祭献；\EditorialNote{撒慕尔纪上7}达味在耶步斯人敖尔难的禾场上的祭献；撒罗满在基贝红和耶路撒冷的祭献；厄里亚在加尔默耳山的祭献；\EditorialNote{列王纪上18}厄里叟在荒芜之泉的祭献；\EditorialNote{列王纪下2}约沙法特在战争中的祭献；希则克雅在病中及关于散乃黑黎布的祭献；默纳舍在加色丁地（在他过犯之后）的祭献；\EditorialNote{编年纪下33}约史雅在帕萨的祭献；厄斯德拉在回归时的祭献；\EditorialNote{厄斯德拉上8}达尼尔在狮子坑中的祭献；约纳在鲸鱼腹中的祭献；\EditorialNote{约纳书2}三童在火窑中的祭献；\EditorialNote{达尼尔书3}亚纳在会幕中约柜前的祭献；\EditorialNote{撒慕尔纪上1}乃赫米雅在重建城墙时的祭献；则鲁巴贝耳的祭献；玛塔提雅和他儿子们在热忱中的祭献；雅厄耳在祝福中的祭献。现在也求祢接纳祢子民以知识、藉着基督在圣神内向祢献上的祈祷。\par

% structure: p0079 source=h3 target=subsubsection reason=relative-heading-rank; base=h2

\subsubsection{为求义人助佑的祈祷。}

三十八、全能的主，我们为一切事感谢祢，因为祢没有从我们身上收回祢的仁慈和怜悯；反而在每一世代中，祢拯救、解救、帮助、保护我们：因为祢在厄诺斯和厄诺克的日子、在梅瑟和若苏厄的日子、在民长日子、在撒慕尔和厄里亚及先知的日子、在达味和君王的日子、在厄斯德尔和摩尔德开的日子、在友弟德的日子、在犹大玛加伯及其兄弟的日子，都曾帮助；而在我们的日子，祢藉着祢的大司铎，祢的圣子耶稣基督，帮助了我们。因为祂救我们脱离了刀剑，使我们免于饥荒，并扶持我们；救我们脱离疾病，保护我们免受恶舌之害。为这一切，我们藉着基督感谢祢，祂赐给我们响亮的嗓音以宣认，又加给适合的舌头作为调音的乐器，以及合适的味觉、适当的触觉、用来默观的视觉、听声音的听觉、闻气味的嗅觉、工作的手和行走的脚。祢在母胎中从一个微小的滴点形成了这一切肢体；成形后，祢又赐给它不朽的灵魂，并把它作为理性受造物，就是人，带到光明中。祢以祢的法律教导他，以祢的诫命提升他；当祢暂时带来解体时，祢又应许了复活。因此，什么生命足够，什么样的年岁长度能足以使人感恩？完全地感恩是不可能的，但按我们的能力去感恩是公正合理的。因为祢拯救我们脱离了多神论的不虔敬，以及杀害基督者的异端；祢拯救我们脱离了错误和愚昧；祢将基督作为人派遣到人间，祂是独生天主；祢使护慰者居住在我们中间；祢为我们设立了天使；祢使魔鬼蒙羞；祢在我们不存在时创造了我们。祢在我们受造后照顾我们；祢为我们量定生命；祢赐给我们食物；祢应许了悔改。愿荣耀和敬拜因这一切归于祢，藉着耶稣基督，从现在直到永远，万世无疆。阿们。弟兄们，要默想这些事；愿主在地上与你们同在，并在祂父的国里与你们同在，祂父派遣了祂，并且祂\BibleQuote{拯救我们脱离腐败的奴役，进入祂荣耀的自由；}\BibleRef{罗 8:21}并且应许永生给那些藉着祂信了全世界之上主的人。\par

% structure: p0081 source=h2 target=subsection reason=relative-heading-rank; base=h2

\subsection{第三节　关于教导慕道者及使他们进入圣洗圣事}

在前面的指示中，已经说明了那些加入基督的人应当如何生活，以及他们应当藉着基督向天主献上怎样的感恩。但也不应让那些尚未加入的人得不到帮助，这是合理的。\par

% structure: p0083 source=h3 target=subsubsection reason=relative-heading-rank; base=h2

\subsubsection{如何教导慕道者基本要理}

三十九、因此，那将要接受关于虔敬之真理教导的人，在领受圣洗之前，应当接受关于无始之天主的认识、关于祂独生子的理解、关于圣神的确实承认的教导。让他学习创造的各个部分的次序、天主眷顾的系列、祢律法的各种安排。让他受教为何世界被造，为何人被任命为其中的公民；也要让他认识自己的本性为何；让他受教天主如何以水和火惩罚恶人，并在每一世代中荣耀圣人——我是指舍特、厄诺斯、厄诺克、诺厄、亚巴郎及其后裔、默基瑟德、约伯、梅瑟、若苏厄、加肋布、司铎丕乃哈斯，以及每一世代中的圣人；以及天主如何仍然照顾而不抛弃人类，而是在各个时期从他们的错误和虚妄中召唤他们归向真理的承认，从奴役和不虔敬中带领他们进入自由和虔敬，从不义到正义，从永死到永生。让那自愿领受圣洗的人在他是慕道者的期间学习这些和类似的事；让那覆手在他身上的人朝拜全世界的主宰天主，并感谢祂的创造，感谢祂派遣了基督，祂的独生子，为藉着涂抹过犯拯救人，为赦免不虔敬和罪恶，并为\BibleQuote{洁净他脱离肉体和心灵的一切污秽，}\BibleRef{格后 7:1}并按祂仁慈的美意圣化人，为启发他以祂旨意的知识，并照亮他心中的眼睛以思想祂奇妙的作为，并使他认识正义的审判，以致他憎恨一切不义之道，行走在真理之路上，使他堪当承受重生的圣洗，以得着儿子的名分，这儿子名分是在基督内的，为\BibleQuote{与基督同钉死相似，}\BibleRef{罗 6:5}怀着荣耀分享的盼望，他可以向罪而死，而向天主活着，在心思、言语和行为上如此，并一同被记在生命册上。在这感恩之后，让他教导他关于我们主降生成人的道理，以及关于祂的苦难、从死者中复活和升天的道理。\par

% structure: p0085 source=h3 target=subsubsection reason=relative-heading-rank; base=h2

\subsubsection{一项典章：在慕道者入门时，司铎应如何祝福他们，以及应教导他们什么}

XL . 当\OriginalTerm{慕道者}{catechumen}即将受洗时，让他学习关于\OriginalTerm{弃绝}{renunciation}魔鬼和与基督结合的事；因为他应先戒绝相反之事，然后才被接纳进入\OriginalTerm{奥迹}{mysteries}。他必须事先从心中清除一切邪恶的意向、一切污点和皱纹，然后才领受圣事；正如最熟练的农夫先清除田里长出的荆棘，然后才播种小麦，同样，你们也应先除去他们的一切不敬，然后在他们内播种虔诚的种子，并恩赐他们圣洗。因为我们的主也曾这样劝诫我们，先说：\Quote{使万民成为门徒；\BibleRef{玛 28:19}} 然后加上这句话：\Quote{因父、及子、及圣神之名给他们授洗。} 因此，请求受洗者应这样宣示他的弃绝：——\par

% structure: p0087 source=h3 target=subsubsection reason=relative-heading-rank; base=h2

\subsubsection{弃绝仇敌，并归顺于天主的基督。}

XLI. 我弃绝\OriginalTerm{撒殚}{Satan}，和他的作为，和他的\OriginalTerm{排场}{pomps}，和他的\OriginalTerm{敬拜}{worships}，和他的天使，和他的发明，以及他之下的一切事物。在他弃绝之后，让他在他的结合中说：我就与基督结合，我相信，并受洗归于一个\OriginalTerm{无始的存在}{unbegotten Being}，唯一真天主全能者，基督的圣父，万物的创造者和造作者，万物都来自他；并归于主耶稣基督，他的\OriginalTerm{独生子}{Only-begotten Son}，整个受造界的\OriginalTerm{首生者}{First-born}，他在万世之前因圣父的善意所生，万物藉他而造，无论天上的还是地上的，可见的与不可见的；他在末世从天降下，取了肉身，由童贞玛利亚所生，并按照他的天主和圣父的法律圣洁地生活，在\OriginalTerm{般雀·比拉多}{Pontius Pilate}手下被钉十字架，为我们而死，受难后第三天从死里复活，升天，坐在圣父的右边，将在世界末日带着荣耀再来审判活人死人，他的王国没有终结。我又受洗归于圣神，就是\OriginalTerm{护慰者}{Comforter}，他从创世之初就在所有圣人中工作，但后来按照我们的救主和主耶稣基督的应许，由圣父派遣给宗徒；并在宗徒之后，赐给所有相信\OriginalTerm{至圣天主教会}{Holy Catholic Church}的人；归于肉身的复活、罪的赦免、天国和来世的生命。在此誓愿之后，他来到圣油傅油。\par

% structure: p0089 source=h3 target=subsubsection reason=relative-heading-rank; base=h2

\subsubsection{\OriginalTerm{圣油}{Mystical Oil}傅油的感恩经。}

XLII. 现在，大司祭祝圣这\OriginalTerm{圣油}{mystical oil}，以赦免罪过，并作为圣洗的最初准备。因为他这样呼求\OriginalTerm{无始的天主}{unbegotten God}、基督的圣父、一切可见与不可见本性的君王，使他因主耶稣的名祝圣这圣油，并将属灵的恩宠和有效的力量、罪的赦免和圣洗宣认的最初准备赐给它，以便慕道者在傅油后能摆脱一切不敬，并按照独生子的命令成为配得\OriginalTerm{入门}{initiation}的人。\par

% structure: p0091 source=h3 target=subsubsection reason=relative-heading-rank; base=h2

\subsubsection{\OriginalTerm{圣水}{Mystical Water}的感恩经。}

XLIII. 此后，他来到水前，祝福并光荣主全能的天主、\OriginalTerm{独生子}{only begotten God}的圣父；司铎感谢他派遣了他的圣子为了我们而成为人，以拯救我们；使他服从那\OriginalTerm{降生成人}{incarnation}的一切法律，宣讲天国、罪的赦免和死人的复活。此外，他朝拜独生子本身，在其圣父之后，并为他感谢，因为他担当了为全人类死在十字架上，其预像指定为\OriginalTerm{重生的洗礼}{baptism of regeneration}。他也光荣他，因为那统治全世界的天主，因基督的名并藉他的圣神，没有抛弃人类，而是按照季节的不同调整了他的\OriginalTerm{眷顾}{providence}：起初赐给亚当本人一个\OriginalTerm{乐园}{paradise}，作为欢乐的居所，后来又出于眷顾赐予一条诫命，并公正地驱逐了犯罪者，但藉着他的良善没有完全抛弃他，而是在后来的世代中以各种方式教导他的后代；为了他们的缘故，在世界终结时，他派遣了他的圣子为人的缘故成为人，并承受一切无罪的人类情感。因此，让司铎甚至现在在洗礼中呼求他，并说：从天垂顾，祝圣这水，赐予它恩宠和能力，使那要受洗的人，按照你的基督的命令，与他同钉十字架，与他同死，与他同葬，与他同复活，获得在他内的\OriginalTerm{义子}{adoption}名分，使他死于罪而活于义。此后，当他因父、及子、及圣神之名给他施洗后，他要用香膏傅油，并接着说：——\par

% structure: p0093 source=h3 target=subsubsection reason=relative-heading-rank; base=h2

\subsubsection{\OriginalTerm{圣膏}{Mystical Ointment}的感恩经。}

XLIV. 主天主，无始无终、无上至尊者，普世的主宰，祢将福音知识的馨香散布于万民之中，求祢此时使这傅油对受洗者有效，以致祢基督的馨香坚定而稳固地常留于他；如今他既与祂同死，也愿与祂同复活。让他这样说及类似的话，因为这是为每个人按手的效果；除非一位虔诚的司铎为每个人如此诵念，否则候洗者只是像犹太人一样下到水里，只除去身体的污秽，而非灵魂的污秽。此后，让他站起来，祈祷主教导我们的那篇祈祷文。但复生者理应站起来祈祷，因为复活的人站立。因此，那与基督同死又同复活的人，让他站起来。但让他面向东方祈祷。因为在《编年纪下》第二卷中也记载：主的圣殿由撒罗满王建成后，在奉献节那天，司祭、肋未人和歌咏者面向东方站立，用钹和琴瑟赞美感谢天主，说：\BibleQuote{你们要赞美主，因为祂是良善的；因为祂的仁慈永远常存。} \BibleRef{编下 5:13}\par

% structure: p0095 source=h3 target=subsubsection reason=relative-heading-rank; base=h2

\subsubsection{为初果的祈祷。}

XLV. 但在前面的祈祷之后，让他这样祈祷说：全能的天主，祢的基督、祢独生子的父，赐给我无玷的身体、纯洁的心灵、警醒的理智、无误的认知、圣神的影响，以获得并确信享有真理，藉着祢的基督，愿荣耀藉着祂归于祢，在圣神内，直到永远。阿们。我们认为作这些关于慕道者的规定是合理的。\par

% structure: p0097 source=h2 target=subsection reason=relative-heading-rank; base=h2

\subsection{第四节 宗徒所规定的名单}

% structure: p0098 source=h3 target=subsubsection reason=relative-heading-rank; base=h2

\subsubsection{圣宗徒们派遣并祝圣了谁？}

XLVI. 关于在我们有生之年所祝圣的主教，我们让你们知道他们是：——耶路撒冷的主教雅各伯，我们的主的兄弟；他去世后，第二位是克罗帕的儿子西默盎；之后第三位是雅各伯的儿子犹达。凯撒勒雅（巴勒斯坦）的第一位是匝凯，他曾是税吏；之后是科尔乃略，第三位是特敖斐罗。安提约基雅的第一位是厄瓦狄约，由我伯多禄祝圣；依纳爵由保禄祝圣。亚历山大的第一位是阿尼雅诺，由马尔谷福音吏祝圣；第二位是阿维略，由路加（他也是福音吏）祝圣。罗马教会的第一位是克罗狄亚的儿子利诺，由保禄祝圣；\BibleRef{弟后 4:21} 利诺死后，第二位是克肋孟，由我伯多禄祝圣。厄弗所的第一位是弟茂德，由保禄祝圣；若望由我若望祝圣。斯米纳的第一位是阿黎斯托；之后是罗依的儿子斯特拉特阿斯；\BibleRef{弟后 1:5} 第三位是阿黎斯托。培尔加摩的加约。非拉铁非的德默特琉，由我。恳克勒雅的路基约，由保禄。克里特的弟铎。雅典的狄约尼削。腓尼基的的黎波里的玛拉托内。夫黎基雅的劳狄刻雅的阿尔希波。哥罗森的费肋孟。马其顿的波勒亚的敖乃息摩，他曾是费肋孟的仆人。迦拉达众教会的克勒斯肯。亚细亚诸堂区的阿桂拉和尼则达。厄基纳教会的克黎斯波。这些是受托管理主内各堂区的主教；要常记他们的教导，遵守我们的话。愿主与你们同在，从现在直到无穷的世代，正如祂将被提升到祂自己的天主和父那里时对我们所说的。因为祂说：\BibleQuote{看，我天天与你们在一起，直到今世的终结。阿们。} \BibleRef{玛 28:20}\par

% structure: p0100 source=h2 target=subsection reason=relative-heading-rank; base=h2

\subsection{第五节 日常祈祷}

% structure: p0101 source=h3 target=subsubsection reason=relative-heading-rank; base=h2

\subsubsection{晨祷。}

XLVII. \BibleQuote{光荣归于至高之天主，平安于地，恩宠于人。} \BibleRef{路 2:14} 我们赞美祢，我们歌颂祢，我们称颂祢；我们光荣祢，我们藉着祢的大司祭朝拜祢；祢是真天主，是唯一无始者，唯一不可接近者。因为祢的伟大光荣，主，天上的君王，全能天主圣父，主天主，无玷羔羊基督的父，祂除免世罪，求祢收纳我们的祈祷，祢坐在革鲁宾之上。因为惟有祢是圣的，惟有祢是主耶稣，普世受造之物的天主的基督，我们的君王，愿荣耀、尊威和朝拜藉着祢归于祢。\par

% structure: p0103 source=h3 target=subsubsection reason=relative-heading-rank; base=h2

\subsubsection{晚祷。}

XLVIII. \BibleQuote{孩子们，你们要赞美主：赞美主的名字。} 我们赞美祢，我们歌颂祢，我们为祢的伟大光荣称颂祢，主我们的君王，无玷羔羊基督的父，祂除免世罪。赞美宜于祢，歌颂宜于祢，光荣宜于祢，天主圣父，藉着圣子，在至圣圣神内，直到永远。阿们。\BibleQuote{主啊，现在可以照祢的话，让祢的仆人平安离去；因为我亲眼看见了祢的救恩，就是祢在万民面前所预备的，为启示外邦人的光，和祢子民以色列的荣耀。}\par

% structure: p0105 source=h3 target=subsubsection reason=relative-heading-rank; base=h2

\subsubsection{餐时祈祷。}

XLIX. 主，祢是可称颂的，祢从我幼年就养育我，祢赐食物给一切血肉。求祢使我们的心充满喜乐和欢愉，使我们常能知足，能多行各种善工，在基督耶稣我们的主内，愿荣耀、尊威和能力藉着祂归于祢，直到永远。阿们。\par

\SourceRecord{Translated by James Donaldson.；Ante-Nicene Fathers；Vol. 7.；Edited by Alexander Roberts, James Donaldson, and A. Cleveland Coxe.；Buffalo, NY: Christian Literature Publishing Co.,；1886.；<http://www.newadvent.org/fathers/07157.htm>.}
